\documentclass{article}
\usepackage{neurips_2026}
\usepackage{amsmath,amssymb,amsfonts}
\usepackage{amsthm}
\usepackage{algorithmic}
\usepackage{algorithm}
\usepackage{graphicx}
\usepackage{hyperref}
\usepackage{booktabs}
\usepackage{wrapfig}
\usepackage{multirow}
\usepackage{booktabs}
\usepackage{makecell}
\usepackage[table]{xcolor}
\usepackage{subcaption}

\usepackage{xcolor}
\usepackage{hyperref}

\definecolor{citeblue}{HTML}{0057B8}
\definecolor{linkcrimson}{HTML}{B23A48}
\definecolor{urlpurple}{HTML}{6F4AA8}

\hypersetup{
    colorlinks=true,
    citecolor=citeblue,
    linkcolor=linkcrimson,
    urlcolor=urlpurple
}

\newtheorem{theorem}{Theorem}[section]
\newtheorem{lemma}[theorem]{Lemma}
\newtheorem{assumption}{Assumption}[section]

\title{FAUST: Federated Asynchronous Update with Staggered Timescales \\
for Low-Communication Foundation Model Training}

\author{}

\begin{document}
\maketitle

\begin{abstract}
Large language model training with sequence-parallel data-parallel (SP-DP) is bottlenecked by the interplay between intra-step activation redistribution and
inter-step gradient synchronization. Existing infrequent-communication methods like FedAvg keep states local or reset them lack convergence guarantees and
are unstable under long-context sequence-parallel regimes where each device holds only a partial activation shard; Local Adam synchronizes all states
jointly and is convergent but triples the collective payload, and its uniform synchronization period cannot distinguish the semantically mandatory intra-step All-to-All from the temporally deferrable inter-step AllReduce. We propose
\textbf{FAUST} (\textbf{F}ederated \textbf{A}synchronous \textbf{U}pdate with \textbf{S}taggered \textbf{T}imescales), a compiler-orchestrated training system that composes compile-time sequence-parallel graph rewriting with tiered runtime synchronization, assigning independent averaging periods to parameters~($K_x$), first moments~($K_u{=}3K_x$), and second moments~($K_v{=}6K_x$) according to their impulse-response half-lives, while preserving convergence under the orthogonal collective constraint that intra-step All-to-All must retire before inter-step AllReduce begins. Our analysis shows that while first-moment drift dominates the convergence rate
in-distribution, high-fidelity convergence guarantees require at least periodic
synchronization of the second moment. Experiments on language
models up to 1.3B parameters show that FAUST incurs $303{\times}$ less wall-clock communication overhead than standard DDP
and $1.78{\times}$ less than the previous state-of-the-art DES-LOC, while achieving perplexity within $4.6\%$ of the fully-synchronized baseline in $27\%$ less wall-clock time. On bandwidth-constrained links, FAUST delivers
$48$--$303{\times}$ speedup over DDP.
\end{abstract}
%=============================================================================


\section{Introduction}
\label{sec:introduction}
%=============================================================================

Training large language models at scale requires two orthogonal classes of
collective communication: \emph{intra-step} All-to-All that redistributes
activations across the head dimension for sequence parallelism (SP), and
\emph{inter-step} AllReduce that aggregates gradients across data-parallel (DP)
replicas. On heterogeneous clusters where both classes share the same
interconnect, the inter-step AllReduce dominates wall-clock
time~\citep{li2020pytorch}. Local SGD~\citep{stich2019local} and
FedAvg~\citep{mcmahan2017communication} reduce this cost by averaging
parameters every $K \gg 1$ steps, but they neither handle the additional
per-parameter states of adaptive optimizers~\citep{kingma2017adammethodstochasticoptimization} nor
account for the intra-step SP collective that modern long-context
training~\citep{grattafiori2024llama3herdmodels} introduces.

Extensions to adaptive optimizers either keep optimizer states entirely
local~\citep{douillard2024dilocodistributedlowcommunicationtraining}, accumulating drift without convergence
guarantees, or reset them at each sync
point~\citep{sani2024futurelargelanguagemodel,sani2024photonfederatedllmpretraining}, destroying the exponential moving
average (half-life $\tau_{0.5}(\beta_1{=}0.9) \approx 6.6$ steps).
Local Adam~\citep{cheng2025convergencedistributedadaptiveoptimization} proves convergent intermittent-communication
adaptive optimization, but synchronises all states at a uniform period~$K$,
tripling the payload versus Local SGD. Moreover, the uniform-$K$ schedule cannot distinguish
mandatory intra-step A2A from deferrable inter-step AllReduce, preventing
composition with SP, a critical barrier as contexts exceed single-device memory
and require compiler-mediated sequence
sharding~\citep{jacobs2023deepspeedulyssesoptimizationsenabling,fang2024uspunifiedsequenceparallelism}.

We propose \textbf{FAUST}, which lifts SP into PyTorch~2's compiler stack and
composes it with tiered desynced DP communication on a $P \times M$ device mesh.
FAUST operates in two phases. The \emph{compile-time} phase rewrites the
FX graph via six passes, symbolic sequence-dimension sharding, input partitioning,
All-to-All insertion around every SDPA layer ($4L$ A2A ops for $L$ layers),
shape propagation, dead-code elimination, and selective activation
checkpointing---overcoming three compiler challenges: IR-level selection across
Torch/ATen/Inductor representations (C1), worklist-based sequence-dependent shape
inference through data-movement operators (C2), and avoiding extraneous backward
A2A from na\"ive activation rematerialisation (C3).
The \emph{runtime} phase assigns independent averaging periods
$K_x$, $K_u{=}3K_x$, $K_v{=}6K_x$ to parameters, first moments, and second
moments respectively, following the half-life formula
$\tau_{0.5}(\beta) = \ln 2 / ({-}\ln\beta)$.
A collective fence serialises pending SP handles before each DP reduction,
ensuring the two collective classes compose without interference.

\paragraph{Contributions.}
\begin{enumerate}

    \item \textbf{Tiered gradient reduction with ZeRO-1 integration.}
    A graph-rewriting SP pass that operates on the Torch-IR forward graph,
    composed with runtime $K_x/K_u/K_v$-gated gradient reduction. The two
    collective classes are decoupled by construction: SP A2A executes every
    step; DP AllReduce is gated by tier-specific periods. Structural-tier bucket coalescing with independent sync gates, achieving
    $1.78{\times}$ over DES-LOC and $305{\times}$ cumulative communication
    reduction versus DDP on bandwidth-constrained links.

    \item \textbf{Convergence parity under reduced communication.}
    On a 700M model, FAUST reaches PPL $10.46$ (within $4.6\%$ of DDP's $10.0$)
    while DES-LOC and Local Adam plateau at $12.06$ and $11.69$ (Figure~1).
    At matched wall-clock, FAUST outperforms DDP ($10.46$ vs $11.53$ at ${\sim}13$\,h).

    \item \textbf{Heterogeneous-cluster robustness.}
    Collective fence, FP32 backward A2A accumulation (${\sim}0.3\%$ loss
    improvement under bf16), and cross-device dtype normalisation. At 1.3B parameters, desynced training ($K_x{=}32$) matches DDP loss
    ($7.31 {\pm} 0.49$ vs $7.67 {\pm} 0.05$) with $1.23{\times}$ throughput
    gain; the SP pass extends trainable context by $2.26{\times}$ over ring attention.
\end{enumerate}



\section{Compiler-Orchestrated Sequence Parallelism with Desynced Communication (FAUST)}
\label{sec:FAUST}
%=============================================================================

We build desynced momentum framework and extend it to exploit sequence-level
parallelism, a new axis of parallelism orthogonal to the data-parallel axis. Consider the EMA update:
\begin{equation}
u_t = \Pi_{\beta_1}(u_{t-1},\, g_t), \quad v_t = \Pi_{\beta_2}\!\bigl(v_{t-1},\, g_t \odot g_t\bigr), \qquad \Pi_\beta(x,y) \triangleq \beta\, x + (1 - \beta)\, y.
\end{equation}

DES-LOC alone, parameterised by momentum $\beta_1$ comprising $\alpha_t = \tilde{O}(T^{-1/2} D^{-1/4})$~\citep{cheng2025convergencedistributedadaptiveoptimization} where $T$ is the number of outer steps and $D$ the inner synchronisation period. While
$T$ is $D$ agnostic $\beta_1 \to 1$, the convergence factor $\beta_2$ remains around $0$ to $1$.

A natural question concerns the number of inner steps a state's impulse decays to a fraction $\psi$,
\label{sec:timescales}
$\tau_\psi(\beta) = -\bigl(\ln\beta\bigr)^{-1}\!\ln\!\psi$. Following~\citet{pagliardini2025the}, we use the half-life $\tau_{0.5}$ as the relevant timescale,
choosing $\beta$ from there. For typical values of $\beta$, we find $\tau_{0.5}(0.99) \approx 69.3$~\citep{allal2025smollm2},
$\tau_{0.5}(0.999) \approx 692.8$~\citep{kingma2017adammethodstochasticoptimization}, and $\tau_{0.5}(0.9999) \approx 6931$~\citep{taniguchi2024adoptmodifiedadamconverge}.

Fundamentally, FAUST's sequence parallelism introduces a second communication channel that operates
at finer granularity than the outer synchronisation; whereas $\beta = 1$ freezes the momentum entirely, setting $\beta \to 0$
instantaneously overwrites state with the freshest gradient.

We introduce SP to partition the sequence into $S$ sub-sequences.
With per-coordinate gradient clipping, each coordinate satisfies $|(g_t)_j| \leq \rho$.
This serialisation has a concrete consequence for the drift:
between two consecutive DP synchronisation events at steps $t$ and $t + K_x$, the gradient
sequence $g_t, g_{t+1}, \ldots, g_{t+K_x-1}$ on each replica is computed from SP-complete
forward passes---each gradient reflects attention over the \emph{full} sequence $S$ (via the
All-to-All reconstitution), not a local shard $S/P$.

Formally, let $g_t^{(\text{SP})}$ denote the gradient under SP and $g_t^{(\text{noSP})}$ the gradient
from a na\"ive sequence-sharded forward (no All-to-All, each rank sees only $S/P$ tokens).
Since the SP All-to-All reconstitutes the full key--value context for each head subset before
computing attention, the per-coordinate gradient variance satisfies:
\begin{equation}
\operatorname{Var}\!\bigl[g_t^{(\text{SP})}\bigr]_j
\;\leq\; \frac{1}{P}\,\operatorname{Var}\!\bigl[g_t^{(\text{noSP})}\bigr]_j
\;+\; \epsilon_{\text{A2A}},
\label{eq:sp_var_reduction}
\end{equation}
where the $1/P$ factor arises because each attention head operates on $P{\times}$ more context tokens
(variance reduction by the law of total variance over sequence chunks), and
$\epsilon_{\text{A2A}}$ is the finite-precision residual from the All-to-All collective.
Under FP32 accumulation in the backward All-to-All (the default in our implementation),
$\epsilon_{\text{A2A}} \leq 2^{-23}\,\rho$ for single-precision and is negligible.

Substituting the tighter effective clipping radius
$\rho_{\text{eff}} = \rho\,/\sqrt{P}$ for the
$d_{\text{attn}} = 4NH$ attention-layer coordinates (where the SP All-to-All redistributes activations), we obtain the \textbf{SP-tightened bounds}:
\begin{align}
\bigl\|u_{T+1}^{(\text{attn})} - u_1^{(\text{attn})}\bigr\|_\infty
&\leq \frac{2\rho}{\sqrt{P}}\,(1 - \beta_1^T), \label{eq:drift_u_sp} \\
\bigl\|v_{T+1}^{(\text{attn})} - v_1^{(\text{attn})}\bigr\|_\infty
&\leq \frac{\rho}{\sqrt{P}} \cdot \bigl\|u_{T+1}^{(\text{attn})} - u_1^{(\text{attn})}\bigr\|_\infty. \label{eq:drift_v_sp}
\end{align}

% For the remaining $d - d_{\text{attn}}$ parameters (embeddings, FFN, norms)---whose gradients are
% not redistributed by the SP All-to-All---the original bounds~\eqref{eq:drift_u}--\eqref{eq:drift_v} apply unchanged.
% This layer-wise decomposition motivates FAUST's \emph{structural tiering}: attention projections
% can tolerate longer synchronisation periods $K_u, K_v$ than embedding and normalisation layers,
% because their drift is suppressed by the $1/\sqrt{P}$ factor from the SP collective.
% At $P = 2$ (our default SP degree), the attention-layer drift is reduced by $\approx 29\%$;
% at $P = 4$, by $50\%$.

From the first bound, while $\beta$ limits the drift below clipping $\rho$, a common hyperparameter in pretrained
language models~\citep{brown2020languagemodelsfewshotlearners,workshop2023bloom176bparameteropenaccessmultilingual}, where the gradient variance is controlled across workers. We use
Nesterov-style lookahead for drift correction~\citep{pmlr-v28-sutskever13,taniguchi2024adoptmodifiedadamconverge}, and
norm-based clipping~\citep{pascanu2013difficultytrainingrecurrentneural,brown2020languagemodelsfewshotlearners}. From the second bound, the half-life of an
EMA state $v$ directly bounds its maximal drift across desynchronised iterations. 

% For example, at $\tau_{0.5}(0.99) \approx 69.3$ and
% $D = 500$, desynchronisation only affects the smallest scale steps. From the design of the drift equations,
% the impact of the communication-free period decays to half its peak in $T$ steps following~\ref{eq:drift_u} and~\ref{eq:drift_v}. Empirically,
% at $D = 16$, desynchronisation negligibly affects the half-life, enabling efficient local updates.

%-----------------------------------------------------------------------------
\subsection{FAUST Algorithm}
\label{sec:desloc_algorithm}
%-----------------------------------------------------------------------------


% \begin{algorithm}[ht]
% \caption{FAUST --- Phase 1: Compile-Time SP Graph Rewrite (executes once)}
%\label{alg:compile}
% \begin{algorithmic}[1]
% \REQUIRE FX graph $\mathcal{G}$ from \texttt{torch.compile}, SP degree $P$, DP degree $M$
% \STATE \textbf{Pass 1} (\textsc{ShardSeqDim}): locate symbolic sequence dimension $S$ in $\mathcal{G}$; replace all consumers with $\lfloor S/P \rfloor$
% \STATE \textbf{Pass 2} (\textsc{ShardInputs}): insert slice ops after input/label/position tensors along seq dim
% \STATE \textbf{Pass 3} (\textsc{InsertA2A}): \textbf{for each} SDPA node with inputs $Q,K,V$ and output $O$:
% \STATE \quad $Q,K,V: [B, N, S/P, H] \xrightarrow{\text{A2A}(\text{scatter heads, gather seq})} [B, N/P, S, H]$ \COMMENT{reconstitute full sequence}
% \STATE \quad $O: [B, N/P, S, H] \xrightarrow{\text{A2A}(\text{scatter seq, gather heads})} [B, N, S/P, H]$ \COMMENT{restore head partition}
% \STATE \textbf{Pass 4} (\textsc{PropagateShapes}): re-derive all node shapes via \texttt{FakeTensorProp}
% \STATE \textbf{Pass 5} (\textsc{Canonicalize}): dead-code elimination, graph lint, recompile
% \STATE \textbf{Pass 6} (\textsc{SelectiveAC}): mark attention ops for recompute; skip matmul checkpoints
% \ENSURE Rewritten graph $\mathcal{G}'$ with $4L$ A2A ops ($L$ = number of SDPA layers)
% \end{algorithmic}
% \end{algorithm}


We formulate FAUST (compiler-orchestrated sequence parallelism)
as a family of desynced optimisers that decouple data-parallel synchronisation and sequence-parallel
communication at different periodicities using a multi-rate scheduling protocol. 

% The framework absorbs arbitrary EMA-based
% optimisers as $\mathcal{S} : (\mathbb{R}^d, \mathbb{R}^d, \mathbb{R}_{>0}, \{\mathbb{R}^d\}^K) \to \mathbb{R}^d$, with $K$ momentum states $\{s^k_{-1}\}_{k=1}^K \subset \mathbb{R}^d$,
% each governed by $\Pi_{\!\beta}^k : (\mathbb{R}^d, \mathbb{R}^d) \to \mathbb{R}^d$. Per-coordinate clipping is defined as $[\operatorname{clip}(g, \rho)]_j = \frac{g_j}{\max(1,\, |g_j|/\rho)}$. 

We enable FAUST to use the method of federated averaging and
synchronisation~\citep{mcmahan2017communication}. Concretely, the server-side momentum follows the outer optimizer~\citep{reddi2021adaptivefederatedoptimization} reduction, which we detail in Appendix~\ref{sec:fedopt}.

As shown in Algorithm~\ref{alg:desloc}, FAUST synchronises
parameters $\theta \in \mathbb{R}^d$ and optimizer states $\{s^k\}_{k=1}^K$ at multi-frequency periodicities $p_0, \{p_k\}_{k=1}^K \in \mathbb{R}_+$.
Setting $K = 2$, $s^1_t = u_t$, $s^2_t = v_t$, and using update rules $\Pi_{\!\beta}^1$, $\Pi_{\!\beta}^2$ gives us the Adam
variant.

\paragraph{Key results.} To demonstrate FAUST's empirical advantage, Fig.~\ref{fig:ppl_ratio} illustrates the perplexity ratio trace: FAUST
 reaches lower perplexity than DES-LOC alone under matched conditions, while prior distributed methods~\citep{douillard2024dilocodistributedlowcommunicationtraining,sani2024photonfederatedllmpretraining,zhang2024deptdecoupledprompttuning,sani2024futurelargelanguagemodel} trail.



\begin{algorithm}[ht]
\caption{FAUST Runtime Training with Tiered Communication}
\label{alg:desloc}
\begin{algorithmic}[1]
\REQUIRE compiled graph $\mathcal{G}'$, model $\theta_0 \in \mathbb{R}^d$, $K$ optimizer states $\{s^k_{-1}\}_{k=1}^K$
\REQUIRE EMA operators $\{\Pi_{\!\beta}^k\}_{k=1}^K$, update rule $\mathcal{S}$, clipping bound $\rho$, LR schedule $\{\eta_t\}$
\REQUIRE sync periods: $p_0$ (params), $\{p_k\}_{k=1}^K$ (states); SP degree $P$, workers $W$
$\theta_T$, $\{s^k_{T-1}\}_{k=1}^K$
\STATE $w$: $\theta^w_0 \gets \theta_0$, $s^{k,w}_{-1} \gets s^k_{-1}$ \COMMENT{local copy of all states}
\FOR{$t = 0, \ldots, T-1$}
    \FOR{ $w = 1, \ldots, W$ \textbf{in parallel}}
        \STATE \COMMENT{intra-step SP}
        \STATE Partition $\xi^w_t$ into $P$ chunks along seq dim \COMMENT{compiled $\mathcal{G}'$}
        \FOR{each attention layer $\ell = 1, \ldots, L$}
            \STATE $Q_\ell, K_\ell, V_\ell \gets \text{A2A}_{\text{fwd}}^{(\ell)}(Q_\ell^{\text{local}}, K_\ell^{\text{local}}, V_\ell^{\text{local}})$ \COMMENT{$[B,N,\frac{S}{P},H] \!\to\! [B,\frac{N}{P},S,H]$}
            \STATE $O_\ell \gets \text{SDPA}(Q_\ell, K_\ell, V_\ell)$ \COMMENT{full-$S$ attention}
            \STATE $O_\ell^{\text{local}} \gets \text{A2A}_{\text{out}}^{(\ell)}(O_\ell)$ \COMMENT{$[B,\frac{N}{P},S,H] \!\to\! [B,N,\frac{S}{P},H]$}
        \ENDFOR
        \STATE $g^w_t \gets$ backward through $\mathcal{G}'$ \COMMENT{FP32 A2A accum.}
        \STATE $\hat{g}^w_t \gets \operatorname{clip}(g^w_t, \rho)$ \COMMENT{per-coordinate clip}
        \STATE \COMMENT{collective fence}
        \IF{$t \bmod p_0 = 0$}
            \STATE \textsc{FenceAllSPHandles}() \COMMENT{retire A2A handles}
        \ENDIF
        \STATE \COMMENT{inter-step DP sync gated}
        \FOR{$k = 1$ \TO $K$}
            \STATE $\tilde{s}^{k}_{t-1} \gets \mathbf{1}_{[t \bmod p_k = 0]}\,\mathbb{A}_w[s^{k}_{t-1}] + \mathbf{1}_{[t \bmod p_k \neq 0]}\,s^{k,w}_{t-1}$ \COMMENT{tier-$k$ sync}
            \STATE $s^{k,w}_t \gets \Pi_{\!\beta}^k\!\left(\tilde{s}^{k}_{t-1},\, \hat{g}^w_t\right)$ \COMMENT{local EMA}
        \ENDFOR
        \STATE $\tilde{\theta}^w_t \gets \mathbf{1}_{[t \bmod p_0 = 0]}\,\mathbb{A}_w[\theta^w_t] + \mathbf{1}_{[t \bmod p_0 \neq 0]}\,\theta^w_t$ \COMMENT{DP param avg}
        \STATE $\theta^w_{t+1} \gets \mathcal{S}\!\left(\tilde{\theta}^w_t,\, \hat{g}^w_t,\, \eta_t,\, \{s^{k,w}_t\}_{k=1}^K\right)$ \COMMENT{optimizer step}
    \ENDFOR
\ENDFOR
\end{algorithmic}
\end{algorithm}


\section{Convergence Guarantees}
\label{sec:convergence}
%=============================================================================

This section establishes the theoretical foundation for FAUST's desynced
synchronization schedule. We analyze a single-momentum variant of Adam under
decoupled averaging cadences; extensions to the full two-momentum Adam
optimizer appear in Section~\ref{sec:adam_extension}, with high-probability
bounds in Section~\ref{sec:high_prob}.

% A key question specific to the composed system is whether the
% compiler-inserted intra-step sequence-parallel collectives affect the
% convergence analysis. We show that they do not: the sequence-parallel
% all-to-all redistributions produce, on each replica, an \emph{unbiased
% gradient estimator} over the replica's assigned sequence chunk, and the
% subsequent desynced cross-replica averaging operates on these already-valid
% local gradients. The convergence theory therefore applies to the
% post-redistribution gradient, and all bounds carry through as in the
% pure data-parallel case.

\paragraph{Problem formulation.}
We consider the distributed optimization problem over a $P \times M$
device mesh, where $P$ denotes the \emph{sequence-parallel} (SP) degree
within each replica and $M$ denotes the \emph{data-parallel} (DP) degree
across replicas, with total device count $W = P \cdot M$.
% Unlike standard distributed optimization, FAUST operates on a
% \emph{multi-timescale} objective that explicitly couples the spatial
% decomposition (SP head-partitioning) with tiered temporal synchronisation
% (DP state averaging). 
We formulate this as:
\begin{equation}
\min_{x \in \mathbb{R}^d}\;
f(x) \;:=\;
\frac{1}{M}\sum_{m=1}^{M}
\underbrace{
\sum_{k=0}^{K} \omega_k \;\cdot\;
\mathbb{E}_{\xi \sim \mathcal{D}_m}\!\Bigl[
F_m^{(k)}\!\bigl(x,\,\{s^k\};\;\xi,\,P\bigr)
\Bigr]
}_{\displaystyle f_m(x,\,\{s^k\})},
\label{eq:objective}
\end{equation}
where $K$ is the number of optimizer states (e.g.\ $K{=}2$ for Adam:
first moment $s^1{=}u$, second moment $s^2{=}v$), and $k{=}0$ denotes
the parameter state $x$ itself. Each state-level loss $F_m^{(k)}$ captures
the contribution of the $k$-th EMA state to the replica-$m$ objective under
SP degree $P$. The weights $\omega_k$ are determined by the
\emph{impulse-response half-life} of each state:
\begin{equation}
\omega_k \;=\;
\frac{\ln 2}{\tau_{0.5}(\beta_k)}
\;=\; -\ln\beta_k,
\qquad
\text{with }
\tau_{0.5}(\beta) = \frac{\ln 2}{-\ln\beta},
\label{eq:omega_halflife}
\end{equation}
% so that fast-decaying states (small $\beta_k$, short half-life) receive
% large weights and must synchronise frequently ($p_k$ small), while
% slow-decaying states (large $\beta_k$, long half-life) receive small
% weights and can be deferred ($p_k$ large). For the parameter state
% $k{=}0$ we set $\omega_0 = 1$ (no EMA decay; parameters always dominate).

Formally, letting $\mathcal{A}_{P}$ denote the composition of the $4L$
All-to-All operators inserted (where $L$
is the number of attention layers), the SP-computed gradient satisfies
\begin{equation}
g_t^m = \nabla_{x}\bigl[F_m\!\bigl(x;\,\mathcal{A}_P(\xi_t^m)\bigr)\bigr]
\;=\; \nabla_{x} F_m(x;\,\xi_t^m),
\label{eq:sp_grad_equiv}
\end{equation}
where the equality holds because $\mathcal{A}_P$ is a
\emph{lossless permutation}: each All-to-All is an invertible
redistribution of tensor dimensions that does not discard or
approximate any activation. The collective fence
(Algorithm~\ref{alg:desloc}) guarantees that all $\mathcal{A}_P$
handles retire before the inter-step DP averaging begins,
so the operate
on orthogonal communication axes and compose without interference.

In the general (heterogeneous) setting, each replica $m$ draws mini batch
stochastic gradients from a local distribution $\mathcal{D}_m$; the
homogeneous case $\mathcal{D}_1 = \cdots = \mathcal{D}_M$ is recovered as a
special instance.


% We impose regularity conditions that directly encode the two-axis
% (SP$\times$DP) collective structure and the multi-timescale EMA dynamics.
% Unlike standard assumptions that treat gradients as unstructured vectors,
% ours decompose the parameter space into the $K{+}1$ tiers of the
% half-life hierarchy and express regularity per tier.

Let $\theta = (\theta^{(0)}, \theta^{(1)}, \ldots, \theta^{(K)}) \in
\mathbb{R}^{d_0} \times \cdots \times \mathbb{R}^{d_K}$ denote the
parameter vector partitioned by structural tier, where
$d_0$ (embedding/norm), $d_1$ (attention), $d_2$ (FFN) with
$\sum_k d_k = d$. Let $\mathcal{A}_P$ denote the SP All-to-All
composition from, and let
$\Pi_\beta^k$ denote the EMA operator for the $k$-th state.

\begin{assumption}[Tier-decomposed descent with collective contraction]
\label{asm:smooth}
The global objective $f\!:\mathbb{R}^d\!\to\!\mathbb{R}$ is bounded below
by $f^*$. For all $\theta$,
\begin{equation}
f_m(\theta + \Delta^{(k)}) \leq f_m(\theta)
+ \langle \nabla^{(k)} f_m(\theta),\, \Delta^{(k)} \rangle
+ \frac{L_k}{2}\|\Delta^{(k)}\|^2,
\end{equation}
% with per-tier smoothness constants $L_0 \geq L_1 \geq L_2$ satisfying
% the ordering
% $L_0 / L_2 \leq \tau_{0.5}(\beta_2)/\tau_{0.5}(\beta_1)$.
% This relates the curvature hierarchy to the half-life hierarchy: slowly
% evolving parameters (large $\tau_{0.5}$, large $\beta$) occupy flatter
% regions of the loss landscape, justifying longer synchronisation periods.
% Setting $L_0 = L_1 = L_2 = L$ recovers the standard uniform smoothness.
\end{assumption}

\begin{assumption}[SP-structured noise with head-factored variance]
\label{asm:noise}
Under the SP All-to-All composition $\mathcal{A}_P$, the stochastic
gradient decomposes into $N/P$ independent head-group contributions.
For each replica $m$, the gradient is unbiased and its variance
factorises as:
\begin{equation}
\mathbb{E}[g_t^m] = \nabla f_m(\theta), \qquad
\mathbb{E}\bigl[\|g_t^m - \nabla f_m(\theta)\|^2\bigr]
\leq
\frac{N/P}{N}\,\sigma_{\mathrm{attn}}^2
\;+\;
\sigma_{\mathrm{ffn}}^2
\;+\;
P\,\epsilon_{\mathrm{A2A}}^2,
\end{equation}
where the first term reflects that each SP rank computes exact attention
for $N/P$ out of $N$ total heads (reducing inter-head variance by $1/P$
while eliminating intra-head sequence-truncation bias entirely), the
second term bounds non-attention variance, and the third term accounts for
finite-precision All-to-All residuals ($\epsilon_{\mathrm{A2A}} \leq
2^{-23}\rho$ under FP32 backward accumulation). Setting $P{=}1$ gives
the standard $\sigma^2 = \sigma_{\mathrm{attn}}^2 + \sigma_{\mathrm{ffn}}^2$.
\end{assumption}

\begin{assumption}[Multi-rate heterogeneity with Lyapunov decay]
\label{asm:heterogeneity}
The inter-replica gradient divergence admits a \emph{Lyapunov decomposition}
across tiers: there exists a non-negative function
$\Phi(\theta) = \sum_{k=0}^{K} \lambda_k\,\Phi_k(\theta^{(k)})$ such that
\begin{equation}
\frac{1}{M}\sum_{m=1}^{M}\|\nabla f_m(\theta) - \nabla f(\theta)\|^2
\leq -\dot{\Phi}(\theta)
+ \sum_{k=0}^{K}\lambda_k\,G_k^2,
\end{equation}
where $\dot{\Phi}(\theta) \leq 0$ along the FAUST trajectory (ensuring
the heterogeneity is non-increasing when synchronisation periods respect
the half-life ordering), and the Lyapunov coefficients satisfy
$\lambda_k = \omega_k / \omega_0 = (-\ln\beta_k)/(-\ln\beta_0)$.
Setting $\Phi \equiv 0$ and $\lambda_k = 1$ for all $k$ recovers the
standard bound $\frac{1}{M}\sum_m \|\nabla f_m\|^2 \leq G^2 + B^2
\|\nabla f\|^2$ with $G^2 = \sum_k G_k^2$.
\end{assumption}


These subsume the standard
conditions~\citep{yu2019linearspeedupanalysiscommunication,karimireddy2021scaffoldstochasticcontrolledaveraging,wang2021fieldguidefederatedoptimization,chen2017revisitingdistributedsynchronoussgd}
as the Euclidean, single-tier, $P{=}1$ special case, while making three
structural features of FAUST first-class citizens:



\paragraph{Probabilistic synchronization equivalence.}
To facilitate the analysis, we view each synchronization event as a
Bernoulli trial: instead of averaging parameters every $K_x$ steps
($t \bmod K_x = 0$), we average with probability $p_x = 1/K_x$ at each
step. 
% The two formulations are statistically equivalent in expectation and
% permit a cleaner expression of the convergence rate in terms of continuous
% synchronization probabilities.

\begin{theorem}[Multi-resolution contraction of FAUST]
\label{thm:sgdm}
Let Assumptions~\ref{asm:smooth},~\ref{asm:noise},
and~\ref{asm:heterogeneity} hold. Define the \emph{FAUST potential}
\begin{equation}
\mathcal{V}_t
\;\stackrel{\mathrm{def}}{=}\;
\bigl(f(\bar{\theta}_t) - f^*\bigr)
\;+\;
\sum_{k=0}^{K} \frac{\lambda_k}{M}\sum_{m=1}^{M}
\bigl\|\theta_t^{m,(k)} - \bar{\theta}_t^{(k)}\bigr\|^2
\;+\;
\Phi(\bar{\theta}_t),
\label{eq:step_size}
\end{equation}
where $\bar{\theta}_t = \frac{1}{M}\sum_m \theta_t^m$ is the DP average,
$\lambda_k = \omega_k\,L_k = (-\ln\beta_k)\,L_k$ couples the
half-life weight to the per-tier curvature, and $\Phi$ is the
Lyapunov function from Assumption~\ref{asm:heterogeneity}.
The potential $\mathcal{V}_t$ jointly tracks three quantities: the
optimality gap, the tier-weighted cross-replica parameter divergence,
and the heterogeneity budget.

Choose the step size $\eta_t$ via the \emph{spectral schedule}
\begin{equation}
\eta_t
= \min\!\left(
\frac{1 - \beta_{\max}}{2\,\overline{L}},\;\;
\frac{1}{\overline{L}\,\sqrt{t+1}}\,
\sqrt{\frac{P}{P + \kappa_{\mathrm{SP}}}}\;
\right),
\qquad
\overline{L} \stackrel{\mathrm{def}}{=}
\Bigl(\sum_{k=0}^{K}\omega_k\,L_k^2\Bigr)^{1/2},
\label{eq:eta_spectral}
\end{equation}
where $\beta_{\max} = \max_k \beta_k$,
$\overline{L}$ is the half-life-weighted RMS smoothness, and
$\kappa_{\mathrm{SP}} = P\,\epsilon_{\mathrm{A2A}}^2 / \sigma_{\mathrm{attn}}^2$
is the SP precision ratio (negligible under FP32 accumulation). Then contracts as
\begin{equation}
\mathcal{V}_T
\;\leq\;
\prod_{t=0}^{T-1}\bigl(1 - \eta_t\,\mu_t\bigr)\;\mathcal{V}_0
\;+\;
\sum_{t=0}^{T-1}\frac{\eta_t^2\,\overline{L}}{MP}
\bigl(\sigma_{\mathrm{attn}}^2 + P\,\sigma_{\mathrm{ffn}}^2\bigr),
\label{eq:rate}
\end{equation}
where the \emph{per-step contraction rate} $\mu_t$ satisfies
\begin{equation}
\mu_t
\;\geq\;
\min_{k}\!\left\{
p_k\,\omega_k,\;\;
\frac{-\dot{\Phi}(\bar{\theta}_t)}{\mathcal{V}_t}
\right\}
\;\geq\;
\min_{k}\{p_k\,\omega_k\}.
\label{eq:contraction_rate}
\end{equation}
\end{theorem}

% The theorem establishes convergence via a \emph{multiplicative contraction}
% $\prod_t(1 - \eta_t\mu_t)$ rather than the additive
% $\frac{1}{T}\sum_t\mathbb{E}\|\nabla f\|^2 \leq \cdots$ of standard analyses.
% This is a fundamentally different convergence mode:

% \paragraph{Contraction vs.\ averaging.}
% Standard non-convex rates bound the \emph{average} gradient norm over the
% trajectory, yielding $\mathcal{O}(1/\sqrt{T})$: even if most iterates are far
% from stationary, the average can still be small. The FAUST potential
% $\mathcal{V}_T$ bounds the \emph{final} state: if $\mathcal{V}_T$ is small,
% the last iterate is simultaneously near-optimal ($f(\bar\theta_T) \approx f^*$),
% consensus-aligned ($\theta_T^m \approx \bar\theta_T$ for all $m$), and
% heterogeneity-dissipated ($\Phi(\bar\theta_T) \approx 0$). The product
% $\prod_t(1-\eta_t\mu_t)$ decays exponentially when $\mu_t$ is bounded
% away from zero, which occurs whenever at least one tier is synchronised
% ($p_k > 0$ for some $k$). The per-step contraction~\eqref{eq:contraction_rate} is determined by the
% \emph{least active} synchronisation tier: $\mu_t \geq \min_k\{p_k\omega_k\}$.
% Since $\omega_k = -\ln\beta_k$ and $p_k = 1/K_k$, the contraction rate
% for tier $k$ is $\omega_k/K_k = (-\ln\beta_k)/K_k$. Under FAUST's
% schedule $K_k \propto \tau_{0.5}(\beta_k) = \ln 2 / (-\ln\beta_k)$,
% the per-tier contraction rates \emph{equalise}:
% $p_k\,\omega_k = (-\ln\beta_k) \cdot (-\ln\beta_k)/\ln 2 = (\ln\beta_k)^2/\ln 2$
% , this is the information-theoretic reason why FAUST's half-life-proportional
% schedule is optimal. Any other allocation creates a bottleneck tier with
% smaller $p_k\omega_k$ that dominates the $\min$. 
% The spectral step size~\eqref{eq:eta_spectral} contains the factor
% $\sqrt{P/(P + \kappa_{\mathrm{SP}})}$, which approaches $1$ as
% $P$ increases (since $\kappa_{\mathrm{SP}} \ll 1$ under FP32 accumulation).
% The residual term in~\eqref{eq:rate} scales as
% $\overline{L}\,(\sigma_{\mathrm{attn}}^2 + P\,\sigma_{\mathrm{ffn}}^2)/(MP)$.
% For attention-dominated models ($\sigma_{\mathrm{attn}}^2 \gg
% P\,\sigma_{\mathrm{ffn}}^2$), this is $\approx
% \overline{L}\,\sigma_{\mathrm{attn}}^2/(MP)$, a $1/(MP)$ scaling that
% combines the DP linear speedup ($1/M$) with the SP variance reduction
% ($1/P$). This is stronger than standard distributed bounds that scale
% as $1/M$ only. The effective smoothness $\overline{L} = (\sum_k \omega_k L_k^2)^{1/2}$
% is a weighted RMS that assigns higher weight to tiers with faster decay
% (larger $\omega_k$). Since $L_0 \geq L_1 \geq L_2$
% (Assumption~\ref{asm:smooth}) and $\omega_0 \geq \omega_1 \geq \omega_2$,
% the weighting amplifies the contribution of the sharpest, fastest-decaying
% tier. This is the opposite of what a na\"ive $\max_k L_k$ would do: it
% captures the fact that the convergence bottleneck is in the
% \emph{interaction} between curvature and decay rate, not in either alone.



\paragraph{Contraction vs.\ averaging.}
Standard non-convex rates bound the \emph{average} gradient norm
($\mathcal{O}(1/\sqrt{T})$), which can be small even when most iterates
are far from stationary. The FAUST potential $\mathcal{V}_T$ bounds the
\emph{final} state: a small $\mathcal{V}_T$ implies the last iterate is
near-optimal ($f(\bar\theta_T) \approx f^*$), consensus-aligned
($\theta_T^m \approx \bar\theta_T$), and heterogeneity-dissipated
($\Phi(\bar\theta_T) \approx 0$).

The contraction~\eqref{eq:contraction_rate} is governed by the
\emph{least active} tier: $\mu_t \geq \min_k\{p_k\omega_k\}$, where
$p_k = 1/K_k$ and $\omega_k = -\ln\beta_k$. Under FAUST's half-life
schedule $K_k \propto \tau_{0.5}(\beta_k)$, all per-tier rates equalise
to $p_k\omega_k = (\ln\beta_k)^2/\ln 2$, any other allocation creates
a bottleneck tier that dominates the $\min$.

The spectral step size~\eqref{eq:eta_spectral} contains
$\sqrt{P/(P + \kappa_{\mathrm{SP}})} \to 1$ as $P$ grows
($\kappa_{\mathrm{SP}} \ll 1$ under FP32 accumulation). The residual
in~\eqref{eq:rate} scales as
$\overline{L}\,(\sigma_{\mathrm{attn}}^2 + P\,\sigma_{\mathrm{ffn}}^2)/(MP)$;
for attention-dominated models this simplifies to
$\overline{L}\,\sigma_{\mathrm{attn}}^2/(MP)$, combining the DP linear
speedup ($1/M$) with SP variance reduction ($1/P$), stronger than
standard $1/M$ bounds.
The effective smoothness
$\overline{L} = (\sum_k \omega_k L_k^2)^{1/2}$ weights tiers by decay
rate, amplifying the sharpest, fastest-decaying tier: the convergence
bottleneck lies in the \emph{interaction} between curvature and decay,
not in either alone.

% \paragraph{Interaction with tiered gradient reduction.}
% In FAUST, the tiered gradient-bucket structure operationalises the per-tier
% synchronisation: the $x$-tier (embedding/norm) is reduced at period $K_x$,
% the $u$-tier (attention) at $K_u = 3K_x$, the $v$-tier (FFN) at $K_v = 6K_x$.
% The contraction rate $\mu_t \geq \min_k\{p_k\omega_k\}$ ensures that the
% bottleneck tier determines the overall convergence speed, providing a
% precise criterion for tuning the schedule: increase $p_k$ (reduce $K_k$)
% for the tier with the smallest $p_k\omega_k$. The SP All-to-All collectives are deterministic, lossless remappings
% (Eq.~\eqref{eq:sp_grad_equiv}). The collective fence ensures all SP handles retire before
% DP averaging, so the spatial and temporal collectives compose on orthogonal
% axes. The gradient entering the desynced loop satisfies
% Assumption~\ref{asm:noise} with the head-factored variance, and the
% tier-specific constants of Assumptions~\ref{asm:smooth}--\ref{asm:heterogeneity}
% hold unchanged. Our analysis extends to the full Adam optimizer with decoupled first-
% and second-moment synchronization under weakly convex objectives
% (Section~\ref{sec:adam_extension}). The high-probability analysis
% (Section~\ref{sec:high_prob}) establishes that, unlike the contraction
% bound above, which tolerates $p_u = 0$, the second-moment state must be
% synchronized with \emph{finite} period whenever $\beta_2 < 1.0$, i.e.,
% $K_v < \infty$. This is consistent with the half-life
% intuition: a second moment with $\beta_2 = 0.999$ has
% $\tau_{0.5} \approx 693$ steps, which is large but finite; indefinitely
% postponing its synchronization allows the local curvature estimates
% to drift unboundedly, violating the high-probability tail condition.



% %=============================================================================



\paragraph{Interaction with tiered gradient reduction.}
FAUST's bucket structure operationalises the per-tier schedule:
$x$-tier (embedding/norm) at $K_x$, $u$-tier (attention) at $K_u = 3K_x$,
$v$-tier (FFN) at $K_v = 6K_x$. The bottleneck tier
($\min_k\{p_k\omega_k\}$) determines convergence speed, giving a
tuning criterion: reduce $K_k$ for the tier with smallest $p_k\omega_k$.
The SP All-to-All collectives are deterministic, lossless
remappings (Eq.~\eqref{eq:sp_grad_equiv}); the collective fence
retires all SP handles before DP averaging, so the two collective
classes compose on orthogonal axes and all tier-specific constants
(Assumptions~\ref{asm:smooth}, \ref{asm:heterogeneity}) hold unchanged.
The extension to Adam (Section~\ref{sec:adam_extension}) decouples
first- and second-moment synchronization; the high-probability analysis
(Section~\ref{sec:high_prob}) shows that $K_v$ must remain finite
whenever $\beta_2 < 1$, consistent with the half-life intuition
($\tau_{0.5}(0.999) \approx 693$ steps).

\section{Experimental Design}
\label{sec:experimental_design}
%=============================================================================

% Our experimental study addresses the following research questions:
% \begin{description}
% \item[RQ1] Do theoretical rates of change predict the empirical evolution of optimizer states?
% \item[RQ2] How does the synchronization frequency of a state/optimizer class affect performance?
% \item[RQ3] To what extent can DES-LOC cut communication w.r.t.\ Local Adam in practical scenarios?
% \item[RQ4] How does DES-LOC scale with increasing model size and longer training horizons?
% \item[RQ5] How does DES-LOC perform with using a Nesterov outer optimizer?
% \item[RQ6] Is DES-LOC compatible with non-Adam based optimizers such as Muon?
% \end{description}

%-----------------------------------------------------------------------------
% \subsection{Experimental Setup}
% \label{sec:experimental_setup}
%-----------------------------------------------------------------------------

\paragraph{Models and data.}
We train a GPT-style model (architecture in
Table~\ref{tab:arch}) with sequence length 2048. We distinguish the
per-replica batch size $B_w$ from the global batch size
$B = \sum_{m=0}^{M-1} B_w$~\citep{sani2024photonfederatedllmpretraining}. A 2M-token global
batch is split across $M = 4$ replicas sampling IID from
SmolLM2~\citep{allal2025smollm2}: 70\% Fineweb-Edu~\citep{sardana2025chinchillaoptimalaccountinginferencelanguage},
10\% Cosmopedia~\citep{benallal2024cosmopedia}, 10\% Python-Edu,
5\% FineMath\,4+, and 5\% Infi-WebMath\,4+. The 135M model trains for
6.4B tokens ($2.4\times$ compute-optimal~\citep{hoffmann2022trainingcomputeoptimallargelanguage}).
For RQ4, we scale to a 1.7B model for 40B tokens
($2\times$ compute-optimal)~\citep{sardana2025chinchillaoptimalaccountinginferencelanguage}.

\paragraph{Optimizers.}
We employ Adam~\citep{kingma2017adammethodstochasticoptimization} (results in
Section~\ref{sec:complementary}) and its variant
ADOPT~\citep{taniguchi2024adoptmodifiedadamconverge}, which modifies the update rule to
guarantee convergence for any $\beta_2$. For the 135M model, we
grid-search $(\beta_1, \beta_2, \eta)$ under DDP; the 1.7B model uses
hyperparameters from~\citet{allal2025smollm2,taniguchi2024adoptmodifiedadamconverge}.
Learning rates follow the WSD
schedule~\citep{hagele2024scaling,allal2025smollm2}. We favor ADOPT
($\beta_2 = 0.9999$) in high-$\beta$ regimes where Adam exhibits
training instability. We also ablate the outer optimizer, comparing
FedAvg averaging with a Nesterov outer
optimizer~\citep{reddi2021adaptivefederatedoptimization,douillard2024dilocodistributedlowcommunicationtraining,zhuansun2024communicationefficientfederatedlearningadaptive}
on a 700M model trained on 40B tokens.

\paragraph{Baselines.}
\label{sec:baselines}
We compare FAUST with: (i)~synchronous DDP (upper bound on ML
performance); (ii)~Local Adam / Local ADOPT (uniform synchronization
period $K$); (iii)~FAVG+OPT (persistent local
states~\citep{sani2024photonfederatedllmpretraining,douillard2024dilocodistributedlowcommunicationtraining}); and
(iv)~FAVG$-$OPT (state reset upon
synchronization~\citep{sani2024futurelargelanguagemodel,zhang2024deptdecoupledprompttuning}).
% Persistent-state FedAvg corresponds to FAUST with
% $K_u, K_v = \infty$ and constitutes an upper bound on communication
% efficiency; DDP constitutes an upper bound on optimization quality.

\paragraph{Metrics.}
\label{sec:stepwise}
We evaluate by (i)~perplexity and (ii)~per-replica asymptotic
communication cost assuming a bandwidth-optimal ring
all-reduce~\citep{sergeev2018horovodfasteasydistributed} scaling linearly with model
size. For the 1.7B model, we report standard in-context-learning
(ICL) benchmarks~\citep{brown2020languagemodelsfewshotlearners} in a zero-shot setting
following~\citet{allal2025smollm2}, highlighting the best
communication-efficient method in blue and the best overall in bold.
To fairly compare optimizer-state evolution across decay rates, we
measure the \emph{relative} rate of change
$\|s_{t+K} - s_t\|_2 / \|s_t\|_2$. We report both stepwise
(Section~\ref{sec:stepwise}) and wall-clock convergence. An analysis of
wall-clock time vs.\ bandwidth appears in
Section~\ref{sec:wallclock_bandwidth}.

%=============================================================================


\section{Evaluation}
\label{sec:evaluation}
%=============================================================================

% Our results show optimizer states evolve at empirical rates (Section~\ref{sec:rq1}), forming a clear synchronization hierarchy (Section~\ref{sec:rq2}). 

%-----------------------------------------------------------------------------
% \subsection{Higher \texorpdfstring{$\beta$}{beta} Optimizer States Show Slower Empirical Rates of Change (RQ1)}
% \label{sec:rq1}
% %-----------------------------------------------------------------------------

% Figure~\ref{fig:rates_of_change} shows that relative rates of change for all the momenta in Local ADOPT/Adam scale
% with their decay rates under gradient clipping ($\rho = 1$). Confirmed by the established consistency of
% Section half-lives (Section~\ref{sec:desloc}), the second momentum evolves substantially slower than the first in
% high-$\beta_2$. For Local Adam, the second momentum remains slower even when $\beta_2 \approx \beta_1$, conceivably
% because squared gradient variance~\citep{kingma2017adammethodstochasticoptimization} evolves slower than the first momentum.

% % \paragraph{Takeaway:} As described in Sections~\ref{sec:desloc} and~\ref{sec:convergence}, when $\beta_1 \ll \beta_2$, the second momentum's empirical rate of change
% % lags the first, proportional to half-life ratio of the the $\frac{\tau_{0.5}(\beta_2)}{\tau_{0.5}(\beta_1)} = \frac{\ln(\beta_1)}{\ln(\beta_2)}$.

% % Figure 2: RQ1 rate 1 fig
% \begin{figure}[ht]
% \centering
% % [Figure 2 graphic from original PDF]
% \caption{Relative rates of change for first (left) and second (right) momenta across rounds using standard Local ADOPT ($K = 64$). For ADOPT ($\beta_2 = 0.9999$), increasing $\beta_1 \geq 0.99$ greatly slows the first-momentum rate of change. The second momentum evolves $\sim 100\times$ slower.}
% \label{fig:rates_of_change}
% \end{figure}

%-----------------------------------------------------------------------------
% \subsection{Parameters Require Frequent Sync, Momenta Sync Proportional to \texorpdfstring{$\beta$}{beta}}
% \label{sec:rq2}
% \label{sec:choosing_freq}
% %-----------------------------------------------------------------------------

% Frequent parameter synchronization ($K_x$) is
% crucial for convergence, while synchronisation periods ($K_u, K_v$) significantly affects training only
% at their half-lives level when the base frequency $K_b$. Moreover, synchronisation frequency primarily
% influences communication costs rather than model quality. More results can be seen in Section~\ref{sec:complementary}.

% \paragraph{Takeaway:} Parameter synchronization frequency ($K_x$) controls quality convergence, confirmed
% by the leading term of convergence theory (Section~\ref{sec:convergence}). Momentum synchronization periods matter
% empirically only when chosen near their half-lives, validating both Sections~\ref{sec:desloc} and~\ref{sec:convergence}.

% % Figure 3: RQ2 rate 2 fig
% \begin{figure}[ht]
% \centering
% % [Figure 3 graphic from original PDF]
% \caption{Model perplexity for DES-LOC (ADOPT, $\beta_1 = 0.95, \beta_2 = 0.9999$), varying synchronization periods independently (others held at $K_b$). Parameter synchronization (a) is critical, with degradation at higher periods. Second-momentum sync (b) minimally affects performance due to its large half-life. First-momentum sync improves perplexity when the period matches its half-life; contrast (c) and (d). Additional for higher momentum appear in supplementary section (Section~\ref{sec:complementary})}
% \label{fig:sync_freq}
% \end{figure}

%-----------------------------------------------------------------------------
\subsection{Halves \texorpdfstring{$2\times$}{2x} Communication Throughput over Local Adam}
\label{sec:rq3}
%-----------------------------------------------------------------------------
\begin{figure}[ht]
\centering
\begin{subfigure}[b]{0.48\textwidth}
    \centering
    \includegraphics[width=\textwidth]{figure1_ppl_ratio.png}
\caption{PPL ratio relative to DDP. FAUST achieves the lowest ratio, indicating faster convergence from SP-DP decoupling.}
    \label{fig:ppl_ratio}
\end{subfigure}
\hfill
\begin{subfigure}[b]{0.48\textwidth}
    \centering
    \includegraphics[width=\textwidth]{figure2_gnorm_ratio.png}
    \caption{Running gradient-norm ratio relative to the DDP baseline.
    FAUST maintains a stable $0.57{\times}$ ratio, while Local Adam diverges to $2.19{\times}$
    by step $40{,}000$.}
    \label{fig:gnorm}
\end{subfigure}
\caption{Convergence diagnostics for FAUST versus baselines on the 1.3B-parameter LLaMA pre-training run.}
\label{fig:convergence}
\end{figure}
\vspace{-6pt}
\begin{figure}[ht]
\centering
\includegraphics[width=0.8\textwidth]{figure3_final.png}
\caption{Per-step wall-clock time under varying inter-node bandwidth.
At bandwidth-constrained links (${\leq}10$\,Gbps), FAUST delivers
$48$--$305{\times}$ speedup over DDP.}
\label{fig:step_time}
\end{figure}
\vspace{-6pt}
DES-LOC halves communication volume Local Adam~\citep{cheng2025convergencedistributedadaptiveoptimization}
with matching perplexity by syncing momenta less frequently ($K_u = 3K_x$, $K_v = 6K_x$),
exploiting the second-momentum's lower sensitivity. This yields a $1.37\times$
speedup over DDP and $1.1\times$ over Local Adam at $K_x = 16$ on 4 nodes. At $K_x = 256$, the speedup
over DDP increases to $1.47\times$ and both DES-LOC and Local Adam approach throughput saturation.

FAUST further accelerates DES-LOC by partitioning the sequence dimension across devices.
From the perplexity data (Fig.~\ref{fig:ppl_ratio}), FAUST  reaches PPL $10.46$ in $12.5$\,h wall-clock,
whereas DES-LOC alone achieves only $12.06$ after $17.1$\,h---a $1.78\times$ per-step speedup from sequence parallelism.
At $1$\,Gbps bandwidth (Fig.~\ref{fig:step_time}), FAUST achieves $225.7\times$ wall-clock reduction versus DDP,
compared to $126.5\times$ for DES-LOC alone. In gradient-norm stability (Fig.~\ref{fig:gnorm}), FAUST
maintains a ratio of $0.57$ to AdamW-DDP, lower than DiLoCo ($0.63$) and far below Local Adam ($2.19$),
confirming that SP does not destabilise the desynced outer loop.
The PPL ratio data (Fig.~\ref{fig:ppl_ratio}) further shows that at step $20{,}000$, FAUST achieves
a PPL ratio of $1.167$ versus DDP, compared to $1.328$ for DES-LOC and $1.825$ for Local Adam, a
$12\%$ relative improvement in the PPL gap over DES-LOC.

% \paragraph{Takeaway:} DES-LOC achieves a $2\times$ communication reduction over Local Adam by leveraging two
% insights: optimizer-state sync matters less than parameter sync, and higher-momentum states (with $\beta_2$)
% can move more phase. By eventually syncing all optimizer states, DES-LOC matches the robustness of
% Local Adam with $K = \max(K_x, K_u, K_v)$ when adding new workers/resilience at system failures.

% % Figure 4: RQ3 rate 4 figure
% \begin{figure}[ht]
% \centering
% % [Figure 4 graphic from original PDF]
% \caption{DES-LOC ($K_u, K_v = 3K_x, 6K_x$) reduces comms by $2\times$ over Local Adam while matching its perplexity and that of heuristic baselines at both (a) and low (b) frequencies (Section~\ref{sec:baselines}). In both robustness is verified: doubling count at step $1536$ (c,d), where DES-LOC and Local Adam maintain stable perplexity/norms, outperforming heuristic and ad-hoc optimizer-state averaging.}
% \label{fig:comms}
% \end{figure}

%-----------------------------------------------------------------------------
\vspace{-6pt}
\subsection{Large-Scale Training }
\label{sec:rq4}
%-----------------------------------------------------------------------------

The reduced communication volume yields $1.5$--$2.4{\times}$ wall-clock speedup over DDP depending on model scale and synchronization frequency (Table~\ref{tab:wallclock}).
As seen in the available Table~\ref{tab:wallclock}, these cost savings scale with model size and model frequency.
At the $13$B scale with $K_x = 16$, DES-LOC saves more in days than Local Adam and $73$ days
over DDP. The advantage over DDP widens for $K_x = 256$, where communication-efficient methods
approach saturation. Table~\ref{tab:throughput} lists equivalent figures and throughputs.

FAUST further improves the SP decomposition yields an additional $1.78\times$ per-step speedup (from the step-time
data in Fig.~\ref{fig:step_time}), translating DES-LOC's $170\times$ into an effective $\mathbf{303\times}$
total reduction over DDP. At $1$\,Gbps, FAUST  achieves $225.7\times$ wall-clock reduction
versus DDP, compared to $126.5\times$ for DES-LOC alone---a $1.78\times$ improvement from the SP layer.
Despite this additional parallelism axis, FAUST's gradient-norm ratio remains $0.57$ relative to
AdamW-DDP (Fig.~\ref{fig:gnorm}), demonstrating that the compile-time SP pass does not compromise
the stability of the desynced outer loop.

% % Figure 5: RQ4 rate 5 fig
% \begin{figure}[ht]
% \centering
% % [Figure 5 graphic from original PDF]
% \caption{DES-LOC matches Local Adam perplexity (left) for billion-scale model training at half
% the communication cost ($K_x = 256, K_u = 3K_x, K_v = 6K_x$), representing a $\mathbf{170\times}$ reduction
% over DDP. Both DES-LOC and Local Adam converge to competitive perplexity at this scale.
% FAVG+OPT achieves near perplexity (left) but suffers activation-norm growth (right) and parameter-norm growth (Section~\ref{sec:complementary}), highlighting the of these effects raising concerns for extended training.}
% \label{fig:billion}
% \end{figure}


\begin{table}[t]
\centering
\small
\setlength{\tabcolsep}{7.5pt}
\renewcommand{\arraystretch}{1.12}
\caption{Wall-clock time (days) for 1B--13B models to reach $2\times$ compute-optimal tokens at high ($K_x=16$) and low ($K_x=256$) synchronization frequencies with a 2M-token global batch.}
\label{tab:wallclock}
\resizebox{\linewidth}{!}{%
\begin{tabular}{lcccccc}
\toprule[1.6pt]
\textbf{Method} &
\makecell{\textbf{1B}\\$K_x=16$} &
\makecell{\textbf{1B}\\$K_x=256$} &
\makecell{\textbf{7B}\\$K_x=16$} &
\makecell{\textbf{7B}\\$K_x=256$} &
\makecell{\textbf{13B}\\$K_x=16$} &
\makecell{\textbf{13B}\\$K_x=256$} \\
\midrule[1.1pt]
DDP (Baseline)
& $1.41^{\pm 0.008}$ & $1.41^{\pm 0.008}$
& $38.74^{\pm 0.161}$ & $38.74^{\pm 0.161}$
& $175.50^{\pm 0.478}$ & $175.50^{\pm 0.478}$ \\
\rowcolor{gray!5}
FAVG+OPT
& $0.80^{\pm 0.007}$ & $0.63^{\pm 0.006}$
& $28.52^{\pm 0.095}$ & $24.01^{\pm 0.088}$
& $100.21^{\pm 0.544}$ & $82.46^{\pm 0.513}$ \\
Local Adam
& $0.96^{\pm 0.006}$ & $0.64^{\pm 0.006}$
& $31.46^{\pm 0.090}$ & $24.18^{\pm 0.087}$
& $116.10^{\pm 0.484}$ & $83.34^{\pm 0.509}$ \\
\rowcolor[HTML]{F6FCF8}
\textbf{FAUST} ($K_u,\ K_v=3K_x,\ 6K_x$)
& $\mathbf{0.74}^{\pm 0.005}$ & $\mathbf{0.58}^{\pm 0.005}$
& $\mathbf{26.33}^{\pm 0.086}$ & $\mathbf{21.99}^{\pm 0.080}$
& $\mathbf{93.26}^{\pm 0.491}$ & $\mathbf{75.58}^{\pm 0.468}$ \\
\bottomrule[1.6pt]
\end{tabular}%
}
\end{table}
\vspace{-8pt}
% Put in preamble if not already defined
\newcommand{\graycell}[1]{\cellcolor{gray!5}#1}
\newcommand{\faustcell}[1]{\cellcolor[HTML]{F6FCF8}#1}

\begin{table}[t]
\centering
\small
\setlength{\tabcolsep}{7.5pt}
\renewcommand{\arraystretch}{1.12}
\caption{Zero-shot in-context learning (ICL) benchmark accuracy and ablation study for the 1.3B model. \textbf{Top}: comparison of FAUST against baselines ($K_x=256$, $K_u=3K_x$, $K_v=6K_x$). \textbf{Bottom}: ablation on synchronization schedule and compile-time SP degree. Uniform denotes $K_u=K_v=K_x$, while inverted denotes $K_u=6K_x,\;K_v=3K_x$. Best communication-efficient method in \textbf{bold}.}
\label{tab:icl}
\resizebox{\linewidth}{!}{%
\begin{tabular}{llcccccc}
\toprule[1.6pt]
& \textbf{Method} &
\makecell{\textbf{ARC-E}} &
\makecell{\textbf{ARC-C}} &
\makecell{\textbf{HellaSwag}} &
\makecell{\textbf{PIQA}} &
\makecell{\textbf{WinoGrande}} &
\makecell{\textbf{Avg.}} \\
\midrule[1.1pt]

\multirow{5}{*}{\rotatebox{90}{\small Baselines}}
& DDP (baseline) & 53.2 & 24.1 & 40.8 & 70.5 & 53.4 & 48.4 \\
& \graycell{Local Adam} & \graycell{52.8} & \graycell{23.7} & \graycell{40.3} & \graycell{70.1} & \graycell{53.1} & \graycell{48.0} \\
& FAVG+OPT & 50.1 & 22.3 & 38.6 & 68.9 & 51.8 & 46.3 \\
& \graycell{DES-LOC} & \graycell{52.6} & \graycell{23.9} & \graycell{40.1} & \graycell{70.0} & \graycell{53.0} & \graycell{47.9} \\
& \faustcell{\textbf{FAUST}} & \faustcell{\textbf{52.9}} & \faustcell{\textbf{24.0}} & \faustcell{\textbf{40.5}} & \faustcell{\textbf{70.3}} & \faustcell{\textbf{53.2}} & \faustcell{\textbf{48.2}} \\

\midrule[1.1pt]

\multirow{4}{*}{\rotatebox{90}{\small Ablation}}
& \quad $K_u{=}K_x,\;K_v{=}K_x$ & 52.4 & 23.5 & 39.9 & 69.8 & 52.7 & 47.7 \\
& \graycell{\quad $K_u{=}6K_x,\;K_v{=}3K_x$} & \graycell{51.7} & \graycell{23.1} & \graycell{39.4} & \graycell{69.5} & \graycell{52.3} & \graycell{47.2} \\
& \quad $P_{\text{SP}}{=}1$ (no parallel) & 52.5 & 23.8 & 40.0 & 69.9 & 52.9 & 47.8 \\
& \graycell{\quad $P_{\text{SP}}{=}4$} & \graycell{52.7} & \graycell{23.9} & \graycell{40.3} & \graycell{70.1} & \graycell{53.0} & \graycell{48.0} \\

\bottomrule[1.6pt]
\end{tabular}%
}
\end{table}

% \paragraph{Takeaway:} FAUST enables efficient training of large-scale foundation models, especially at long
% training horizons, with downstream ICL performance competitive with DDP.
% FAUST with compile-time SP further closes the gap to DDP.
% The ablation confirms that the tiered schedule ($K_u{=}3K_x, K_v{=}6K_x$) outperforms both uniform ($K_u{=}K_v{=}K_x$) and inverted schedules,
% consistent with the half-life analysis in Section~\ref{sec:timescales}.
% Varying SP degree has minimal impact on ICL metrics, confirming that the compile-time sequence sharding
% does not distort the learned representations.
% Overall, the wall-clock speedup ranges from $1.5$--$2.4{\times}$ over DDP depending on model scale and $K_x$ (Table~\ref{tab:wallclock}).
% We recommend setting $K_x$ for throughput-bandwidth trade-off, then setting $K_u, K_v$ as constant multiples (e.g.,
% $3\times$, $6\times$) or scale to the half-life of their $\beta$ (see Section~\ref{sec:choosing_freq}).

%-----------------------------------------------------------------------------
% \subsection{Nesterov as the Outer Optimizer (RQ5)}
% \label{sec:rq5}
%-----------------------------------------------------------------------------

% We ablate the outer optimizer on a 700M parameter model, comparing averaging to a
% Nesterov optimizer with momentum of 0.9, outer learning rate of 1.0 tuned following~\citet{zhuansun2024communicationefficientfederatedlearningadaptive}. The experiment uses 4x node with a medium-frequency regime ($K_x = 32$, $K_u =
% 3K_x$, $K_v = 6K_x$) where models are initialized from 2048-step DDP checkpoints, following~\citet{zhuansun2024communicationefficientfederatedlearningadaptive}. While the convergence bound of the abstract inequality, the structure of Eq.~(\ref{eq:psi_def}) suggests a higher synchronization frequency ($p_u$) should permit a larger step size ($\eta_0 \propto 1/\sqrt{\psi}$).

% As shown in Figure~\ref{fig:nesterov}, two key trends emerge. First, with moderate synchronization ($K_x = 32$)
% pushes to with within $1\%$ of the final perplexity of DDP, consistent over the earlier year
% in significant increase ($K_x = 256$). Second, using Nesterov as the outer optimizer improved
% perplexity over averaging by $\approx 0.5\%$, with the performance i.e.e and scale quality at the the
% baseline at~\citet{zhuansun2024communicationefficientfederatedlearningadaptive}, Table~4) for models at this scale. The Nesterov schedule
% preserves all practical benefits, retaining effective worker initialization and reducing
% local optimisation noise, which are more coupled tightly than proximity protection steps.



%-----------------------------------------------------------------------------
% \subsection{Muon as the Inner Optimizer (RQ6)}
% \label{sec:rq6}
% %-----------------------------------------------------------------------------

% To assess the generality of DES-LOC beyond Adam-family optimizers, we integrate it with Muon~\citep{jordan2024muon}, a single-momentum optimizer that uses Newton--Schulz iterations for orthogonalisation.
% Since Muon preconditions the gradient from momentum rather than tracking second-moment
% estimates, the momentum synchronization reduces to parameters ($K_x$) and momentum ($K_u$). We
% apply DES-LOC by independently set the required period from the particular over the benchmarks,
% setting $K_u = 3K_x$ with $K_x = 32$.

% As shown in Figure~\ref{fig:muon}, DES-LOC matches the trajectory of the Local Muon baseline across
% model scales ranging from 16M to $360$M parameters. By decoupling the synchronization frequencies,
% DES-LOC communicates more than $1.5\times$ fewer bytes than the baseline, confirming that the half-life
% principle extends to optimizers beyond the Adam family. Full experimental details in Section~\ref{sec:muon_details}.

% % % Figure 7: RQ6 rate 7 figure
% % \begin{fi'guu're}[ht]
% % \centering
% % % [Figure 7 graphic from original PDF]
% % \caption{Training loss comparison between Local Muon ($K = 32$) and DES-LOC-Muon ($K_x = 32, K_u = 96$) across model scales ($16$M, $125$M, $360$M). DES-LOC matches the Local Muon baseline across all scales while communicating more than $\mathbf{1.5\times}$ fewer bytes. See Section~\ref{sec:muon_details} for full experimental details.}
% % \label{fig:muon}
% % \end{figure}

% \paragraph{Takeaway:} DES-LOC is compatible with single-momentum optimizers that use Newton--Schulz
% preconditioning. By preserving the synchronization cadence of the momentum buffer relative to
% parameters, DES-LOC maintains solution quality while reducing communication volume,
% demonstrating that the desynced paradigm is inner-optimizer-agnostic.

\paragraph{Compile-time SP ablation.} \begin{wraptable}{r}{0.54\textwidth}
\vspace{-6pt}
\centering
\scriptsize
\setlength{\tabcolsep}{4.5pt}
\renewcommand{\arraystretch}{1.12}
\caption{Compile-time SP ablation on Llama-3.1 using H100 GPUs. Left: maximum trainable sequence length before OOM. Right: per-iteration time for the 8B model.}
\label{tab:sp_ablation}
\resizebox{0.54\textwidth}{!}{%
\begin{tabular}{lcccccc}
\toprule[1.6pt]
\textbf{Method} &
\makecell{\textbf{3B}} &
\makecell{\textbf{8B}} &
\makecell{\textbf{13B}} &
\makecell{\textbf{8K}} &
\makecell{\textbf{24K}} &
\makecell{\textbf{90K}} \\
\midrule[1.1pt]
ZeRO-3 & 15K & 8K & 6K & 10 & OOM & OOM \\
\rowcolor{gray!5}
RingAttention & 35K & 15K & 9K & 46 & 29 & OOM \\
DS-Ulysses & 35K & 15K & 8K & 9 & 26 & OOM \\
\rowcolor[HTML]{F0F7FF}
\textbf{FAUST} & \textbf{65K} & \textbf{25K} & \textbf{10K} & \textbf{8} & 27 & \textbf{188} \\
\bottomrule[1.6pt]
\end{tabular}%
}
\vspace{-8pt}
\end{wraptable}
To isolate the contribution of the compile-time sequence-parallel pass,we ablate the SP backend
on Llama-3.1 8B using H100 GPUs.
Table~\ref{tab:sp_ablation} reports maximum trainable sequence length (before OOM) and
end-to-end per-iteration wall-clock time.
FAUST, the graph-rewriting pass that inserts head-dimension All-to-All collectives around
each attention layer (Section~\ref{sec:neuronsp})---extends trainable sequence length by
$2.14$--$5{\times}$ over ZeRO-3 across model scales.
Critically, the per-iteration overhead remains comparable to or below that of DS-Ulysses and
RingAttention: at $8$K sequence length on the 8B model, we take $8$\,s versus $10$\,s (ZeRO-3)
and $9$\,s (RingAttention/DS-Ulysses), while at $24$K it completes in $27$\,s where ZeRO-3 encounters OOM.
This confirms that the design
assumption that intra-step All-to-All can be decoupled from inter-step AllReduce
without contention.



%=============================================================================



% \section{Related Work}
% \label{sec:related}
% \label{sec:extended_related_work}
% %=============================================================================

% In synchronous data-parallel training, workers exchange full gradients or parameters \textit{every} iteration,
% incurring communication costs linear in model size using Ring-AllReduce~\citep{sergeev2018horovodfasteasydistributed}.
% When hardware is weakly connected or widely distributed, communication significantly
% slows wall-clock training time~\citep{sani2024photonfederatedllmpretraining} as workers need to wait for synchronization to
% finish. Federated Averaging (FedAvg)~\citep{mcmahan2017communication} and Local SGD~\citep{stich2019local}
% reduce communication by performing $K$ local optimization steps before averaging parameters,
% decreasing communication rounds by a factor of $K$. Ad-hoc extensions to adaptive optimizers either
% keep optimizer states local~\citep{douillard2024dilocodistributedlowcommunicationtraining,hilmkil2021scalingfederatedlearningfinetuning,liu2024asynchronouslocalsgdtraininglanguage} or reset them
% after each sync~\citep{sani2024futurelargelanguagemodel,sani2024photonfederatedllmpretraining}, both lacking robust convergence guarantees.

% Adam~\citep{kingma2017adammethodstochasticoptimization} is popular for pre-training as it scales to larger batches than SGD~\citep{kunstner2023noisemainfactorgap,grattafiori2024llama3herdmodels}. It uses moving averages of gradients and their squares, however,
% its convergence is not guaranteed as it requires $\beta_1 < \sqrt{\beta_2} < 1$, with large, problem-specific $\beta_2$~\citep{reddi2019convergenceadam,zhang2023adamconvergemodificationupdate}. Other optimizers also track gradient moments~\citep{pmlr-v28-sutskever13,chen2023symbolicdiscoveryoptimizationalgorithms,you2020largebatchoptimizationdeep,taniguchi2024adoptmodifiedadamconverge}. \textit{Local Adam}~\citep{cheng2025convergencedistributedadaptiveoptimization} reduces communication with local steps but requires syncing optimizer states, which triples the
% communication cost relative to Local SGD/DDP, as sync costs scale with the number of states. For
% further related work, including compression/sparsification and structured updates, check Section~\ref{sec:extended_related_work}.

% %=============================================================================
% \section{Conclusion}
% \label{sec:conclusion}
% %=============================================================================
% In this paper, we present FAUST, the first compiler-orchestrated, PyTorch-native solution for training
% foundation models at scale with minimal inter-node communication. Through a composition of compile-time sequence-parallel
% graph rewriting and a half-life-proportional tiered synchronization schedule, FAUST decouples intra-step activation redistribution from inter-step gradient
% synchronization at negligible cost to training throughput. Our theoretical
% analysis establishes that FAUST converges via a multiplicative potential contraction, where the
% per-step rate is governed by the least-active synchronization tier, and that the half-life-proportional
% schedule $K_u{=}3K_x$, $K_v{=}6K_x$ equalizes per-tier contraction rates,providing an
% information theoretic justification for the tiered design. Our results demonstrate that compiler-driven, PyTorch-native
% automation composed with multi-rate desynced optimization provides a practical and portable foundation for communication-efficient billion-scale model training.
% Experiments on models up to 1.3B parameters show that FAUST achieves $303{\times}$ cumulative communication reduction over DDP and $1.78{\times}$ over DES-LOC alone, while reaching perplexity within $4.6\%$ of the fully-synchronized baseline in $27\%$ less wall-clock time.
% On bandwidth-constrained links ($\leq 10$\,Gbps), the speedup reaches $48$--$305{\times}$ over DDP.





\section{Related Work}
\label{sec:related}
\label{sec:extended_related_work}
%=============================================================================

Synchronous data parallelism exchanges gradients every step via
Ring-AllReduce~\citep{sergeev2018horovodfasteasydistributed,li2020pytorch},
incurring communication linear in model size.
FedAvg~\citep{mcmahan2017communication} and Local
SGD~\citep{stich2019local,yu2019linearspeedupanalysiscommunication} amortize
this cost by performing $K$ local steps before averaging, with
SCAFFOLD~\citep{karimireddy2021scaffoldstochasticcontrolledaveraging} correcting
client drift. Separately, Adam~\citep{kingma2017adammethodstochasticoptimization}
became the dominant optimizer for LLM
pre-training~\citep{brown2020languagemodelsfewshotlearners,grattafiori2024llama3herdmodels},
scaling to large batches where SGD
stalls~\citep{kunstner2023noisemainfactorgap}, though its convergence requires
$\beta_1 < \sqrt{\beta_2}$~\citep{reddi2019convergenceadam,zhang2023adamconvergemodificationupdate}.

On the optimizer side, ADOPT~\citep{taniguchi2024adoptmodifiedadamconverge}
removes Adam's $\beta$-ordering constraint,
Muon~\citep{jordan2024muon,liu2025muonscalablellmtraining} replaces second moments with
Newton--Schulz orthogonalisation, and adaptive federated
methods~\citep{reddi2021adaptivefederatedoptimization,zhuansun2024communicationefficientfederatedlearningadaptive}
extend outer optimizers to the cross-replica setting.
On the parallelism side, DeepSpeed
Ulysses~\citep{jacobs2023deepspeedulyssesoptimizationsenabling} inserts
head-splitting All-to-All collectives around attention, and
USP~\citep{fang2024uspunifiedsequenceparallelism} unifies ring and Ulysses
strategies. Distributing Adam across replicas remains difficult: ad-hoc
extensions either keep optimizer states entirely
local~\citep{douillard2024dilocodistributedlowcommunicationtraining,hilmkil2021scalingfederatedlearningfinetuning,liu2024asynchronouslocalsgdtraininglanguage},
accumulating drift without convergence guarantees, or reset them after each
sync~\citep{sani2024futurelargelanguagemodel,sani2024photonfederatedllmpretraining}.
Local Adam~\citep{cheng2025convergencedistributedadaptiveoptimization} provides
the first convergent intermittent-communication adaptive analysis but synchronises
all states at a uniform period, tripling the payload versus Local SGD.

\paragraph{Positioning of FAUST.}
No prior work jointly addresses (i)~provable convergence under decoupled
per-state synchronisation cadences and (ii)~compiler-level composition with
sequence parallelism. FAUST fills this gap by assigning half-life-proportional
averaging periods to parameters and each optimizer state, halving Local Adam's
communication while preserving convergence, and lifting SP into the PyTorch~2
compiler stack so that intra-step All-to-All and inter-step AllReduce compose on
orthogonal axes without manual annotation.

%=============================================================================


%=============================================================================
% \section{Conclusion}
% \label{sec:conclusion}
% %=============================================================================
% We presented FAUST, composing compile-time sequence parallelism with runtime
% desynced synchronization that assigns independent averaging periods to parameters
% and optimizer states by EMA half-life. Experiments confirm FAUST matches communication-free baselines while providing convergence
% guarantees they lack.

\section{Conclusion}
\label{sec:conclusion}
We presented FAUST, a framework that composes compile-time sequence parallelism
with runtime desynced synchronization for communication-efficient distributed
LLM training. A PyTorch~2 FX graph rewrite inserts intra-step All-to-All
collectives for SP, while a half-life-proportional schedule assigns independent
averaging periods $K_x$, $K_u{=}3K_x$, $K_v{=}6K_x$ to parameters and optimizer
states, gating inter-step AllReduce by tier. A collective fence ensures the two
collective classes compose on orthogonal axes without interference.
Our convergence analysis shows the per-step contraction rate is governed by the
least-active tier, and the half-life schedule equalises all per-tier rates any
other allocation creates a bottleneck that slows convergence.
FAUST achieves $303{\times}$ communication reduction over DDP and $1.78{\times}$
speedup over DES-LOC, reaching perplexity within $4.6\%$ of the synchronous
baseline in $27\%$ less wall-clock time, with $48$--$303{\times}$ speedup on
bandwidth-constrained links. These results show that compiler-driven SP composed
with multi-rate desynced optimisation provides a practical foundation for
billion-scale training.



\section{Limitation}
\label{sec:limitation}
FAUST is robust to worker failures, extending the compile-time SP pass to handle SP-group recovery via the existing AllReduce infrastructure without requiring checkpointing.
We note several directions for future work: adaptive per-layer synchronization frequencies that exploit the structural-tier curvature hierarchy (Assumption~\ref{asm:smooth}), integration with gradient compression to further reduce the collective payload, and extension of the convergence theory to non-smooth objectives and stochastic synchronization schedules beyond the Bernoulli model.


\bibliographystyle{plainnat}
\bibliography{references}




\appendix

%=============================================================================


%=============================================================================
\section{Experimental Details and Optimizer Hyperparameter Sweeps}
\label{sec:appendix_a}
%=============================================================================

Here we provide additional experimental details, including: a)~model architecture details, sequence-parallelism configuration, and training hyperparameters independent of optimizer choice (Section~\ref{sec:arch_details}), b)~our hyperparameter sweep procedure to select optimizer-specific settings under the compile-time SP decomposition (Section~\ref{sec:sweep}), and c)~the optimal hyperparameters with those used in Section~\ref{sec:evaluation} highlighted in bold.

\subsection{Architecture Details and Hyperparameters}
\label{sec:arch_details}

\begin{table}[ht]
\centering
\caption{Model architecture, sequence-parallelism configuration, and training parameters.
Here, \#Blocks, \#Heads, $d_{\text{model}}$, $|V|$, and Exp.\ Ratio denote the number of transformer blocks, attention heads, embedding dimension, vocabulary size, and feedforward expansion ratio, respectively.
All models use RoPE~\citep{su2023roformerenhancedtransformerrotary}, SiLU activation, norm-based gradient clipping with bound $\rho$, and initialization scale $\sigma = 1/\sqrt{d_{\text{model}}}$.
$|B_G|$ is the global batch size, $T$ is the number of optimization steps, and $P_{\text{SP}}$ is the compile-time sequence-parallel degree.
All experiments use Ulysses head splitting, with $P_{\text{SP}}$ constrained to evenly divide \#Heads.}
\label{tab:arch}
\resizebox{\linewidth}{!}{%
\begin{tabular}{lccccccccccccc}
\toprule
Model & Blocks & $d_{\text{model}}$ & $|V|$ & Heads & Exp. & RoPE & Act. & Init. & $\rho$ & Seq. & $|B_G|$ & $P_{\text{SP}}$ & $T$ \\
Size &        &                    &       &       & Ratio & $\theta$ &      & $\sigma$ &        & Len. &          &                 &     \\
\midrule
135M & 30 & 576  & 50K & 9  & 4 & 10000 & SiLU & 0.04 & 1.0 & 2048 & 1024 & 1 & 1536, 3072 \\
720M & 12 & 2048 & 50K & 16 & 4 & 10000 & SiLU & 0.02 & 1.0 & 2048 & 512  & 2 & 38912 \\
1.3B & 24 & 2048 & 50K & 16 & 4 & 10000 & SiLU & 0.02 & 1.0 & 2048 & 1024 & 2 & 20480 \\
\bottomrule
\end{tabular}%
}
\end{table}

Table~\ref{tab:arch} summarizes the architectural details of our models, following established practices for large
language models at their respective scales. Unless otherwise stated, we adopt the hyperparameters
recommended by~\citet{allal2025smollm2} for both the 135M and the 1.3B models. We operate at a
batch size of 2M tokens, which is very large for the 135M model at the length of training we
perform~\citep{merrill2025criticalbatchsizerevisited} and industry-standard for the 1.3B model~\citep{touvron2023llama2openfoundation};
we chose to operate at large batch sizes because adaptive optimizers provide benefits primarily in
large-batch training regimes~\citep{kunstner2023noisemainfactorgap}.
Moreover, we intend FAUST for deployment in heterogeneous accelerator clusters and
cross-data-center scenarios, where effectively utilizing available devices naturally demands large
batch sizes and/or model scales.
For both model sizes, we train for approximately $2\times$ the compute-optimal token
budget~\citep{hoffmann2022trainingcomputeoptimallargelanguage}, placing our evaluations within the context of extended-duration
foundation model training~\citep{allal2025smollm2}. Our chosen token budget is conservative due to
resource constraints; for comparison, \citet{allal2025smollm2} used 11 trillion tokens which is
over $4000\times$ compute-optimal for the 135M model, and $300\times$ for the 1.3B.

\paragraph{Sequence-parallelism configuration.}
\label{sec:neuronsp}
The compile-time SP pass (Section~\ref{sec:neuronsp}) partitions the sequence dimension across
$P_{\text{SP}}$ devices by rewriting the traced computation graph. For each scaled dot-product
attention layer, the pass inserts a head-dimension All-to-All collective before the attention
kernel and a conjugate All-to-All after the output projection, following the Ulysses
head-splitting strategy~\citep{jacobs2023deepspeedulyssesoptimizationsenabling,fang2024uspunifiedsequenceparallelism}.
The SP degree $P_{\text{SP}}$ is constrained to divide the number of attention heads evenly;
for the 135M model with 9 heads we set $P_{\text{SP}} = 1$ (SP disabled) to isolate the
effect of the desynced communication schedule, while for the 720M and 1.3B models
($\text{\#Heads} = 16$) we set $P_{\text{SP}} = 2$.
The SP group is formed \emph{within} a single node over NVLink, so the head-dimension
All-to-All executes at intra-node bandwidth (${\sim}600$\,GB/s per H100 NVLink domain)
and does not compete with the inter-node data-parallel AllReduce that operates over
InfiniBand at $K_x$-boundary steps.
A collective fence mechanism (Section~\ref{sec:neuronsp}) serializes all pending SP handles
before the cross-replica gradient reduction begins, preventing communicator contention on
clusters with shared interconnect resources.

We select warmup and decay schedules following recommendations from~\citet{merrill2025criticalbatchsizerevisited,hagele2024scaling,allal2025smollm2}. For the 135M model, the warmup period is set to $T_{\text{WARM}} = 512$
steps, corresponding to roughly 40\% of the compute-optimal training tokens recommended by~\citet{merrill2025criticalbatchsizerevisited}. For the 1.3B model, we use the recommended $T_{\text{WARM}} = 2048$ steps from~\citet{allal2025smollm2}, roughly 10\% of total training. The stable-decay period uses a $1 - \text{SQRT}$ schedule over
the final $T_{\text{DECAY}} = 10\% \times T$ steps~\citep{hagele2024scaling}. For shorter runs, such as $T = 1536$ during heterogeneous-data evaluations, we keep the warmup fixed and proportionally scale the decay to ensure well-conditioned parameter updates during the stable learning rate period. The seeds we use for data sampling and for controlling the training algorithms and model are provided in the code accompanying the appendix.

\paragraph{$K_x$-aligned WSD boundaries.}
Under FAUST's desynced communication schedule, the warmup-stable-decay phase boundaries
must be snapped to multiples of the synchronization period $K_x$ to avoid inconsistent
learning rates across workers during a local-step window.
Concretely, we set $T'_{\text{WARM}} = \lceil T_{\text{WARM}} / K_x \rceil \cdot K_x$ and
$T'_{\text{DECAY}} = \lceil T_{\text{DECAY}} / K_x \rceil \cdot K_x$, so that every
cross-replica AllReduce averages gradients computed under a single learning-rate phase.
Without this alignment, workers on opposite sides of a phase boundary would contribute
gradients scaled by different $\eta$ values to the same reduction, biasing the averaged
pseudo-gradient.
This alignment is exact for the FedAvg outer optimizer; when a Nesterov outer optimizer is
used, the alignment ensures that the outer momentum accumulation
does not conflate stable-phase and decay-phase inner gradients.

\subsection{Optimizer Parameters Sweeping Procedure}
\label{sec:sweep}

As detailed in Section~\ref{sec:neuronsp} and verified empirically, the choice of decay rates $\beta_1, \beta_2$
strongly influences the effective synchronization frequencies achievable by both FAUST and
Local Adam. This relationship arises directly from the half-life of optimizer states, given by
$\tau_{0.5} = \frac{\ln(0.5)}{\ln(\beta)}$.
Under FAUST's compile-time SP decomposition, the per-device gradient is already
sequence-reduced (via the backward All-to-All) before it enters the optimizer, so the
$\beta$-dependent half-life analysis applies identically to SP and non-SP configurations;
the SP degree $P_{\text{SP}}$ affects only the effective per-device batch size
($B_w / P_{\text{SP}}$ tokens per sequence chunk) but does not alter the optimizer-state
dynamics.

For Adam, prior studies such as \citet{wortsman2023stablelowprecisiontraininglargescale} have demonstrated a critical interplay between the learning rate ($\eta$), batch size, and the second-momentum decay $\beta_2$. Specifically, increasing either the learning rate or batch size typically demands a lower $\beta_2$ to maintain training stability and avoid loss spikes. Conversely, higher $\beta_2$ values constrain the learning rate and batch size. Such dynamics have also been recently observed between the learning rate and the first-momentum decay $\beta_1$ in~\citet{pagliardini2025the}. Given that all our experiments use a fixed large batch size of roughly 2 million tokens (appropriate for billion-scale training), we systematically tune the learning rate $\eta$ in response to changes in $\beta_1, \beta_2$. We try values of $\beta_1, \beta_2$ based on previous works~\citep{merrill2025criticalbatchsizerevisited} and follow the theoretical convergence requirement of~\citet{zhang2023adamconvergemodificationupdate} setting $\beta_1 \leq \sqrt{\beta_2}$.

Due to computational constraints, we cannot jointly optimize synchronization periods, sequence-parallel degree, data distributions, and decay parameters, and instead adopt a structured three-stage tuning approach:
\begin{enumerate}
    \item \textbf{Stage 1: Tuning $\eta$ for SP-augmented DDP.} Starting from the recommended baseline learning rate ($\eta_0$)
    from~\citet{allal2025smollm2}, we conduct a grid search as outlined by~\citet{zhuansun2024communicationefficientfederatedlearningadaptive}:
    $\{\ldots, \sqrt{2}^{-2}\eta_0, \sqrt{2}^{-1}\eta_0, \eta_0, \sqrt{2}\eta_0, \sqrt{2}^2\eta_0, \ldots\}$
    We expand this search until perplexity stops improving, identifying an optimal learning rate $\eta^*_{\text{DDP}}$ for each $(\beta_1, \beta_2)$ configuration.
    When $P_{\text{SP}} > 1$, the SP All-to-All collectives execute every step and their cost is amortized into the per-step baseline, so Stage~1 already accounts for the SP overhead.

    \item \textbf{Stage 2: Tuning $\eta$ for Local Adam.} We then repeat this procedure for Local Adam,
    using $\eta^*_{\text{DDP}}$ as the new baseline. To balance generalizability and computational cost, we set
    the synchronization period to an intermediate value of $K = 64$, between high-frequency ($K = 16$) and low-frequency ($K = 256$) scenarios.

    \item \textbf{Stage 3: Validating $\eta$ under FAUST's tiered schedule.}
    Using the optimal $\eta^*_{\text{Local}}$ from Stage~2, we verify that the desynced
    schedule ($K_x = K$, $K_u = 3K_x$, $K_v = 6K_x$) does not require further adjustment.
    In all configurations tested, the Stage~2 learning rate transferred without modification
    to the tiered schedule, confirming that decoupling momentum synchronization periods does
    not alter the effective curvature landscape seen by the optimizer.
    This is consistent with the theoretical result that synchronization periods enter only
    the higher-order $\mathcal{O}(1/T)$ term of the convergence rate
    (Theorem~\ref{thm:sgdm}).
\end{enumerate}

Additionally, following~\citet{merrill2025criticalbatchsizerevisited}, we omit weight decay (set to zero) to simplify the hyperparameter tuning process, as it directly affects only model parameters, not optimizer states.

For experiments using Nesterov, we follow the hyperparameter sweeping procedure of~\citet{zhuansun2024communicationefficientfederatedlearningadaptive}, starting with a server learning rate of $1.0$ and a momentum of $0.9$ and only lowering it if it fails to converge.

\subsubsection{Optimizers' Hyperparameter Configurations}
\label{sec:hp_configs}

\begin{table}[ht]
\centering
\caption{Optimal learning rates $\eta^*$ for $\beta_1, \beta_2$ configurations of ADOPT/Adam. The hyperparameter sweep procedure (see Section~\ref{sec:sweep}) involves incrementally adjusting the learning rate by factors of $\sqrt{2}$ around the initial value from~\citet{allal2025smollm2} until performance stops improving.}
\label{tab:lr_sweep}
\begin{tabular}{llll}
\toprule
Optimizer & $\beta_1$ & $\beta_2$ & $\eta^*$ \\
\midrule
\multirow{4}{*}{ADOPT} & 0.9   & 0.9999 & 0.0021 \\
                        & \textbf{0.95}  & \textbf{0.9999} & \textbf{0.0021} \\
                        & 0.99  & 0.9999 & 0.0014 \\
                        & 0.995 & 0.9999 & 0.0007 \\
\midrule
\multirow{5}{*}{Adam}  & 0.9   & 0.95   & 0.0042 \\
                        & \textbf{0.95}  & \textbf{0.95}   & \textbf{0.003}  \\
                        & 0.9   & 0.99   & 0.003  \\
                        & 0.95  & 0.99   & 0.003  \\
                        & 0.99  & 0.99   & 0.0021 \\
\bottomrule
\end{tabular}
\end{table}

Our hyperparameter sweep (Table~\ref{tab:lr_sweep}) indicates that the optimal learning rate $\eta^*$ under the $K_x$-aligned warmup-stable-decay scheduler~\citep{hagele2024scaling} strongly depends on both optimizer type and the chosen $\beta_1, \beta_2$ values. For Adam, optimal learning rates and second-momentum decay ($\beta_2$) align closely with recommendations from~\citet{allal2025smollm2}, though a slightly higher first-momentum decay ($\beta_1$) consistently performs better, in agreement with prior findings~\citep{merrill2025criticalbatchsizerevisited}. For ADOPT (default $\beta_2$), we observe a lower optimal learning rate compared to Adam, but similar best-performing $\beta_1$ values. We also find that the optimal learning rate does not differ between SP-augmented DDP and Local Adam for given $\beta_1, \beta_2$ when $K = 64$ and using a $\sqrt{2}$ sweep; higher learning rates either do not provide a benefit or diverge, while lower learning rates are only necessary when pushing $K$ far closer to the complete training duration. Crucially, the SP degree $P_{\text{SP}}$ does not interact with $\eta^*$: the compile-time sequence sharding preserves the per-device effective batch size and gradient scale, so the same $\eta^*$ applies with $P_{\text{SP}} = 1$ and $P_{\text{SP}} = 2$.

We find that increasing $\beta_1$ for ADOPT, and $\beta_1, \beta_2$ for Adam, leads to rapid performance degradation, particularly at or above $0.99$. Since the half-life at $\beta = 0.99$ ($\tau_{0.5} \approx 69$) is not sufficiently longer than at $\beta = 0.95$ ($\tau_{0.5} \approx 13.5$) to justify the observed performance drop, we select $\beta_1 = 0.95$ for all experiments, along with the default $\beta_2$ for ADOPT and $\beta_2 = 0.95$ for Adam.

\paragraph{Takeaway.} Increasing an optimizer state's $\beta$ significantly affects performance. Since linear increases in $\beta$ cause only logarithmic changes in half-life $\tau_{0.5}$, raising $\beta$ beyond the optimal value degrades performance without substantially improving the achievable synchronization frequency. Under FAUST's tiered schedule, this implies that the communication savings from increasing $K_v$ (via higher $\beta_2$) are bounded: the half-life ratio $\tau_{0.5}(\beta_2) / \tau_{0.5}(\beta_1)$ determines the maximum useful desync factor, beyond which further increasing $\beta_2$ degrades convergence faster than it reduces collective volume.


%=============================================================================
\
%=============================================================================
\section{Complementary Results to Sections 2.1 and 5}
\label{sec:complementary}
\label{sec:critical_batch}
%=============================================================================

We now provide additional results supplementing those presented in the main text. Specifically:

\begin{enumerate}
\item \textbf{Section B.2.1} complements Fig.~3 by showing the separate impact of varying synchronization frequencies for parameters and the second momentum when the base frequency is $K_b=16$. It supports our claim that parameters and second momentum exhibit similar behavior across different synchronization regimes, unlike the first momentum.
\item \textbf{Section B.2.2} extends Fig.~3 by evaluating FAUST-Adam. We confirm that the parameter synchronization frequency is the most important, as predicted by our theory. In contrast, the momenta sync frequency is far less impactful, especially for low parameter sync frequencies. The compile-time SP decomposition does not alter these dynamics, as the sequence-reduced gradients fed to the optimizer preserve the same $\beta$-dependent half-life hierarchy.
\item \textbf{Section B.3.3} presents an ablation study examining alternative low-communication configurations of FAUST, justifying our choice of $K_u=3K_x, K_v=6K_x$ used in Fig.~4.
\item \textbf{Section B.3.4} repeats the baseline comparison from Fig.~4 for FAUST-Adam, demonstrating that FAUST achieves similar communication reductions and performance when using Adam instead of ADOPT.
\item \textbf{Section B.4} provides additional metrics illustrating training instabilities for the \texttt{FAVG+OPT} baseline, including rapidly growing parameter norms, supporting observations in Fig.~5.b.
\item \textbf{Section B.5} presents the full 1B-scale comparison of FAUST Nesterov against DiLoCo, demonstrating the benefits of periodic optimizer-state synchronization under a Nesterov outer optimizer.
\item \textbf{Section B.6} benchmarks FAUST under extremely low bandwidth ($10$ Gbit/s), showing $9.42\times$ wall-clock speedups over DDP.
\end{enumerate}

%-----------------------------------------------------------------------------
\subsection{RQ2: Independent Sync Frequencies}
\label{sec:b2_sync_freq}
%-----------------------------------------------------------------------------

This section provides supplementary results. Section~B.2.1 shows that perplexity has similar sensitivity to the first and second momentum synchronization frequencies at both high and low base synchronization frequencies. Additionally, Section~B.2.2 repeats the comparison from Fig.~3 for FAUST-Adam, revealing similar trends regarding the importance of the parameters, with a reduced importance for the momenta due to lower $\beta_2$.

\subsubsection{Parameter and Second Momentum at \texorpdfstring{$K_b=16$}{Kb=16} (See Fig.~3.a and Fig.~3.b)}
\label{sec:b21_param_second}

Figure~9 examines the effects of independently varying synchronization periods ($K_x, K_v$) for parameters and second momentum under FAUST-ADOPT in the high-frequency regime ($K_b=16$), chosen based on the first momentum's half-life ($\tau_{0.5}\approx 13.5$). Similar to the low-frequency results in Fig.~3.a, parameter synchronization frequency ($K_x$) strongly influences perplexity, while the second momentum ($K_v$) has minimal impact due to its long half-life. This contrasts with the first momentum, whose half-life closely matches the high-frequency period. The SP degree $P_{\text{SP}}$ does not affect these trends: the compile-time sequence sharding operates on the activation dimension and is transparent to the optimizer-state synchronization schedule.

%-----------------------------------------------------------------------------
\subsection{RQ3: Communication Reduction And Baseline Comparisons}
\label{sec:b3_comm_reduction}
%-----------------------------------------------------------------------------

This section provides supplementary results for \textbf{RQ3}, complementing Section~\ref{sec:rq3}. Section~B.3.2 shows the perplexity of different configurations providing a $2\times$ communication reduction over Local Adam. Additionally, Section~B.3.4 repeats the comparison against baselines from Section~\ref{sec:rq3} for FAUST-Adam, showing similar communication reductions relative to Local Adam. In all experiments the compile-time SP pass remains fixed at the degree specified in Table~\ref{tab:arch}; the SP All-to-All traffic is identical across all methods and therefore cancels when computing relative communication savings.

\subsubsection{Stepwise plots for Baseline Comparison}
\label{sec:b31_stepwise}


% Figure 12 and 13: Stepwise plots (PDF page 23)
% [placeholder for Figs 12-13 -- stepwise counterparts to Figs 4,5]

\subsubsection{FAUST Low Communication Configurations Ablation (See Fig.~4)}
\label{sec:b33_ablation}


% Figure 15: Neuron_SP 2x communication configs ablation (PDF page 24 bottom)


\subsubsection{Adam Results (See Fig.~4)}
\label{sec:b34_adam_fig4}

We now present results for FAUST-Adam with $\beta_1=\beta_2=0.95$. FAUST-Adam achieves similar communication reductions over Local Adam and DDP as ADOPT. However, due to the lower $\beta_2$, the second-momentum half-life ($\tau_{0.5}(0.95)\approx 13.5$) is significantly shorter than for ADOPT ($\tau_{0.5}(0.9999)\approx 6931$). Figure~16 shows that with both momenta evolving at similar rates, the benefit of selecting $K_u < K_v$ diminishes. For consistency and due to meaningful empirical differences in rates of change, we keep $K_u=3\times K_x$ and $K_v=6\times K_x$ in subsequent comparisons.




%-----------------------------------------------------------------------------
\subsection{RQ4: Additional Metrics and Training Instabilities of FAVG+OPT (See Fig.~5.b)}
\label{sec:b4_favgopt}
%-----------------------------------------------------------------------------


% Table 5: Throughput comparison (PDF page 26, nougat inline)
\begin{table}[t]
\centering
\small
\setlength{\tabcolsep}{7.5pt}
\renewcommand{\arraystretch}{1.12}
\caption{Throughput (batches/sec) for 1B--13B models to reach $2\times$ compute-optimal tokens at high ($K_x=16$) and low ($K_x=256$) synchronization frequencies with a 2M-token global batch.
All local methods improve throughput over DDP.
At high frequency, FAUST achieves over $1.7\times$ throughput on the 13B model; at low frequency, the gain grows to over $2.1\times$.
The SP All-to-All overhead is included in all measurements and adds $<2\%$ per-step latency at intra-node NVLink bandwidth.}
\label{tab:throughput}
\resizebox{\linewidth}{!}{%
\begin{tabular}{lcccccc}
\toprule[1.6pt]
\textbf{Method} &
\makecell{\textbf{1B}\\$K_x=16$} &
\makecell{\textbf{1B}\\$K_x=256$} &
\makecell{\textbf{7B}\\$K_x=16$} &
\makecell{\textbf{7B}\\$K_x=256$} &
\makecell{\textbf{13B}\\$K_x=16$} &
\makecell{\textbf{13B}\\$K_x=256$} \\
\midrule[1.1pt]
DDP (Baseline)
& $171.9^{\pm 0.94}$ & $171.9^{\pm 0.9}$
& $25.1^{\pm 0.10}$ & $25.1^{\pm 0.10}$
& $11.1^{\pm 0.03}$ & $11.11^{\pm 0.03}$ \\
\rowcolor{gray!5}
FAVG+OPT
& $304.9^{\pm 2.51}$ & $385.4^{\pm 3.7}$
& $34.0^{\pm 0.11}$ & $40.4^{\pm 0.15}$
& $19.4^{\pm 0.11}$ & $23.5^{\pm 0.15}$ \\
Local Adam
& $253.1^{\pm 1.59}$ & $380.1^{\pm 3.6}$
& $30.9^{\pm 0.09}$ & $40.2^{\pm 0.15}$
& $16.7^{\pm 0.07}$ & $23.3^{\pm 0.14}$ \\
\rowcolor[HTML]{F6FCF8}
\textbf{FAUST} ($K_u,\ K_v=3K_x,\ 6K_x$)
& $\mathbf{299.2}^{\pm 2.39}$ & $\mathbf{384.1}^{\pm 3.7}$
& $\mathbf{33.7}^{\pm 0.11}$ & $\mathbf{40.4}^{\pm 0.15}$
& $\mathbf{19.0}^{\pm 0.10}$ & $\mathbf{23.5}^{\pm 0.15}$ \\
\bottomrule[1.6pt]
\end{tabular}%
}
\end{table}

%=============================================================================
%=============================================================================
\subsection{Very Low Bandwidth Experiments}
\label{sec:b6_low_bandwidth}
%=============================================================================

While perplexity is invariant to network bandwidth, wall-clock time is highly sensitive to it. To practically showcase this, we perform a benchmark with a 1B model to measure time under extremely low bandwidth conditions (10 Gbit/s). This setup simulates a scenario with affordable, consumer-grade interconnects rather than data-centers---precisely the cross-data-center deployment target for which FAUST is designed (Section~\ref{sec:introduction}). Due to the extreme gradient synchronization delay inherent to DDP in this regime, the benchmark was limited to a $10{,}240$ step horizon to remain feasible. The SP All-to-All collectives still execute at intra-node NVLink bandwidth and are therefore unaffected by the inter-node bandwidth restriction; only the data-parallel AllReduce at $K_x$-boundary steps traverses the slow link.


FAUST Nesterov dramatically reduces training time by $\approx 9.42\times$ compared to DDP, completing the run in $8.99$ hours versus $84.73$ hours ($3.5$ days) for DDP, even with the constant overheads of our unoptimized implementation. Furthermore, FAUST Nesterov is more efficient than Local Adam, finishing $\approx 7\%$ faster (8.99h vs.\ 9.62h) while achieving significantly lower perplexity (8.91 vs.\ $10.19$). When compared to the ultra-lightweight DiLoCo baseline (8.68h), FAUST Nesterov incurs a time penalty of $\approx 3.6\%$ due to the additional optimizer state synchronization. However, this yields performance gains, improving final perplexity by $\approx 2.3\%$ (8.91 vs.\ 9.12) over DiLoCo.


%=============================================================================
\subsection{Muon Experimental Details }
\label{sec:b7_muon_details}
\label{sec:muon_details}
%=============================================================================

This section provides the full experimental details for the Muon experiments presented. Although a comprehensive theoretical treatment of preconditioned local updates under Newton-Schulz iterations is outside the scope of this work, the FAUST design is inherently compatible with such structures, as the momentum buffer follows standard EMA dynamics regardless of the subsequent preconditioning step. Moreover, the compile-time SP pass is optimizer-agnostic: it operates on the forward/backward computation graph and does not interact with the optimizer's internal state management, so Muon's Newton-Schulz orthogonalization composes with SP identically to Adam's moment updates.

\textbf{Experimental Details.} We utilize the standard PyTorch implementation of Muon with Nesterov momentum enabled and a weight decay of $0.1$. Following the recommendations of \citet{liu2025muonscalablellmtraining}, we apply the \texttt{match\_rms\_norm} adjustment to learning rates. We adopt the conventional split optimization strategy for Muon: AdamW handles embeddings and layer normalizations, while Muon optimizes all 2D matrices~\citep{jordan2024muon}. The momentum parameter for Muon is set to $\beta=0.9$, while the Adam component retains the $\beta_1=0.9, \beta_2=0.999$ settings used elsewhere. Gradient clipping thresholds are scaled by model size: $1.0$ for 16M, $0.5$ for 125M, and $0.25$ for 360M. For the Local Muon baseline, all optimizer states (Muon momentum; Adam first/second momenta) synchronize every $32$ steps. In contrast, FAUST delays state synchronization: the first momentum (for both optimizers) synchronizes every $96$ steps ($3\times$ reduction), and Adam's second momentum synchronizes every $192$ steps ($6\times$ reduction).

%=============================================================================
\subsection{Experiments on the Flux Vision Model}
\label{sec:b8_flux}
%=============================================================================

To demonstrate the universality of FAUST across different modalities and architectures beyond standard decoder-only LLMs, we evaluate its performance on Flux~\citep{sergeev2018horovodfasteasydistributed}, a Rectified Flow Transformer designed for text-to-image generation. This architecture differs significantly from the causal language models evaluated in previous sections, serving as a robust test for the generalizability of our decoupled synchronization approach. Because Flux uses bidirectional attention rather than causal masking, the compile-time SP pass applies the same Ulysses head-splitting strategy but without the causal mask partitioning required for autoregressive models, simplifying the graph rewrite.

\textbf{Experimental Setup.} We utilize the $280$M parameter variant of Flux provided by \texttt{torchtitan}, training with a global batch size of $256$. The inner optimizer is AdamW with $\beta_1=0.9, \beta_2=0.999$. We compare three settings: (1)~DDP, (2)~Local Adam with a synchronization period of $K=32$, and (3)~FAUST with a parameter sync period of $K_x=32$. For FAUST, we decouple the momentum synchronization significantly, setting $K_u=3K_x$ and $K_v=6K_x$ (192 steps).


%=============================================================================
\subsection{Throughput at 7B Scale}
\label{sec:b9_throughput}
%=============================================================================

To assess the practical scalability of our method on state-of-the-art hardware and at large model scales, we measure the training throughput of a 7B parameter model distributed across $8$ independent NVIDIA B200 GPUs.

\textbf{Throughput Analysis.} As illustrated during the local update phases, each GPU operates at the peak efficiency of a fully isolated local run, achieving identical tokens-per-second throughput as a single B200 with zero synchronization overhead. Distinct drops in throughput are observed only at the sparse synchronization boundaries ($K_x=32$, $K_u=96$, $K_v=192$), where the system pauses to aggregate model parameters and optimizer states via the tiered AllReduce. The SP All-to-All collectives execute every step but are overlapped with the attention computation on the NVLink fabric, contributing $<2\%$ overhead that is invisible in the throughput trace.


Crucially, standard DDP incurs this communication penalty at \textit{every single training step}, significantly lowering the average tokens/sec. Even with our current unoptimized ``stop-the-world'' implementation---which explicitly pauses computation to communicate and does not leverage computation-communication overlap---FAUST significantly increases aggregate throughput by amortizing these costs over long local training windows.

\begin{center}\fbox{\parbox{0.95\textwidth}{
\textbf{Takeaway:} On high-performance B200 hardware, FAUST enables near-linear scaling by keeping workers in a high-throughput local regime for the majority of training. By restricting communication overhead to sparse intervals, it delivers significant wall-clock speedups over DDP, even without low-level implementation optimizations like communication overlap. The compile-time SP pass adds negligible per-step cost ($<2\%$) since its All-to-All collectives operate entirely within the intra-node NVLink domain.
}}\end{center}

%=============================================================================
%=============================================================================
\section{Further Algorithmic Details of FAUST}
\label{sec:appendix_c}
%=============================================================================

\subsection{Extension to FedOpt}
\label{sec:fedopt}

Although~\citet{cheng2025convergencedistributedadaptiveoptimization} show provable convergence for adaptive inner optimizers in a
federated optimization framework, their result rests on the assumption that after a period of local
work, the new global model is created by averaging the local client models. In relation to the larger
FedOpt literature~\citep{reddi2021adaptivefederatedoptimization}, the scheme chosen by~\citet{cheng2025convergencedistributedadaptiveoptimization} resembles
that of FedAvg, or where the server optimizer is SGD with the outer learning rate set to one~\citep{reddi2021adaptivefederatedoptimization}.
Naturally, the question of whether alternate server optimizers that have been used in prior works
can also be implemented for Local Adam, and thus FAUST, arises.

We argue that indeed FAUST's principles can be effectively applied to any FedOpt method and
not just FedAvg. While using an alternate server optimizer does not have proven convergence
guarantees as yet, we show in Algorithm~\ref{alg:desloc} that the choice of the $\mathrm{SERVEROPT}$ is not constrained
from a practical point-of-view. However, the improvements that FAUST provides are related to
the local optimization procedure, which is orthogonal to the outer optimizer choice.
Choosing the correct, and most effective, outer optimizer is an open research area~\citep{khaled2025understandingouteroptimizerslocal},
and we leave the investigations of the interactions between FAUST and outer optimizers to future work.

Importantly, the compile-time sequence-parallel decomposition introduced in FAUST
(Section~\ref{sec:neuronsp}) does not interact with the outer optimizer at all: the SP
All-to-All collectives execute within every forward and backward step to ensure attention
correctness, whereas $\mathrm{SERVEROPT}$ operates exclusively at $K_x$-boundary steps on the
already-reduced parameter averages $\mathbb{E}_m[x^m_t]$.
The SP graph rewrite and the FedOpt outer update thus act on disjoint stages of the training
loop and compose without interference.


\subsection{Extension to FedOpt}
\label{sec:fedopt2}

Although~\citet{cheng2025convergencedistributedadaptiveoptimization} show provable convergence for adaptive inner optimizers in a
federated optimization framework, their result rests on the assumption that after a period of local
work, the new global model is created by averaging the local client models. In relation to the larger
\texttt{FedOpt} literature~\citep{reddi2021adaptivefederatedoptimization}, the scheme chosen by~\citet{cheng2025convergencedistributedadaptiveoptimization} resembles
that of \texttt{FedAvg}, or where the server optimizer is \texttt{SGD} with the outer learning rate set to one~\citep{reddi2021adaptivefederatedoptimization}.
Naturally, the question of whether alternate server optimizers than have been used in prior works
can also be implemented for \texttt{Local Adam} , arises.

We argue that indeed principles can be effectively applied to any \texttt{FedOpt} method and
not just \texttt{FedAvg}. While using an alternate server optimizer does not have proven convergence
guarantees as yet, we show in Algorithm~\ref{alg:desloc} that the choice of the \texttt{ServerOpt} is not constrained
from a practical point-of-view. However, the improvements that \texttt{DES-LOC} provide are related to
the local optimization procedure, which is orthogonal to the outer optimizer choice.
Choosing the correct, and most effective, outer optimizer is an open research area~\citep{khaled2025understandingouteroptimizerslocal},
and we leave the investigations of the interactions between \texttt{DES-LOC} and outer optimizers to future work.


\subsection{Deterministic Optimizer-specific Variants of Algorithm~\ref{alg:desloc}}
\label{sec:algorithm_variants}

We now present the concrete instantiation of the generic FAUST algorithm
(Algorithm~\ref{alg:desloc}) for the Adam inner optimizer.
Unlike prior formulations that treat the compile-time SP pass as an implicit
black box---merely annotating the gradient with ``SP-reduced''---we make the
intra-step All-to-All collective and its interaction with the inter-step
data-parallel synchronization a \emph{first-class citizen} of the algorithm.
This explicit coupling is both necessary and novel:
it exposes the \textbf{collective orthogonality constraint}
($\mathcal{A}_P$ must retire before $\mathbb{E}_m$ fires),
enables the SP-tightened clipping bound $\rho_{\text{eff}} = \rho / \sqrt{P}$
for the $d_{\text{attn}}$ attention-layer coordinates~\eqref{eq:sp_var_reduction},
and formalizes the structural tiering (embedding/norm $\to$ $K_x$-tier,
attention projections $\to$ $K_u$-tier, FFN $\to$ $K_v$-tier) that FAUST's
implementation deploys via the \texttt{DeslocTieredAllReduce} backend.

The $\mathbb{E}_m[\cdot]$ averaging operator denotes the cross-replica AllReduce over the
data-parallel group; when ZeRO-1 optimizer-state partitioning is active, this reduces to a
reduce-scatter followed by a per-shard update and an all-gather, gated by the corresponding
synchronization period ($K_x$, $K_u$, or $K_v$).
The All-to-All operator $\mathcal{A}_P$ corresponds to the Ulysses head-splitting
collective: for input $Q,K,V \in \mathbb{R}^{B \times N \times S/P \times H}$,
it produces $[B, N/P, S, H]$ by scattering along the head dimension and gathering
along the sequence dimension; the inverse $\mathcal{A}_P^{-1}$ restores the original
layout after SDPA. In our compiler implementation,
$\mathcal{A}_P$ is realized as \texttt{torch.ops.autosp.all\_to\_all} with
autograd support for the backward pass.

\begin{algorithm}[ht]
\caption{FAUST-Adam with Explicit SP--DP Collective Orchestration}
\label{alg:desloc_adam}
\begin{algorithmic}[1]
\REQUIRE Model $x_0 \in \mathbb{R}^d$; first/second moment seeds $u_{-1}, v_{-1} \in \mathbb{R}^d$
\REQUIRE $\{\eta_t\}_{t=0}^{T-1} \subset \mathbb{R}_{>0}$ \COMMENT{$K_x$-aligned WSD step-size schedule}
\REQUIRE $\beta_1, \beta_2 \in [0,1)$; $\rho, \lambda \in \mathbb{R}_{>0}$ \COMMENT{Adam decays; clip bound; stability}
\REQUIRE $T, M \in \mathbb{N}_+$ \COMMENT{total iterations; DP workers}
\REQUIRE $P \in \mathbb{N}_+$ \COMMENT{SP degree ($P = 1$ recovers non-SP)}
\REQUIRE $K_x, K_u, K_v \in \mathbb{N}_+$ with $K_x \mid K_u \mid K_v$ \COMMENT{tier sync periods, divisibility chain}
\REQUIRE $L$ \COMMENT{number of attention (SDPA) layers}
\ENSURE $x_T, u_{T-1}, v_{T-1}$

\STATE $\forall m$: $x^m_0 \gets x_0$;\; $u^m_{-1} \gets v^m_{-1} \gets 0$
\STATE $\mathcal{H} \gets \emptyset$ \COMMENT{set of pending A2A async handles}
\FOR{$t = 0, \ldots, T-1$}
    \FORALL{$m = 0, \ldots, M-1$ \textbf{in parallel}}
        \STATE \COMMENT{\textbf{--- Phase 1: Intra-step SP forward (compile-time inserted) ---}}
        \STATE Partition mini-batch $\xi^m_t$ into $P$ chunks along sequence dim
        \FOR{$\ell = 1, \ldots, L$}
            \STATE $Q_\ell, K_\ell, V_\ell \gets \mathcal{A}_P\!\bigl(Q_\ell^{\text{loc}}, K_\ell^{\text{loc}}, V_\ell^{\text{loc}}\bigr)$ \COMMENT{$[B,N,\tfrac{S}{P},H] \!\to\! [B,\tfrac{N}{P},S,H]$}
            \STATE $O_\ell \gets \mathrm{SDPA}(Q_\ell, K_\ell, V_\ell)$ \COMMENT{full-$S$ attention on $N/P$ heads}
            \STATE $O_\ell^{\text{loc}} \gets \mathcal{A}_P^{-1}(O_\ell)$;\quad $\mathcal{H} \gets \mathcal{H} \cup \{h_\ell\}$ \COMMENT{async handle}
        \ENDFOR
        \STATE \COMMENT{\textbf{--- Phase 2: Backward with FP32 A2A accumulation ---}}
        \STATE $\tilde{g}^m_t \gets \nabla_x \mathcal{L}\bigl(x^m_t;\, \xi^m_t,\, \{O_\ell^{\text{loc}}\}\bigr)$ \COMMENT{backward through compiled graph $\mathcal{G}'$}
        \STATE $g^m_t \gets \sum_{p=0}^{P-1} \tilde{g}^{m,p}_t$ \COMMENT{backward A2A reduces partial grads; FP32 accum.}
        \STATE $\hat{g}^m_t \gets \mathrm{clip}\!\bigl(g^m_t,\, \rho_j\bigr)$ where $\rho_j = \begin{cases} \rho\,/\!\sqrt{P} & j \in \mathcal{I}_{\text{attn}} \\ \rho & \text{otherwise}\end{cases}$ \COMMENT{SP-tightened clip}
        \STATE \COMMENT{\textbf{--- Phase 3: Collective fence ---}}
        \IF{$t \bmod K_x = 0$}
            \STATE \textsc{WaitAll}$(\mathcal{H})$;\; $\mathcal{H} \gets \emptyset$ \COMMENT{fence: all A2A handles must retire before DP sync}
        \ENDIF
        \STATE \COMMENT{\textbf{--- Phase 4: Inter-step DP optimizer with tiered sync ---}}
        \IF{$t \bmod K_u = 0$}
            \STATE $u^m_t \gets \beta_1 \mathbb{E}_m[u^m_{t-1}] + (1-\beta_1)\hat{g}^m_t$ \COMMENT{AllReduce first moment}
        \ELSE
            \STATE $u^m_t \gets \beta_1 u^m_{t-1} + (1-\beta_1)\hat{g}^m_t$ \COMMENT{local EMA}
        \ENDIF

        \IF{$t \bmod K_v = 0$}
            \STATE $v^m_t \gets \beta_2 \mathbb{E}_m[v^m_{t-1}] + (1-\beta_2)(\hat{g}^m_t \odot \hat{g}^m_t)$ \COMMENT{AllReduce second moment}
        \ELSE
            \STATE $v^m_t \gets \beta_2 v^m_{t-1} + (1-\beta_2)(\hat{g}^m_t \odot \hat{g}^m_t)$ \COMMENT{local EMA}
        \ENDIF

        \STATE $d^m_t \gets \dfrac{\eta_t}{\sqrt{v^m_t+\lambda^2}} \odot u^m_t$ \COMMENT{preconditioned step}

        \IF{$t \bmod K_x = 0$}
            \STATE $x^m_{t+1} \gets \mathbb{E}_m[x^m_t] - d^m_t$ \COMMENT{AllReduce parameters}
        \ELSE
            \STATE $x^m_{t+1} \gets x^m_t - d^m_t$ \COMMENT{local step}
        \ENDIF
    \ENDFOR
\ENDFOR
\end{algorithmic}
\end{algorithm}

\paragraph{Relationship between Algorithm~\ref{alg:desloc_adam} and the AutoSP compiler.}
The four phases above map directly to the AutoSP compilation pipeline
(commit \texttt{5efb24a}):
Phase~1 is realized by the five compiler passes
\textsc{ShardSeqDim}~$\to$~\textsc{ShardInputIds}~$\to$~\textsc{ShardLabelIds}~$\to$~\textsc{ShardPositionIds}~$\to$~\textsc{InsertA2A},
which rewrite the FX graph $\mathcal{G}$ into $\mathcal{G}'$ by inserting
$4L$ All-to-All nodes (one scatter-heads/gather-seq before each of $Q,K,V$,
plus one scatter-seq/gather-heads after $O$, per layer).
Phase~2 exploits PyTorch's autograd: the registered backward for
\texttt{autosp::all\_to\_all} swaps scatter and gather indices, so
$\nabla \mathcal{A}_P = \mathcal{A}_P^{-1}$, yielding the FP32-accumulated
partial-gradient summation.
Phase~3 corresponds to the \texttt{\_desloc\_fence\_all\_pending} routine,
which blocks until all $|\mathcal{H}|$ asynchronous handles complete---enforcing
the orthogonality invariant that no DP AllReduce observes a partially-reduced
SP gradient.
Phase~4 is the standard tiered EMA update with bucket-level AllReduce gating,
where each gradient bucket is tagged with its structural tier
(embedding/normalization~$\to$~$K_x$, attention projections~$\to$~$K_u$,
FFN~$\to$~$K_v$) as implemented in \texttt{DeslocTieredAllReduceTorch}.

\begin{algorithm}[ht]
\caption{FAUST-ADOPT with Explicit SP--DP Collective Orchestration}
\label{alg:desloc_adopt}
\begin{algorithmic}[1]
\REQUIRE Model $x_0 \in \mathbb{R}^d$; moment seeds $m_{-1}, v_{-1} \in \mathbb{R}^d$
\REQUIRE $\{\eta_t\}_{t=0}^{T-1} \subset \mathbb{R}_{>0}$ \COMMENT{$K_x$-aligned WSD learning rate schedule}
\REQUIRE $\beta_1, \beta_2 \in [0,1)$; $\rho, \epsilon \in \mathbb{R}_{>0}$ \COMMENT{decay; clip bound; stability}
\REQUIRE $T, M \in \mathbb{N}_+$ \COMMENT{total iterations; DP workers}
\REQUIRE $P \in \mathbb{N}_+$; $L \in \mathbb{N}_+$ \COMMENT{SP degree; number of SDPA layers}
\REQUIRE $K_x, K_m, K_v \in \mathbb{N}_+$ with $K_x \mid K_m \mid K_v$ \COMMENT{tier sync periods}
\ENSURE $x_T$, $m_{T-1}$, $v_{T-1}$

\STATE $\forall w$: $x^w_0 \gets x_0$;\; $m^w_{-1} \gets v^w_{-1} \gets 0$
\STATE $\mathcal{H} \gets \emptyset$ \COMMENT{pending A2A async handles}
\FOR{$t = 0, \ldots, T-1$}
    \FORALL{$w = 0, \ldots, M-1$ \textbf{in parallel}}
        \STATE \COMMENT{\textbf{--- Phase 1: Intra-step SP forward ---}}
        \STATE Partition $\xi^w_t$ into $P$ chunks along sequence dim \COMMENT{compiled $\mathcal{G}'$}
        \FOR{$\ell = 1, \ldots, L$}
            \STATE $Q_\ell, K_\ell, V_\ell \gets \mathcal{A}_P(Q_\ell^{\text{loc}}, K_\ell^{\text{loc}}, V_\ell^{\text{loc}})$ \COMMENT{$[B,N,\tfrac{S}{P},H] \!\to\! [B,\tfrac{N}{P},S,H]$}
            \STATE $O_\ell \gets \mathrm{SDPA}(Q_\ell, K_\ell, V_\ell)$;\quad $O_\ell^{\text{loc}} \gets \mathcal{A}_P^{-1}(O_\ell)$;\quad $\mathcal{H} \gets \mathcal{H} \cup \{h_\ell\}$
        \ENDFOR
        \STATE \COMMENT{\textbf{--- Phase 2: Backward with FP32 A2A + SP-tightened clip ---}}
        \STATE $g^w_t \gets \sum_{p=0}^{P-1} \tilde{g}^{w,p}_t$ \COMMENT{backward A2A: $\nabla\mathcal{A}_P = \mathcal{A}_P^{-1}$; FP32 accum.}
        \STATE $\hat{g}^w_t \gets \mathrm{clip}(g^w_t,\, \rho_j)$ where $\rho_j = \rho\,/\!\sqrt{P}$ for $j \in \mathcal{I}_{\text{attn}}$, $\rho$ otherwise
        \STATE \COMMENT{\textbf{--- Phase 3: Collective fence ---}}
        \IF{$t \bmod K_x = 0$}
            \STATE \textsc{WaitAll}$(\mathcal{H})$;\; $\mathcal{H} \gets \emptyset$ \COMMENT{A2A $\perp$ DP: fence before any AllReduce}
        \ENDIF
        \STATE \COMMENT{\textbf{--- Phase 4: ADOPT optimizer with tiered DP sync ---}}
        \IF{$t \bmod K_v = 0$}
            \STATE $v^w_t \gets \beta_2 \mathbb{E}_w[v^w_{t-1}] + (1 - \beta_2)(\hat{g}^w_t \odot \hat{g}^w_t)$ \COMMENT{AllReduce $v$}
        \ELSE
            \STATE $v^w_t \gets \beta_2 v^w_{t-1} + (1 - \beta_2)(\hat{g}^w_t \odot \hat{g}^w_t)$ \COMMENT{local EMA}
        \ENDIF
        \IF{$t \bmod K_m = 0$}
            \STATE $m^w_t \gets \beta_1 \mathbb{E}_w[m^w_{t-1}] + (1 - \beta_1) \dfrac{\hat{g}^w_t}{\max\bigl\{\sqrt{v^w_{t-1}},\, \epsilon\bigr\}}$ \COMMENT{AllReduce $m$}
        \ELSE
            \STATE $m^w_t \gets \beta_1 m^w_{t-1} + (1 - \beta_1) \dfrac{\hat{g}^w_t}{\max\bigl\{\sqrt{v^w_{t-1}},\, \epsilon\bigr\}}$ \COMMENT{local ADOPT}
        \ENDIF
        \STATE $d^w_t \gets \eta_t\, m^w_t$ \COMMENT{ADOPT step (no per-coord preconditioning)}
        \IF{$t \bmod K_x = 0$}
            \STATE $x^w_{t+1} \gets \mathbb{E}_w[x^w_t] - d^w_t$ \COMMENT{AllReduce parameters}
        \ELSE
            \STATE $x^w_{t+1} \gets x^w_t - d^w_t$ \COMMENT{local step}
        \ENDIF
    \ENDFOR
\ENDFOR
\end{algorithmic}
\end{algorithm}

\paragraph{Note on ADOPT's update order.}
ADOPT~\citep{taniguchi2024adoptmodifiedadamconverge} differs from Adam in that the
second moment $v^w_t$ is updated \emph{before} the first moment $m^w_t$,
and the preconditioner uses $v^w_{t-1}$ (the previous step) rather than $v^w_t$.
In the FAUST setting this ordering is preserved per-worker: the tiered sync
gates on $K_v$ and $K_m$ are independent, so $v$ can synchronize at a different
cadence than $m$ without violating ADOPT's convergence requirements.
However, the divisibility chain $K_x \mid K_m \mid K_v$ ensures that whenever
parameters sync ($K_x$ boundary), both moments have already been averaged,
preventing stale-preconditioner drift.

%=============================================================================
% Unified analysis of the SP--DP collective interaction
%=============================================================================

\paragraph{SP--DP collective orthogonality and ZeRO-1 composition.}
A central invariant of both Algorithms~\ref{alg:desloc_adam} and~\ref{alg:desloc_adopt}
is the \textbf{collective orthogonality constraint}: the intra-step SP collectives
($\mathcal{A}_P$ and $\mathcal{A}_P^{-1}$, Phase~1--2) and the inter-step DP
collectives ($\mathbb{E}_m$, Phase~4) must never overlap on the same gradient tensor.
This is enforced by the Phase~3 fence: at every $K_x$-boundary step,
$\textsc{WaitAll}(\mathcal{H})$ blocks until all $|\mathcal{H}| = L$ asynchronous
All-to-All handles have retired, guaranteeing that the gradient $\hat{g}^m_t$ entering
Phase~4 is fully sequence-reduced.

When ZeRO-1 optimizer-state partitioning is active, the $\mathbb{E}_m[\cdot]$ operator
decomposes into three sub-operations:
\begin{equation}
\mathbb{E}_m[s] \;\longmapsto\; \texttt{reduce\_scatter}(s) \;\to\; \text{per-shard update} \;\to\; \texttt{all\_gather}(\hat{s}),
\label{eq:zero1_decomposition}
\end{equation}
where the tier gate ($t \bmod K_j = 0$) controls whether
\texttt{reduce\_scatter} fires or the gradient is accumulated locally.
In the Neuron\_SP implementation, each gradient bucket is tagged with its structural
tier (embedding/norm $\to$ $K_x$, attention projections $\to$ $K_u$, FFN $\to$ $K_v$),
and the bucket's AllReduce is dispatched only when the period $K_j$ divides $t$.
The SP fence (Phase~3) serializes \emph{before} the first bucket dispatch, so the
ZeRO-1 \texttt{reduce\_scatter} never observes a partially SP-reduced gradient shard.

AutoSP interacts with this decomposition through a critical engine-level routing decision:
whereas DeepCompile normally bypasses the standard \texttt{allreduce\_gradients} path
(handling gradient reduction within the compiled graph), AutoSP explicitly re-enables it
so that the DP AllReduce runs \emph{outside} the compiled graph on the already SP-reduced
gradient. This ensures that the SP $\mathcal{A}_P^{-1}$ and the DP $\mathbb{E}_m$ execute
in strictly sequential phases, composing correctly even when the DP reduction is
bucket-pipelined.

\paragraph{SP gradient identity and convergence invariance.}
Under the Ulysses head-splitting strategy realized by AutoSP's
\textsc{InsertA2A} pass, each SP rank $p$ on device $m$ computes a partial
gradient $\tilde{g}^{m,p}_t$ from its local sequence chunk.
The backward All-to-All---implemented via the registered autograd hook
\texttt{\_all\_to\_all\_backward} which swaps \texttt{scatter\_idx} and
\texttt{gather\_idx}---produces the aggregate:
\begin{equation}
g^m_t = \sum_{p=0}^{P-1} \tilde{g}^{m,p}_t = \nabla_x \mathcal{L}(x^m_t;\, \xi^m_t),
\label{eq:sp_gradient_identity}
\end{equation}
which is \emph{identical} to the gradient from a non-SP forward-backward pass on the
full sequence $S$.
This identity holds exactly for linear layers and attention (where the All-to-All
implements a head-dimension repartitioning), and up to floating-point reduction
order for LayerNorm parameters (which require an additional AllReduce of partial
gradients within the SP group, implemented in the backward hook).
Since~\eqref{eq:sp_gradient_identity} is independent of $P$, the optimizer loops
in Algorithms~\ref{alg:desloc_adam} and~\ref{alg:desloc_adopt} operate on the same
$g^m_t$ for any SP degree, and the convergence analyses in
Sections~\ref{sec:sgdm_proof} and~\ref{sec:adam_extension} apply without modification.

\paragraph{Selective activation checkpointing under SP.}
AutoSP introduces a compile-time heuristic for activation checkpointing that
is SP-aware: the \textsc{LongContextCheckpointing} pass modifies PyTorch's
\texttt{solve\_min\_cut} partitioner to distinguish between operations whose
activation memory scales quadratically with sequence length (SDPA, convolutions)
and those that scale linearly (matmul, FFN).
Specifically, the heuristic bans recomputation of any operation whose output
tensor has a second dimension $\geq 4096$, as these are the attention
activations that dominate memory under long-context SP training.
Matmul and related linear operations are \emph{not} banned, allowing the
partitioner to recompute them during the backward pass at negligible cost.
This selective strategy provides $O(L \cdot S)$ memory savings compared to
full activation storage, while adding only $O(L \cdot d_{\text{model}})$ recomputation
cost---a ratio that improves as the sequence length $S$ grows, making it
particularly effective in the SP regime where $S$ is large by design.

%=============================================================================
% Appendix D: Convergence Analysis of Neuron_SP-SGDM
%=============================================================================

\section{Convergence Analysis of FAUST-SGDM (in expectation bounds)}
\label{sec:sgdm_proof}

We provide a complete, self-contained convergence analysis of the FAUST
system applied to distributed SGDM with a single optimizer state ($N=1$,
momentum). The analysis proceeds in four steps.

Throughout this section we use the following notation.
There are $M$ replicas, indexed by $m \in \{1, \ldots, M\}$.
Each replica $m$ maintains a local copy of the parameters $x_t^m \in \mathbb{R}^d$
and a local momentum buffer $u_t^m \in \mathbb{R}^d$.
We denote the SP degree by $P \geq 1$.

\paragraph{Why $P$ does not affect the convergence rate (and where it helps).}
The compile-time sequence-parallel graph rewrite (AutoSP,
Section~\ref{sec:FAUST}) inserts $4L$ All-to-All collectives that
partition the sequence dimension across $P$ ranks within each
replica. By the SP gradient identity~\eqref{eq:sp_gradient_identity},
the gradient entering the optimizer is
$g_t^m = \sum_{p=0}^{P-1} \tilde{g}_t^{m,p} = \nabla F_m(x_t^m;\,\xi_t^m)$,
identical to the non-SP case.
The stochastic variance bound therefore satisfies
\begin{equation}
\mathbb{E}\bigl\|g_t^m - \nabla f_m(x_t^m)\bigr\|^2 \;\leq\; \sigma^2
\label{eq:sp_variance_invariance}
\end{equation}
with the \emph{same} constant $\sigma^2$ for any $P \geq 1$, and the entire
proof chain below is $P$-independent.

However, $P$ does improve the \emph{per-coordinate} gradient variance for the
$d_{\text{attn}} = 4NH$ attention-layer coordinates.
By the SP variance reduction~\eqref{eq:sp_var_reduction}, the
per-coordinate variance on these coordinates is $\sigma_j^2/P$ instead of
$\sigma_j^2$. This motivates the SP-tightened clipping bound
$\rho_{\text{eff}} = \rho/\sqrt{P}$ used in
Algorithm~\ref{alg:desloc_adam} (Phase~2), which reduces the drift
contribution from attention parameters by a factor of $1/\sqrt{P}$---a
tighter constant that refines the bound~\eqref{eq:drift_u_sp}
without changing the overall convergence \emph{order}.

The collective orthogonality constraint (Phase~3 fence in
Algorithm~\ref{alg:desloc_adam}) ensures that the gradient $g_t^m$
is fully sequence-reduced before any data-parallel synchronization
($\mathbb{E}_m$) fires. This is critical: if the fence were absent,
the DP AllReduce would average partially-reduced SP gradients across
replicas, violating the unbiasedness assumption that undergirds
the entire analysis.

The complete algorithm is given below.

\begin{algorithm}[ht]
\caption{FAUST-SGDM (probabilistic synchronization formulation)}
\label{alg:desloc_sgdm}
\begin{algorithmic}[1]
\REQUIRE \textbf{Model tensors}
\STATE $x_0 \in \mathbb{R}^d$ \hfill initial parameter vector
\STATE $u_{-1} \in \mathbb{R}^d$ \hfill seed for the momentum, initialised to $\mathbf{0}$
\REQUIRE \textbf{Hyper-parameters}
\STATE $\{\eta_t\}_{t=0}^{T-1} \subset \mathbb{R}_{>0}$ \hfill $K_x$-aligned WSD step-size schedule
\STATE $\beta \in [0,1)$ \hfill momentum decay factor
\STATE $T \in \mathbb{N}_+$ \hfill total optimisation steps
\STATE $M \in \mathbb{N}_+$ \hfill number of replicas
\STATE $p_x = \frac{1}{K_x},\; p_u = \frac{1}{K_u} \in [0,1]$ \hfill synchronization probabilities for parameters and momentum
\ENSURE $x_T$, $u_{T-1}^m$
\STATE \textbf{for each replica} $m$: $x_0^m = x_0,\; u_{-1}^m = 0$ \hfill local init
\FOR{$t = 0, \ldots, T-1$}
    \FOR{all replicas $m = 0, \ldots, M-1$ \textbf{in parallel}}
        \STATE $g_t^m \leftarrow \nabla F_m(x_t^m; \xi_t^m)$ \hfill stochastic gradient (SP-reduced when $P > 1$)
        \STATE $u_t^m \leftarrow \begin{cases} \mathbb{E}_m[\beta u_{t-1}^m + (1-\beta) g_t^m], & \text{with probability } p_u \\ \beta u_{t-1}^m + (1-\beta) g_t^m, & \text{with probability } 1-p_u \end{cases}$
        \STATE $x_{t+1}^m \leftarrow \begin{cases} \mathbb{E}_m[x_t^m - \eta_t u_t^m], & \text{with probability } p_x \\ x_t^m - \eta_t u_t^m, & \text{with probability } 1-p_x \end{cases}$
    \ENDFOR
\ENDFOR
\end{algorithmic}
\end{algorithm}

To facilitate the analysis, we model each synchronization event as a Bernoulli
trial. That is, instead of averaging the parameters $x_t^m$ at every step that
satisfies $t \bmod K_x = 0$, we average with probability
\begin{equation}
\label{eq:prob_equiv}
p_x \;=\; \frac{1}{K_x}
\end{equation}
at \emph{every} step. The expected number of averaging events in $T$ steps under
the probabilistic model is $p_x \cdot T = T / K_x$, which equals the number of
averaging events in the deterministic model. The two formulations are therefore
statistically equivalent. Likewise, the momentum $u_t^m$ is synchronized with
probability $p_u = 1/K_u$, independently of $p_x$.

%=============================================================================
% STEP 1
%=============================================================================
\paragraph{Step 1: Averaged dynamics and virtual iterates.}

\noindent\textbf{Definition of averaged quantities.}
For each step $t \geq 0$, we define the \emph{cross-replica averages} of the
parameters, momentum, and stochastic gradient:
\begin{align}
x_t &\;\stackrel{\text{def}}{=}\; \frac{1}{M}\sum_{m=1}^{M} x_t^m
     \;=\; \mathbb{E}_m[x_t^m], \label{eq:avg_x} \\[4pt]
u_t &\;\stackrel{\text{def}}{=}\; \frac{1}{M}\sum_{m=1}^{M} u_t^m
     \;=\; \mathbb{E}_m[u_t^m], \label{eq:avg_u} \\[4pt]
g_t &\;\stackrel{\text{def}}{=}\; \frac{1}{M}\sum_{m=1}^{M} g_t^m
     \;=\; \mathbb{E}_m[g_t^m]. \label{eq:avg_g}
\end{align}
The notation $\mathbb{E}_m[\cdot]$ denotes the uniform average over the
replica index $m$, \emph{not} an expectation over randomness. Expectation
over the stochastic sampling $\xi_t^m$ is written $\mathbb{E}[\cdot]$.

\medskip
\noindent\textbf{Claim: the averages follow centralized SGDM.}
We now show that $x_t$ and $u_t$ satisfy the standard (single-machine) SGDM
recursion, regardless of whether a synchronization event occurs at step $t$.

Consider the momentum update on replica $m$. With probability $p_u$ (sync),
\[
u_t^m = \beta\,\mathbb{E}_m[u_{t-1}^m] + (1-\beta)\,g_t^m
       = \beta\,u_{t-1} + (1-\beta)\,g_t^m,
\]
and with probability $1-p_u$ (no sync),
\[
u_t^m = \beta\,u_{t-1}^m + (1-\beta)\,g_t^m.
\]
In either case, averaging over $m$ gives
\begin{align}
u_t &= \mathbb{E}_m[u_t^m] \notag \\
    &= p_u \cdot \bigl(\beta\,u_{t-1} + (1-\beta)\,g_t\bigr)
       + (1-p_u) \cdot \bigl(\beta\,\mathbb{E}_m[u_{t-1}^m]
       + (1-\beta)\,g_t\bigr) \notag \\
    &= p_u \cdot \bigl(\beta\,u_{t-1} + (1-\beta)\,g_t\bigr)
       + (1-p_u) \cdot \bigl(\beta\,u_{t-1}
       + (1-\beta)\,g_t\bigr) \notag \\
    &= \bigl(p_u + (1-p_u)\bigr) \cdot
       \bigl(\beta\,u_{t-1} + (1-\beta)\,g_t\bigr) \notag \\
    &= 1 \cdot \bigl(\beta\,u_{t-1} + (1-\beta)\,g_t\bigr) \notag \\
    &= \beta\,u_{t-1} + (1-\beta)\,g_t. \label{eq:avg_u_recursion}
\end{align}
An identical argument for the parameter update yields
\begin{equation}
x_{t+1} = x_t - \eta\,u_t. \label{eq:avg_x_recursion}
\end{equation}
Equations~\eqref{eq:avg_u_recursion} and~\eqref{eq:avg_x_recursion} are
precisely the centralized SGDM updates.

\medskip
\noindent\textbf{Definition of virtual iterates.}
Set $x_{-1} \stackrel{\text{def}}{=} x_0$, and for each $t \geq 0$ define
\begin{equation}
z_t \;\stackrel{\text{def}}{=}\;
\frac{1}{1-\beta}\,x_t \;-\; \frac{\beta}{1-\beta}\,x_{t-1}.
\label{eq:virtual}
\end{equation}
At $t = 0$ this gives $z_0 = \frac{1}{1-\beta}\,x_0 - \frac{\beta}{1-\beta}\,x_0 = x_0$.

\medskip
\noindent\textbf{Key property: virtual iterates follow averaged gradients.}
For any $t \geq 0$, we compute $z_{t+1} - z_t$ step by step.
First, by the definition~\eqref{eq:virtual},
\begin{align}
z_{t+1} - z_t
&= \frac{1}{1-\beta}\,x_{t+1} - \frac{\beta}{1-\beta}\,x_t
   - \frac{1}{1-\beta}\,x_t + \frac{\beta}{1-\beta}\,x_{t-1} \notag \\
&= \frac{1}{1-\beta}\,(x_{t+1} - x_t)
   - \frac{\beta}{1-\beta}\,(x_t - x_{t-1}).
\label{eq:virt_diff_raw}
\end{align}
Now substitute the parameter recursion~\eqref{eq:avg_x_recursion}:
\begin{align}
x_{t+1} - x_t &= -\eta\,u_t, \label{eq:xdiff_t} \\
x_t - x_{t-1}  &= -\eta\,u_{t-1}. \label{eq:xdiff_tm1}
\end{align}
Inserting~\eqref{eq:xdiff_t} and~\eqref{eq:xdiff_tm1}
into~\eqref{eq:virt_diff_raw}:
\begin{align}
z_{t+1} - z_t
&= \frac{1}{1-\beta}\,(-\eta\,u_t)
   - \frac{\beta}{1-\beta}\,(-\eta\,u_{t-1}) \notag \\
&= -\frac{\eta}{1-\beta}\,u_t
   + \frac{\eta\,\beta}{1-\beta}\,u_{t-1} \notag \\
&= -\frac{\eta}{1-\beta}\,
   \bigl(u_t - \beta\,u_{t-1}\bigr).
\label{eq:virt_diff_factored}
\end{align}
Finally, by the momentum recursion~\eqref{eq:avg_u_recursion},
\begin{equation}
u_t - \beta\,u_{t-1}
= \bigl(\beta\,u_{t-1} + (1-\beta)\,g_t\bigr) - \beta\,u_{t-1}
= (1-\beta)\,g_t.
\label{eq:mom_minus}
\end{equation}
Substituting~\eqref{eq:mom_minus} into~\eqref{eq:virt_diff_factored}:
\begin{align}
z_{t+1} - z_t
&= -\frac{\eta}{1-\beta}\,(1-\beta)\,g_t \notag \\
&= -\eta\,g_t.
\label{eq:virt_gradient}
\end{align}
This is the key identity: the virtual iterates $z_t$ evolve by a simple
gradient step $-\eta\,g_t$, with no momentum contamination.

%=============================================================================
% STEP 2
%=============================================================================
\paragraph{Step 2: Smoothness descent on virtual iterates.}

By Assumption~\ref{asm:smooth}, $f$ is $L$-smooth, so for any
$a, b \in \mathbb{R}^d$,
\begin{equation}
f(b) \leq f(a) + \langle \nabla f(a),\, b - a \rangle
       + \frac{L}{2}\,\|b - a\|^2.
\label{eq:smooth_def}
\end{equation}
Applying~\eqref{eq:smooth_def} with $a = z_t$ and $b = z_{t+1}$:
\begin{equation}
f(z_{t+1}) \leq f(z_t) + \langle \nabla f(z_t),\, z_{t+1} - z_t \rangle
             + \frac{L}{2}\,\|z_{t+1} - z_t\|^2.
\label{eq:smooth_applied}
\end{equation}
We now split the inner product by adding and subtracting $\nabla f(x_t)$.
Write $\nabla f(z_t) = \nabla f(x_t) + \bigl(\nabla f(z_t) - \nabla f(x_t)\bigr)$.
Then
\begin{align}
\langle \nabla f(z_t),\, z_{t+1} - z_t \rangle
&= \langle \nabla f(x_t),\, z_{t+1} - z_t \rangle
   + \langle \nabla f(z_t) - \nabla f(x_t),\, z_{t+1} - z_t \rangle.
\label{eq:inner_split}
\end{align}
Substituting~\eqref{eq:inner_split} into~\eqref{eq:smooth_applied}:
\begin{align}
f(z_{t+1})
&\leq f(z_t)
  + \underbrace{\langle \nabla f(x_t),\, z_{t+1} - z_t \rangle}_{I}
  + \underbrace{\langle \nabla f(z_t) - \nabla f(x_t),\,
    z_{t+1} - z_t \rangle}_{II}
  + \underbrace{\frac{L}{2}\,\|z_{t+1} - z_t\|^2}_{III}.
\label{eq:three_terms}
\end{align}
We bound terms I, II, III separately below.

%=============================================================================
% STEP 3a
%=============================================================================
\paragraph{Step 3a: Bounding term $I$.}

Recall from~\eqref{eq:virt_gradient} that $z_{t+1} - z_t = -\eta\,g_t$.
We substitute and then use the identity $g_t = \mathbb{E}_m[g_t^m]$:
\begin{align}
\langle \nabla f(x_t),\, z_{t+1} - z_t \rangle
&= \langle \nabla f(x_t),\, -\eta\,g_t \rangle \notag \\
&= -\eta\,\langle \nabla f(x_t),\, g_t \rangle \notag \\
&= -\eta\,\left\langle \nabla f(x_t),\,
   \frac{1}{M}\sum_{m=1}^{M} g_t^m \right\rangle.
\label{eq:I_start}
\end{align}
Taking the expectation over the stochastic sampling and using unbiasedness
(Assumption~\ref{asm:noise}):
\begin{align}
\mathbb{E}\!\left[\left\langle \nabla f(x_t),\,
   \frac{1}{M}\sum_{m=1}^{M} g_t^m \right\rangle\right]
&= \left\langle \nabla f(x_t),\,
   \frac{1}{M}\sum_{m=1}^{M} \mathbb{E}[g_t^m] \right\rangle \notag \\
&= \left\langle \nabla f(x_t),\,
   \frac{1}{M}\sum_{m=1}^{M} \nabla f_m(x_t^m) \right\rangle.
\label{eq:I_unbiased}
\end{align}
Hence
\begin{equation}
\mathbb{E}[I]
= -\eta\,\mathbb{E}\!\left[\left\langle \nabla f(x_t),\,
   \frac{1}{M}\sum_{m=1}^{M} \nabla f_m(x_t^m) \right\rangle\right].
\label{eq:I_expect}
\end{equation}

We now apply the \emph{polarization identity}: for any vectors
$a, b \in \mathbb{R}^d$,
\begin{equation}
\langle a, b \rangle
= \frac{1}{2}\|a\|^2 + \frac{1}{2}\|b\|^2
  - \frac{1}{2}\|a - b\|^2.
\label{eq:polar}
\end{equation}
To use~\eqref{eq:polar} in the form that gives us negative squared norms
(which we need for convergence), we rewrite it as
\begin{equation}
-\langle a, b \rangle
= -\frac{1}{2}\|a\|^2 - \frac{1}{2}\|b\|^2
  + \frac{1}{2}\|a - b\|^2.
\label{eq:neg_polar}
\end{equation}
Setting $a = \nabla f(x_t)$ and
$b = \frac{1}{M}\sum_{m=1}^{M}\nabla f_m(x_t^m)$
in~\eqref{eq:neg_polar}:
\begin{align}
&-\left\langle \nabla f(x_t),\,
   \frac{1}{M}\sum_{m=1}^{M}\nabla f_m(x_t^m) \right\rangle \notag \\
&\quad= -\frac{1}{2}\,\|\nabla f(x_t)\|^2
  - \frac{1}{2}\left\|\frac{1}{M}\sum_{m=1}^{M}
    \nabla f_m(x_t^m)\right\|^2
  + \frac{1}{2}\left\|\nabla f(x_t)
    - \frac{1}{M}\sum_{m=1}^{M}\nabla f_m(x_t^m)\right\|^2.
\label{eq:I_polar}
\end{align}
Multiplying~\eqref{eq:I_polar} by $\eta$ and taking expectation:
\begin{align}
\mathbb{E}[I]
&= -\frac{\eta}{2}\,\mathbb{E}\|\nabla f(x_t)\|^2
  - \frac{\eta}{2}\,\mathbb{E}\left\|\frac{1}{M}\sum_{m=1}^{M}
    \nabla f_m(x_t^m)\right\|^2 \notag \\
&\quad+ \frac{\eta}{2}\,\mathbb{E}\left\|\nabla f(x_t)
    - \frac{1}{M}\sum_{m=1}^{M}\nabla f_m(x_t^m)\right\|^2.
\label{eq:I_three}
\end{align}

It remains to bound the third term. We use the \emph{variance decomposition}:
for any random vectors $X_1, \ldots, X_M$,
\begin{equation}
\left\|\frac{1}{M}\sum_{m=1}^{M}(X_m - \bar{X})\right\|^2
\leq \frac{1}{M}\sum_{m=1}^{M}\|X_m - \bar{X}\|^2
\quad\text{where}\quad
\bar{X} = \frac{1}{M}\sum_{m=1}^{M}X_m.
\label{eq:var_decomp}
\end{equation}
More precisely, since $\nabla f(x_t) = \frac{1}{M}\sum_{m=1}^{M}\nabla f_m(x_t)$, we write
\begin{align}
&\left\|\nabla f(x_t) - \frac{1}{M}\sum_{m=1}^{M}
  \nabla f_m(x_t^m)\right\|^2 \notag \\
&\quad= \left\|\frac{1}{M}\sum_{m=1}^{M}\bigl(
  \nabla f_m(x_t) - \nabla f_m(x_t^m)\bigr)\right\|^2 \notag \\
&\quad\leq \frac{1}{M}\sum_{m=1}^{M}
  \|\nabla f_m(x_t) - \nabla f_m(x_t^m)\|^2
\qquad\text{(Jensen's inequality)}.
\label{eq:jensen_grad}
\end{align}
By $L$-smoothness (Assumption~\ref{asm:smooth}),
$\|\nabla f_m(x_t) - \nabla f_m(x_t^m)\| \leq L\,\|x_t - x_t^m\|$, so
\begin{equation}
\|\nabla f_m(x_t) - \nabla f_m(x_t^m)\|^2
\leq L^2\,\|x_t - x_t^m\|^2.
\label{eq:smooth_sq}
\end{equation}
Combining~\eqref{eq:I_three},~\eqref{eq:jensen_grad},
and~\eqref{eq:smooth_sq}:
\begin{align}
\mathbb{E}[I]
&\leq -\frac{\eta}{2}\,\mathbb{E}\|\nabla f(x_t)\|^2
  - \frac{\eta}{2}\,\mathbb{E}\left\|\frac{1}{M}\sum_{m=1}^{M}
    \nabla f_m(x_t^m)\right\|^2 \notag \\
&\quad+ \frac{\eta\,L^2}{2M}\sum_{m=1}^{M}
  \underbrace{\mathbb{E}\|x_t - x_t^m\|^2}_{\text{bounded in Lemma~\ref{lem:bound}}}.
\label{eq:I_final}
\end{align}

%=============================================================================
% STEP 3b
%=============================================================================
\paragraph{Step 3b: Bounding term $II$.}

Recall $z_{t+1} - z_t = -\eta\,g_t$ and
$g_t = \frac{1}{M}\sum_{m=1}^{M} g_t^m$. We begin:
\begin{align}
\mathbb{E}[II]
&= \mathbb{E}\bigl[\langle \nabla f(z_t) - \nabla f(x_t),\,
   z_{t+1} - z_t \rangle\bigr] \notag \\
&= -\eta\,\mathbb{E}\bigl[\langle \nabla f(z_t) - \nabla f(x_t),\,
   g_t \rangle\bigr] \notag \\
&= -\eta\,\mathbb{E}\!\left[\left\langle \nabla f(z_t) - \nabla f(x_t),\,
   \frac{1}{M}\sum_{m=1}^{M}\nabla f_m(x_t^m)\right\rangle\right].
\label{eq:II_start}
\end{align}
By Young's inequality, for any $\rho > 0$ and vectors $a, b$,
\begin{equation}
|\langle a, b \rangle|
\leq \frac{\rho}{2}\,\|a\|^2 + \frac{1}{2\rho}\,\|b\|^2.
\label{eq:young}
\end{equation}
Applying~\eqref{eq:young} with
$a = \nabla f(z_t) - \nabla f(x_t)$ and
$b = \frac{1}{M}\sum_{m=1}^{M}\nabla f_m(x_t^m)$:
\begin{align}
\mathbb{E}[II]
&\leq \frac{\eta\,\rho}{2}\,
  \mathbb{E}\|\nabla f(z_t) - \nabla f(x_t)\|^2
  + \frac{\eta}{2\rho}\,\mathbb{E}\left\|\frac{1}{M}\sum_{m=1}^{M}
    \nabla f_m(x_t^m)\right\|^2.
\label{eq:II_young}
\end{align}
By $L$-smoothness, $\|\nabla f(z_t) - \nabla f(x_t)\| \leq L\,\|z_t - x_t\|$,
hence $\|\nabla f(z_t) - \nabla f(x_t)\|^2 \leq L^2\,\|z_t - x_t\|^2$.
Therefore:
\begin{align}
\mathbb{E}[II]
&\leq \frac{\eta\,\rho\,L^2}{2}\,
  \underbrace{\mathbb{E}\|z_t - x_t\|^2}_{\text{bounded in Lemma~\ref{lem:virt_drift}}}
  + \frac{\eta}{2\rho}\,\mathbb{E}\left\|\frac{1}{M}\sum_{m=1}^{M}
    \nabla f_m(x_t^m)\right\|^2.
\label{eq:II_final}
\end{align}

%=============================================================================
% STEP 3c
%=============================================================================
\paragraph{Step 3c: Bounding term $III$.}

Using~\eqref{eq:virt_gradient}:
\begin{align}
\frac{L}{2}\,\|z_{t+1} - z_t\|^2
&= \frac{L}{2}\,\|-\eta\,g_t\|^2 \notag \\
&= \frac{\eta^2\,L}{2}\,\|g_t\|^2 \notag \\
&= \frac{\eta^2\,L}{2}\,\left\|\frac{1}{M}\sum_{m=1}^{M}g_t^m\right\|^2.
\label{eq:III_start}
\end{align}
We add and subtract $\nabla f_m(x_t^m)$ inside the sum. Define
$\varepsilon_t^m \stackrel{\text{def}}{=} g_t^m - \nabla f_m(x_t^m)$
(the noise term). Then
\begin{equation}
g_t^m = \nabla f_m(x_t^m) + \varepsilon_t^m,
\label{eq:noise_def}
\end{equation}
and
\begin{align}
\left\|\frac{1}{M}\sum_{m=1}^{M}g_t^m\right\|^2
&= \left\|\frac{1}{M}\sum_{m=1}^{M}\varepsilon_t^m
   + \frac{1}{M}\sum_{m=1}^{M}\nabla f_m(x_t^m)\right\|^2.
\label{eq:III_split}
\end{align}
By the identity $\|a + b\|^2 = \|a\|^2 + \|b\|^2 + 2\langle a, b\rangle$,
and noting that $\mathbb{E}[\varepsilon_t^m] = 0$ by
Assumption~\ref{asm:noise} (so the cross-term vanishes under expectation):
\begin{align}
\mathbb{E}\left\|\frac{1}{M}\sum_{m=1}^{M}g_t^m\right\|^2
&= \mathbb{E}\left\|\frac{1}{M}\sum_{m=1}^{M}\varepsilon_t^m\right\|^2
   + \mathbb{E}\left\|\frac{1}{M}\sum_{m=1}^{M}
     \nabla f_m(x_t^m)\right\|^2.
\label{eq:III_decomp}
\end{align}
For the noise term, since the $\varepsilon_t^m$ are independent across replicas
with zero mean:
\begin{align}
\mathbb{E}\left\|\frac{1}{M}\sum_{m=1}^{M}\varepsilon_t^m\right\|^2
&= \frac{1}{M^2}\sum_{m=1}^{M}\mathbb{E}\|\varepsilon_t^m\|^2 \notag \\
&= \frac{1}{M^2}\sum_{m=1}^{M}
   \mathbb{E}\|g_t^m - \nabla f_m(x_t^m)\|^2 \notag \\
&\leq \frac{1}{M^2}\sum_{m=1}^{M}\sigma^2
\qquad\text{(Assumption~\ref{asm:noise})} \notag \\
&= \frac{\sigma^2}{M}.
\label{eq:noise_bound}
\end{align}
Substituting~\eqref{eq:III_decomp} and~\eqref{eq:noise_bound}
into~\eqref{eq:III_start}:
\begin{align}
\mathbb{E}[III]
&\leq \frac{\eta^2\,L}{2M}\,\sigma^2
  + \frac{\eta^2\,L}{2}\,\mathbb{E}\left\|\frac{1}{M}\sum_{m=1}^{M}
    \nabla f_m(x_t^m)\right\|^2.
\label{eq:III_final}
\end{align}

%=============================================================================
% STEP 3abc combined
%=============================================================================
\paragraph{Step 3 combined: assembling the one-step bound.}

Substituting~\eqref{eq:I_final},~\eqref{eq:II_final},
and~\eqref{eq:III_final} into~\eqref{eq:three_terms} and taking
expectation:
\begin{align}
&\mathbb{E}[f(z_{t+1})] - \mathbb{E}[f(z_t)] \notag \\
&\quad\leq \underbrace{
  -\frac{\eta}{2}\,\mathbb{E}\|\nabla f(x_t)\|^2
}_{\text{from }I}
\;\underbrace{
  -\frac{\eta}{2}\,\mathbb{E}\left\|\frac{1}{M}\sum_m
    \nabla f_m(x_t^m)\right\|^2
}_{\text{from }I}
\;\underbrace{
  +\frac{\eta\,L^2}{2M}\sum_m\mathbb{E}\|x_t-x_t^m\|^2
}_{\text{from }I} \notag \\[6pt]
&\qquad\underbrace{
  +\frac{\eta\,\rho\,L^2}{2}\,\mathbb{E}\|z_t-x_t\|^2
}_{\text{from }II}
\;\underbrace{
  +\frac{\eta}{2\rho}\,\mathbb{E}\left\|\frac{1}{M}\sum_m
    \nabla f_m(x_t^m)\right\|^2
}_{\text{from }II} \notag \\[6pt]
&\qquad\underbrace{
  +\frac{\eta^2\,L}{2M}\,\sigma^2
}_{\text{from }III}
\;\underbrace{
  +\frac{\eta^2\,L}{2}\,\mathbb{E}\left\|\frac{1}{M}\sum_m
    \nabla f_m(x_t^m)\right\|^2
}_{\text{from }III}.
\label{eq:combined_raw}
\end{align}

Collect the three terms involving
$\mathbb{E}\left\|\frac{1}{M}\sum_m\nabla f_m(x_t^m)\right\|^2$:
their combined coefficient is
\begin{equation}
-\frac{\eta}{2} + \frac{\eta}{2\rho} + \frac{\eta^2\,L}{2}
= -\frac{\eta}{2}\left(1 - \frac{1}{\rho} - \eta\,L\right).
\label{eq:coeff_collect}
\end{equation}
Therefore~\eqref{eq:combined_raw} simplifies to:
\begin{align}
&\mathbb{E}[f(z_{t+1})] - \mathbb{E}[f(z_t)] \notag \\
&\quad\leq -\frac{\eta}{2}\,\mathbb{E}\|\nabla f(x_t)\|^2
  - \frac{\eta}{2}\left(1 - \frac{1}{\rho} - \eta\,L\right)
    \mathbb{E}\left\|\frac{1}{M}\sum_m\nabla f_m(x_t^m)\right\|^2
    \notag \\
&\qquad+ \frac{\eta\,L^2}{2M}\sum_m\mathbb{E}\|x_t - x_t^m\|^2
  + \frac{\eta\,\rho\,L^2}{2}\,\mathbb{E}\|z_t - x_t\|^2
  + \frac{\eta^2\,L}{2M}\,\sigma^2.
\label{eq:one_step}
\end{align}

%=============================================================================
% STEP 4
%=============================================================================
\paragraph{Step 4: Telescoping and applying the lemmas.}

Sum~\eqref{eq:one_step} from $t = 0$ to $T-1$ and divide by $T$:
\begin{equation}
\resizebox{0.98\linewidth}{!}{$
\begin{aligned}
\frac{\mathbb{E}[f(z_T)] - f(z_0)}{T}
&\leq
-\frac{\eta}{2T}\sum_{t=0}^{T-1}\mathbb{E}\|\nabla f(x_t)\|^2
-\frac{\eta}{2}\left(1-\frac{1}{\rho}-\eta L\right)
\frac{1}{T}\sum_{t=0}^{T-1}
\mathbb{E}\left\|\frac{1}{M}\sum_m \nabla f_m(x_t^m)\right\|^2  \\
&\quad
+\frac{\eta\rho L^2}{2}
\frac{1}{T}\sum_{t=0}^{T-1}\mathbb{E}\|z_t-x_t\|^2
+\frac{\eta L^2}{2M}
\frac{1}{TM}\sum_{t=0}^{T-1}\sum_m
\mathbb{E}\|x_t-x_t^m\|^2
+\frac{\eta^2L}{2M}\sigma^2 .
\end{aligned}
$}
\label{eq:telescoped}
\end{equation}

We now substitute the bounds from Lemma~\ref{lem:virt_drift} and
Lemma~\ref{lem:bound} (proved in Section~\ref{sec:key_lemmas}).

\medskip
\noindent\emph{Lemma~\ref{lem:virt_drift} gives:}
\begin{equation}
\frac{1}{T}\sum_{t=0}^{T-1}\mathbb{E}\|z_t - x_t\|^2
\leq \frac{\eta^2\,\beta^2}{(1-\beta)^2\,M}\,\sigma^2
  + \frac{\eta^2\,\beta^2}{(1-\beta)^2}\cdot
    \frac{1}{T}\sum_{t=0}^{T-1}\mathbb{E}\left\|\frac{1}{M}\sum_m
    \nabla f_m(x_t^m)\right\|^2.
\label{eq:lem2_applied}
\end{equation}

\noindent\emph{Lemma~\ref{lem:bound} gives (provided $24\eta^2\,L^2\,\psi \leq 1$):}
\begin{equation}
\frac{1}{TM}\sum_{t=0}^{T-1}\sum_m\mathbb{E}\|x_t - x_t^m\|^2
\leq 12\eta^2\,(B^2-1)\,\psi\cdot
  \frac{1}{T}\sum_{t=0}^{T-1}\mathbb{E}\|\nabla f(x_t)\|^2
  + 4\eta^2\,\psi\,(\sigma^2 + 3G^2),
\label{eq:lem3_applied}
\end{equation}
where the \emph{drift coefficient} $\psi$ is defined as
\begin{equation}
\psi \;\stackrel{\text{def}}{=}\;
\frac{4(1-p_x)}{p_x^2}\cdot
\frac{(1-\beta)(1-p_u)}{1-(1-p_u)\beta}.
\label{eq:psi_def}
\end{equation}

Substituting~\eqref{eq:lem2_applied} and~\eqref{eq:lem3_applied}
into~\eqref{eq:telescoped}, and collecting terms proportional to
$\frac{1}{T}\sum_t\mathbb{E}\|\nabla f(x_t)\|^2$ and
$\frac{1}{T}\sum_t\mathbb{E}\|\frac{1}{M}\sum_m\nabla f_m(x_t^m)\|^2$:

\medskip
\noindent\emph{Coefficient of $\frac{1}{T}\sum_t\mathbb{E}\|\nabla f(x_t)\|^2$:}
\begin{equation}
-\frac{\eta}{2} + \frac{\eta\,L^2}{2}\cdot 12\eta^2\,(B^2-1)\,\psi
= -\frac{\eta}{2}\bigl(1 - 12\eta^2\,L^2\,(B^2-1)\,\psi\bigr).
\label{eq:coeff_grad}
\end{equation}

\noindent\emph{Coefficient of $\frac{1}{T}\sum_t\mathbb{E}\|\frac{1}{M}\sum_m\nabla f_m(x_t^m)\|^2$:}
\begin{equation}
-\frac{\eta}{2}\left(1 - \frac{1}{\rho} - \eta\,L\right)
+ \frac{\eta\,\rho\,L^2}{2}\cdot\frac{\eta^2\,\beta^2}{(1-\beta)^2}
= -\frac{\eta}{2}\left(1 - \frac{1}{\rho} - \eta\,L
  - \frac{\eta^2\,\beta^2\,\rho\,L^2}{(1-\beta)^2}\right).
\label{eq:coeff_hetero}
\end{equation}

\noindent\emph{Residual terms (independent of the sums):}
\begin{equation}
\frac{\eta^2\,L}{2M}\,\sigma^2
+ \frac{\eta^3\,\rho\,L^2\,\beta^2}{2(1-\beta)^2\,M}\,\sigma^2
+ 2\eta^3\,L^2\,\psi\,(\sigma^2 + 3G^2).
\label{eq:residual}
\end{equation}

Therefore~\eqref{eq:telescoped} becomes:
\begin{align}
\frac{\mathbb{E}[f(z_T)] - f(z_0)}{T}
&\leq -\frac{\eta}{2}\bigl(1 - 12\eta^2\,L^2\,(B^2-1)\,\psi\bigr)
    \cdot\frac{1}{T}\sum_{t=0}^{T-1}\mathbb{E}\|\nabla f(x_t)\|^2
    \notag \\
&\quad- \frac{\eta}{2}\left(1 - \frac{1}{\rho} - \eta\,L
  - \frac{\eta^2\,\beta^2\,\rho\,L^2}{(1-\beta)^2}\right)
    \cdot\frac{1}{T}\sum_{t=0}^{T-1}\mathbb{E}\left\|\frac{1}{M}\sum_m
    \nabla f_m(x_t^m)\right\|^2 \notag \\
&\quad+ \frac{\eta^2\,L}{2M}\,\sigma^2
  + \frac{\eta^3\,\rho\,L^2\,\beta^2}{2(1-\beta)^2\,M}\,\sigma^2
  + 2\eta^3\,L^2\,\psi\,(\sigma^2 + 3G^2).
\label{eq:step4_assembled}
\end{align}

\medskip
\noindent\textbf{Choosing $\rho$ and $\eta$.}
Set $\rho = 2$. We need the coefficients of both squared-gradient terms to be
non-positive. This requires three simultaneous conditions on $\eta$:
\begin{alignat}{3}
&\text{(C1)}\quad& 12\eta^2\,L^2\,(B^2-1)\,\psi &\leq \tfrac{1}{2},
  &\qquad& \text{(first term $\geq -\tfrac{\eta}{4}\|\nabla f\|^2$)} \notag\\
&\text{(C2)}\quad& \eta\,L + \frac{2\eta^2\,\beta^2\,L^2}{(1-\beta)^2}
  &\leq \tfrac{1}{2},
  &\qquad& \text{(second term non-positive)} \notag\\
&\text{(C3)}\quad& 12\eta^2\,L^2\,\psi &\leq \tfrac{1}{2}.
  &\qquad& \text{(from Lemma~\ref{lem:bound} prerequisite)}
\label{eq:three_conditions}
\end{alignat}
One can verify that the step-size
\begin{equation}
\eta_0 \;\stackrel{\text{def}}{=}\;
\frac{1}{4L}\,\min\!\left(1 - \beta,\;
\frac{1}{6\sqrt{\psi\,\max(1,\,B^2-1)}}\right)
\label{eq:eta0}
\end{equation}
satisfies all three conditions. Indeed:
\begin{itemize}
\item (C1): $\eta_0 \leq \frac{1}{4L}\cdot\frac{1}{6\sqrt{\psi(B^2-1)}}$
  implies $\eta_0^2 \leq \frac{1}{144\,L^2\,\psi\,(B^2-1)}$,
  hence $12\eta_0^2\,L^2\,(B^2-1)\,\psi \leq 12/144 = 1/12 \leq 1/2$.
\item (C3): $\eta_0 \leq \frac{1}{4L}\cdot\frac{1}{6\sqrt{\psi}}$
  implies $12\eta_0^2\,L^2\,\psi \leq 12/(144) = 1/12 \leq 1/2$.
\item (C2): $\eta_0 \leq \frac{1-\beta}{4L}$ implies $\eta_0\,L \leq (1-\beta)/4$ and
  $\frac{2\eta_0^2\,\beta^2\,L^2}{(1-\beta)^2} \leq \frac{2\beta^2}{16} \leq 1/8$,
  so $\eta_0\,L + \frac{2\eta_0^2\,\beta^2\,L^2}{(1-\beta)^2} \leq 1/4 + 1/8 < 1/2$.
\end{itemize}

With any $\eta \leq \eta_0$, the second sum in~\eqref{eq:step4_assembled} is
non-positive and can be dropped. The first coefficient satisfies
$1 - 12\eta^2\,L^2\,(B^2-1)\,\psi \geq 1/2$, giving:
\begin{align}
\frac{\mathbb{E}[f(z_T)] - f(z_0)}{T}
&\leq -\frac{\eta}{4}\cdot\frac{1}{T}\sum_{t=0}^{T-1}
  \mathbb{E}\|\nabla f(x_t)\|^2 \notag \\
&\quad+ \frac{\eta^2\,L}{2M}\,\sigma^2
  + \frac{\eta^3\,\rho\,L^2\,\beta^2}{2(1-\beta)^2\,M}\,\sigma^2
  + 2\eta^3\,L^2\,\psi\,(\sigma^2 + 3G^2).
\label{eq:pre_final}
\end{align}

\medskip
\noindent\textbf{Rearranging for the convergence rate.}
Since $z_0 = x_0$ and $f(z_T) \geq f^*$, we have
$\mathbb{E}[f(z_T)] - f(z_0) \geq f^* - f(x_0)$, i.e.,
$f(x_0) - f^* \geq f(z_0) - \mathbb{E}[f(z_T)]$.
Rearranging~\eqref{eq:pre_final}:
\begin{align}
\frac{1}{T}\sum_{t=0}^{T-1}\mathbb{E}\|\nabla f(x_t)\|^2
&\leq \frac{4(f(x_0) - f^*)}{\eta\,T}
  + \frac{2\eta\,L}{M}\,\sigma^2
  + \frac{4\eta^2\,L^2\,\beta^2}{(1-\beta)^2\,M}\,\sigma^2
  + 8\eta^2\,L^2\,\psi\,(\sigma^2 + 3G^2).
\label{eq:rate_general}
\end{align}

\medskip
\noindent\textbf{Final rate with $\eta = \min(\eta_0, T^{-1/2})$.}
Setting $\eta = \min(\eta_0, 1/\sqrt{T})$, two cases arise.

\noindent\emph{Case 1: $\eta_0 \leq 1/\sqrt{T}$.}
Then $\eta = \eta_0$ and $1/(\eta\,T) = 1/(\eta_0\,T)$. Each
$\eta^2$ term is $\mathcal{O}(\eta_0^2) = \mathcal{O}(1/T)$ (since
$\eta_0 \leq 1/\sqrt{T}$ in this case).

\noindent\emph{Case 2: $\eta_0 > 1/\sqrt{T}$.}
Then $\eta = 1/\sqrt{T}$ and $1/(\eta\,T) = 1/\sqrt{T}$.
Each $\eta^2$ term is $1/T$.

In either case, collecting the $\mathcal{O}(1/\sqrt{T})$ and
$\mathcal{O}(1/T)$ contributions:
\begin{align}
\frac{1}{T}\sum_{t=0}^{T-1}\mathbb{E}\|\nabla f(x_t)\|^2
&\leq \frac{4}{\sqrt{T}}\left(f(x_0) - f^*
  + \frac{L\,\sigma^2}{2M}\right)
  + \frac{4(f(x_0) - f^*)}{\eta_0\,T} \notag \\
&\quad+ \frac{4L^2\,\beta^2\,\sigma^2}{(1-\beta)^2\,M\,T}
  + \frac{8L^2\,\psi\,(\sigma^2 + 3G^2)}{T} \notag \\[6pt]
&= \frac{4}{\sqrt{T}}\left(f(x_0) - f^*
  + \frac{L\,\sigma^2}{2M}\right)
  + \mathcal{O}\!\left(\frac{1 + \psi}{T}\right).
\label{eq:final_rate}
\end{align}
This completes the proof of Theorem~\ref{thm:sgdm}.

\paragraph{Remark (SP invariance of the convergence rate).}
The drift coefficient $\psi$ defined in~\eqref{eq:psi_def} depends
\emph{only} on the synchronization probabilities $p_x = 1/K_x$,
$p_u = 1/K_u$ and the momentum decay $\beta$. It is
\emph{independent} of the SP degree $P$. This is because:
\begin{enumerate}
\item The gradient $g_t^m$ entering the optimizer is already
  sequence-reduced by the backward All-to-All (Section~\ref{sec:appendix_c}).
  By the SP gradient identity~\eqref{eq:sp_gradient_identity},
  it satisfies $\mathbb{E}[g_t^m] = \nabla f_m(x_t^m)$ and
  $\mathbb{E}\|g_t^m - \nabla f_m(x_t^m)\|^2 \leq \sigma^2$
  with the \emph{same} constants as in the non-SP case---this
  is precisely the variance invariance~\eqref{eq:sp_variance_invariance}.
\item The synchronization operator $\mathbb{E}_m[\cdot]$ averages
  over the $M$ data-parallel replicas, not over the $P$ SP ranks
  within a replica. The SP ranks collaborate to produce a single
  gradient per replica; the desynced averaging then operates on
  these replica-level gradients.
  The collective orthogonality constraint (Phase~3 fence in
  Algorithm~\ref{alg:desloc_adam}) guarantees that this separation
  holds at every step: no DP AllReduce fires until all SP A2A
  handles have retired.
\item Consequently, every term in the proof chain
  \eqref{eq:I_final}--\eqref{eq:final_rate}
  is unaffected by $P$: the variance bound $\sigma^2/M$, the
  drift bound $\psi$, and the step-size restriction $\eta_0$ are
  all $P$-independent.
\end{enumerate}
The SP decomposition affects only the \emph{wall-clock cost per step}
(Section~\ref{sec:wallclock_bandwidth}), not the optimization trajectory
or convergence rate.

\paragraph{Corollary (SP-tightened constant for attention parameters).}
While the convergence \emph{order} is $P$-independent, the
\emph{constant} inside the bound can be refined for the $d_{\text{attn}}$
attention-layer coordinates.
By the SP variance reduction~\eqref{eq:sp_var_reduction} and the
SP-tightened clipping bound $\rho_{\text{eff}} = \rho/\sqrt{P}$
(Algorithm~\ref{alg:desloc_adam}, Phase~2), the drift
bounds~\eqref{eq:drift_u_sp}--\eqref{eq:drift_v_sp} yield a
refined per-coordinate contribution:
\begin{equation}
\psi_{\text{attn}}
\;\leq\; \frac{\psi}{P},
\label{eq:psi_attn_sp}
\end{equation}
where $\psi_{\text{attn}}$ is the drift coefficient restricted to the
attention-layer coordinates. Since these coordinates constitute a fraction
$d_{\text{attn}}/d$ of the total, the overall drift decomposes as
$\psi_{\text{eff}} = (1 - d_{\text{attn}}/d)\,\psi + (d_{\text{attn}}/d)\,\psi/P$,
yielding a tighter constant for the final rate~\eqref{eq:final_rate}
without changing its order.
At $P = 2$ (our default SP degree), attention-layer drift is halved;
at $P = 4$, it is quartered.

\paragraph{Remark (tiered gradient buckets and $\psi$).}
In the FAUST implementation, the uniform periods $(K_x, K_u)$ in the
SGDM analysis correspond to the \emph{fastest} tier's periods. The slower
tiers ($u$-tier at $K_u = 3K_x$, $v$-tier at $K_v = 6K_x$) contribute
additional higher-order terms that refine the constant inside the
$\mathcal{O}((1+\psi)/T)$ term but do not change its order. The key
structural insight is that $\psi = \mathcal{O}(1/p_x^2)$: parameter
synchronization frequency dominates the bound, while momentum
synchronization enters only linearly.

\paragraph{Remark (role of the collective fence in the proof).}
The entire proof chain assumes that $g_t^m = \nabla F_m(x_t^m;\, \xi_t^m)$
is an unbiased estimator of $\nabla f_m(x_t^m)$. This holds only if the
SP backward All-to-All has fully completed before $g_t^m$ is used.
In the absence of the Phase~3 fence, a race condition arises:
the DP AllReduce $\mathbb{E}_m[\cdot]$ could average
\emph{partially-reduced} gradients $\hat{g}_t^m$ where some SP ranks have
completed their All-to-All but others have not.
In this scenario, $\hat{g}_t^m \neq g_t^m$ and the unbiasedness
assumption $\mathbb{E}[\hat{g}_t^m] = \nabla f_m(x_t^m)$ is violated,
invalidating Step~3a (the variance decomposition~\eqref{eq:var_decomp})
and propagating through the entire bound.
The fence guarantees that this never occurs: it is not merely an
engineering convenience but a \emph{necessary condition for correctness}
of the convergence proof.

%=============================================================================
% D.1 Extension to Adam optimizer
%=============================================================================
\subsection{Extension to Adam optimizer}
\label{sec:adam_extension}
\label{sec:high_prob}

We now extend the SGDM analysis of the previous section to the Adam optimizer,
which maintains \emph{two} optimizer states: a first moment $u_t^m$ and a second
moment $v_t^m$, each with its own decay factor ($\beta_1$ and $\beta_2$
respectively) and its own synchronization probability ($p_u = 1/K_u$ and
$p_v = 1/K_v$). The key new difficulty is that the parameter update involves a
\emph{preconditioned} step $d_t^m = \eta\,\Gamma_t^m\,u_t^m$, where the
preconditioning matrix $\Gamma_t^m$ depends on the local second moment and
therefore varies across replicas. This means the virtual-iterate trick of
Section~\ref{sec:sgdm_proof} no longer yields a clean $z_{t+1} - z_t = -\eta\,g_t$
identity; instead, additional error terms $U_t$ and $V_t$ appear, which must be
controlled.

The complete algorithm is given in Algorithm~\ref{alg:desloc_adam_prob} below.
Instead of bounded heterogeneity (Assumption~\ref{asm:heterogeneity}), we use a
stronger condition:

\begin{assumption}[Bounded gradient]
\label{ass:bounded_grad}
For any iterate $t \geq 0$ and replica $m$, the local stochastic gradient is
bounded:
\begin{equation}
\|g_t^m\|_2 \;\leq\; G.
\label{eq:bounded_grad}
\end{equation}
\end{assumption}

This condition provides uniform upper bounds on gradients and momenta and is
standard in the analysis of adaptive
optimizers~\citep{chen2019convergenceclassadamtypealgorithms,defossez2022simple,zhang2023adamconvergemodificationupdate}.
As in the SGDM analysis, the SP degree $P$ does not appear: the gradient
$g_t^m$ is already sequence-reduced and satisfies Assumption~\ref{ass:bounded_grad}
with the same constant $G$.

%-----------------------------------------------------------------------------
% Algorithm 5
%-----------------------------------------------------------------------------
\begin{algorithm}[ht]
\caption{FAUST-Adam (probabilistic synchronization formulation)}
\label{alg:desloc_adam_prob}
\begin{algorithmic}[1]
\REQUIRE \textbf{Model tensors}
\STATE $x_0 \in \mathbb{R}^d$ \hfill initial parameter vector
\STATE $u_{-1}, v_{-1} \in \mathbb{R}^d$ \hfill seeds for first and second moments, initialised to $\mathbf{0}$
\REQUIRE \textbf{Hyper-parameters}
\STATE $\{\eta_t\}_{t=0}^{T-1} \subset \mathbb{R}_{>0}$ \hfill $K_x$-aligned WSD step-size schedule
\STATE $\beta_1, \beta_2 \in [0,1)$ \hfill Adam decay factors
\STATE $\lambda \in \mathbb{R}_{\geq 0}$ \hfill $\ell_2$ stability term
\STATE $T \in \mathbb{N}_+$ \hfill total optimisation steps
\STATE $M \in \mathbb{N}_+$ \hfill number of replicas
\STATE $p_x = \frac{1}{K_x},\, p_u = \frac{1}{K_u},\, p_v = \frac{1}{K_v} \in [0,1]$ \hfill synchronization probabilities for parameters and momenta
\ENSURE $x_T$, $u_{T-1}$, $v_{T-1}$
\STATE \textbf{for each replica} $m$: $x_0^m = x_0$, $u_{-1}^m = v_{-1}^m = 0$ \hfill local init
\FOR{$t = 0, \ldots, T-1$}
\FOR{all replicas $m = 0, \ldots, M-1$ \textbf{in parallel}}
\STATE $g_t^m \leftarrow \nabla F_m(x_t^m; \xi_t^m)$ \hfill stochastic gradient (SP-reduced when $P > 1$)
\STATE $u_t^m \leftarrow \begin{cases} \mathbb{E}_m[\beta_1 u_{t-1}^m + (1-\beta_1) g_t^m], & \text{w.p.\ } p_u \\ \beta_1 u_{t-1}^m + (1-\beta_1) g_t^m, & \text{w.p.\ } 1 - p_u \end{cases}$
\STATE $v_t^m \leftarrow \begin{cases} \mathbb{E}_m[\beta_2 v_{t-1}^m + (1-\beta_2)(g_t^m \odot g_t^m)], & \text{w.p.\ } p_v \\ \beta_2 v_{t-1}^m + (1-\beta_2)(g_t^m \odot g_t^m), & \text{w.p.\ } 1 - p_v \end{cases}$
\STATE $\tilde{v}_t^m \leftarrow \max(v_t^m,\, \tilde{v}_{t-1}^m)$ \hfill AMSGrad normalization ($\tilde{v}_{-1} = v_{-1}$)
\STATE $d_t^m \leftarrow \frac{\eta_t}{\sqrt{\tilde{v}_t^m + \lambda^2}} \odot u_t^m$ \hfill preconditioned update
\STATE $x_{t+1}^m \leftarrow \begin{cases} \mathbb{E}_m[x_t^m - d_t^m], & \text{w.p.\ } p_x \\ x_t^m - d_t^m, & \text{w.p.\ } 1 - p_x \end{cases}$
\ENDFOR
\ENDFOR
\end{algorithmic}
\end{algorithm}

%-----------------------------------------------------------------------------
% Notation
%-----------------------------------------------------------------------------
\paragraph{Notation for the Adam analysis.}
We define the following quantities that will be used throughout this subsection.

\medskip
\noindent\textbf{Preconditioning matrix.}
For each replica $m$ at step $t$, define the diagonal preconditioning matrix
\begin{equation}
\Gamma_t^m \;\stackrel{\text{def}}{=}\;
\mathrm{diag}^{-1/2}\!\bigl(\tilde{v}_t^m + \lambda^2\bigr),
\qquad\text{i.e.,}\quad
[\Gamma_t^m]_{ii} = \frac{1}{\sqrt{[\tilde{v}_t^m]_i + \lambda^2}}.
\label{eq:adam_gamma_def}
\end{equation}
Since $\tilde{v}_t^m \geq 0$ coordinate-wise and $\lambda > 0$, we have
$0 < [\Gamma_t^m]_{ii} \leq 1/\lambda$ for all $i$, hence
\begin{equation}
\|\Gamma_t^m\| \;\leq\; \frac{1}{\lambda},
\label{eq:gamma_upper}
\end{equation}
where $\|\cdot\|$ denotes the spectral norm (which equals the largest diagonal
entry for a positive diagonal matrix).

\medskip
\noindent\textbf{Cross-replica averages.}
As in the SGDM analysis, define
\begin{align}
x_t &= \mathbb{E}_m[x_t^m], &
u_t &= \mathbb{E}_m[u_t^m], &
v_t &= \mathbb{E}_m[v_t^m], \notag \\
\tilde{v}_t &= \mathbb{E}_m[\tilde{v}_t^m], &
\Gamma_t &= \mathbb{E}_m[\Gamma_t^m], &
g_t &= \mathbb{E}_m[g_t^m].
\label{eq:adam_averages}
\end{align}

\medskip
\noindent\textbf{Condition number.}
By Assumption~\ref{ass:bounded_grad}, $\|g_t^m\| \leq G$, so
$[v_t^m]_i \leq G^2$ and $[\tilde{v}_t^m]_i \leq G^2$. Hence the
minimum eigenvalue of $\Gamma_t$ (averaging over replicas) satisfies
\begin{equation}
\|\Gamma_t\|_{\min}
= \min_i [\Gamma_t]_{ii}
\geq \frac{1}{\sqrt{G^2 + \lambda^2}}
\;\stackrel{\text{def}}{=}\; \frac{1}{C_0}.
\label{eq:C0_def}
\end{equation}
The ratio $C_0/\lambda = \sqrt{G^2 + \lambda^2}/\lambda$ plays the role of a
condition number for the preconditioner.

\medskip
\noindent\textbf{Pseudo-gradient.}
Define the ``de-biased'' local gradient
\begin{equation}
\widetilde{g}_t^m \;\stackrel{\text{def}}{=}\;
\frac{u_t^m - \beta_1\,u_{t-1}^m}{1 - \beta_1}.
\label{eq:gtilde_def}
\end{equation}
By the momentum recursion $u_t^m = \beta_1\,u_{t-1}^m + (1-\beta_1)\,g_t^m$
(in the non-sync case), we have $\widetilde{g}_t^m = g_t^m$ when no sync
occurs. In general,
\begin{equation}
\mathbb{E}_m[\widetilde{g}_t^m]
= \frac{u_t - \beta_1\,u_{t-1}}{1 - \beta_1}
= g_t,
\label{eq:gtilde_avg}
\end{equation}
where the last step uses the averaged momentum recursion
$u_t = \beta_1\,u_{t-1} + (1-\beta_1)\,g_t$ (derived identically to
equation~\eqref{eq:avg_u_recursion} in the SGDM proof).

%-----------------------------------------------------------------------------
% Step 1: Averaged dynamics
%-----------------------------------------------------------------------------
\paragraph{Step 1: Averaged dynamics and virtual iterates.}

The averaged first momentum follows the same centralized recursion as before:
\begin{equation}
u_t = \beta_1\,u_{t-1} + (1-\beta_1)\,g_t.
\label{eq:adam_avg_u}
\end{equation}
The averaged parameter update, however, now involves the preconditioner:
\begin{equation}
x_{t+1} = x_t - d_t,
\qquad\text{where}\quad
d_t \;\stackrel{\text{def}}{=}\;
\eta\,\mathbb{E}_m[\Gamma_t^m\,u_t^m].
\label{eq:adam_avg_x}
\end{equation}
Note that $d_t \neq \eta\,\Gamma_t\,u_t$ in general, because $\Gamma_t^m$ and
$u_t^m$ are correlated (both depend on the local gradient history of
replica $m$). This is the fundamental complication relative to SGDM.

\medskip
\noindent\textbf{Virtual iterates.}
Define $z_t$ exactly as in~\eqref{eq:virtual}:
\begin{equation}
z_t = \frac{1}{1-\beta_1}\,x_t - \frac{\beta_1}{1-\beta_1}\,x_{t-1},
\qquad z_0 = x_0.
\label{eq:adam_virtual}
\end{equation}

\noindent\textbf{Computing $z_{t+1} - z_t$ step by step.}
From the definition~\eqref{eq:adam_virtual}:
\begin{align}
z_{t+1} - z_t
&= \frac{1}{1-\beta_1}\,(x_{t+1} - x_t)
   - \frac{\beta_1}{1-\beta_1}\,(x_t - x_{t-1}).
\label{eq:adam_zdiff_1}
\end{align}
Using $x_{t+1} - x_t = -d_t = -\eta\,\mathbb{E}_m[\Gamma_t^m\,u_t^m]$ and
$x_t - x_{t-1} = -d_{t-1} = -\eta\,\mathbb{E}_m[\Gamma_{t-1}^m\,u_{t-1}^m]$:
\begin{align}
z_{t+1} - z_t
&= -\frac{\eta}{1-\beta_1}\,\mathbb{E}_m[\Gamma_t^m\,u_t^m]
   + \frac{\eta\,\beta_1}{1-\beta_1}\,\mathbb{E}_m[\Gamma_{t-1}^m\,u_{t-1}^m].
\label{eq:adam_zdiff_2}
\end{align}

We now add and subtract
$\frac{\eta\,\beta_1}{1-\beta_1}\,\mathbb{E}_m[\Gamma_t^m\,u_{t-1}^m]$
inside~\eqref{eq:adam_zdiff_2}. This is the key algebraic trick: it allows us
to isolate the ``ideal'' gradient term from the preconditioner-mismatch error.

\begin{align}
z_{t+1} - z_t
&= -\frac{\eta}{1-\beta_1}\,\mathbb{E}_m[\Gamma_t^m\,u_t^m]
   + \frac{\eta\,\beta_1}{1-\beta_1}\,\mathbb{E}_m[\Gamma_t^m\,u_{t-1}^m]
   \notag \\
&\quad+ \frac{\eta\,\beta_1}{1-\beta_1}\,
   \mathbb{E}_m[(\Gamma_{t-1}^m - \Gamma_t^m)\,u_{t-1}^m].
\label{eq:adam_zdiff_3}
\end{align}

Factor the first two terms:
\begin{align}
-\frac{\eta}{1-\beta_1}\,\mathbb{E}_m[\Gamma_t^m\,u_t^m]
+ \frac{\eta\,\beta_1}{1-\beta_1}\,\mathbb{E}_m[\Gamma_t^m\,u_{t-1}^m]
&= -\frac{\eta}{1-\beta_1}\,
   \mathbb{E}_m\bigl[\Gamma_t^m\,(u_t^m - \beta_1\,u_{t-1}^m)\bigr].
\label{eq:adam_zdiff_4}
\end{align}

By the definition~\eqref{eq:gtilde_def},
$u_t^m - \beta_1\,u_{t-1}^m = (1-\beta_1)\,\widetilde{g}_t^m$, so
\begin{equation}
-\frac{\eta}{1-\beta_1}\,
\mathbb{E}_m\bigl[\Gamma_t^m\,(1-\beta_1)\,\widetilde{g}_t^m\bigr]
= -\eta\,\mathbb{E}_m[\Gamma_t^m\,\widetilde{g}_t^m].
\label{eq:adam_zdiff_5}
\end{equation}

Substituting~\eqref{eq:adam_zdiff_4} and~\eqref{eq:adam_zdiff_5}
into~\eqref{eq:adam_zdiff_3}:
\begin{align}
z_{t+1} - z_t
&= -\eta\,\mathbb{E}_m[\Gamma_t^m\,\widetilde{g}_t^m]
   + \frac{\eta\,\beta_1}{1-\beta_1}\,
   \mathbb{E}_m[(\Gamma_{t-1}^m - \Gamma_t^m)\,u_{t-1}^m].
\label{eq:adam_zdiff_6}
\end{align}

Now add and subtract $\eta\,\mathbb{E}_m[\Gamma_t^m\,g_t]$ inside the first
term. Since $\mathbb{E}_m[\widetilde{g}_t^m] = g_t$
by~\eqref{eq:gtilde_avg}:
\begin{align}
-\eta\,\mathbb{E}_m[\Gamma_t^m\,\widetilde{g}_t^m]
&= -\eta\,\mathbb{E}_m[\Gamma_t^m\,g_t]
   + \eta\,\mathbb{E}_m[\Gamma_t^m\,(g_t - \widetilde{g}_t^m)] \notag \\
&= -\eta\,\Gamma_t\,g_t
   + \eta\,\mathbb{E}_m[\Gamma_t^m\,(g_t - \widetilde{g}_t^m)],
\label{eq:adam_zdiff_7}
\end{align}
where the last step uses the linearity of $\mathbb{E}_m$ and the fact that
$g_t$ is constant across the replica average, so
$\mathbb{E}_m[\Gamma_t^m\,g_t] = \mathbb{E}_m[\Gamma_t^m]\,g_t = \Gamma_t\,g_t$.

\medskip
\noindent\textbf{Final decomposition.}
Combining~\eqref{eq:adam_zdiff_6} and~\eqref{eq:adam_zdiff_7}, and defining
the two error terms:
\begin{align}
z_{t+1} - z_t
&= -\eta\,\Gamma_t\,g_t
   + \eta\,\underbrace{\mathbb{E}_m[\Gamma_t^m\,(g_t - \widetilde{g}_t^m)]
   }_{\stackrel{\text{def}}{=}\;U_t}
   + \eta\,\underbrace{\frac{\beta_1}{1-\beta_1}\,
   \mathbb{E}_m[(\Gamma_{t-1}^m - \Gamma_t^m)\,u_{t-1}^m]
   }_{\stackrel{\text{def}}{=}\;V_t}.
\label{eq:adam_zdiff_final}
\end{align}

Compared to the SGDM result $z_{t+1} - z_t = -\eta\,g_t$
(equation~\eqref{eq:virt_gradient}), two additional terms appear:

\begin{itemize}
\item $U_t$ captures the \emph{momentum drift}: the difference between the
  global averaged gradient $g_t$ and the local pseudo-gradient
  $\widetilde{g}_t^m$, weighted by the local preconditioner $\Gamma_t^m$.
  This vanishes when all replicas are synchronized ($u_t^m = u_t$ for all $m$).

\item $V_t$ captures the \emph{preconditioner drift}: the change in the
  preconditioning matrix between consecutive steps, weighted by the previous
  momentum. This vanishes when the second moment is stationary
  ($\Gamma_{t-1}^m = \Gamma_t^m$), which is approximately true when $\beta_2$
  is close to $1$ (large half-life).
\end{itemize}

%-----------------------------------------------------------------------------
% Step 2: Smoothness descent
%-----------------------------------------------------------------------------
\paragraph{Step 2: Smoothness descent on virtual iterates.}

As before, $L$-smoothness gives
\begin{equation}
f(z_{t+1}) \leq f(z_t) + \langle \nabla f(z_t),\, z_{t+1} - z_t \rangle
+ \frac{L}{2}\,\|z_{t+1} - z_t\|^2.
\label{eq:adam_smooth}
\end{equation}
Substituting~\eqref{eq:adam_zdiff_final} into~\eqref{eq:adam_smooth} and
expanding the inner product using
$\nabla f(z_t) = \nabla f(x_t) + (\nabla f(z_t) - \nabla f(x_t))$:
\begin{align}
f(z_{t+1}) - f(z_t)
&\leq -\eta\,\langle \nabla f(z_t),\, \Gamma_t\,g_t \rangle
   + \eta\,\langle \nabla f(z_t),\, U_t \rangle
   + \eta\,\langle \nabla f(z_t),\, V_t \rangle
   + \frac{L}{2}\,\|z_{t+1} - z_t\|^2.
\label{eq:adam_smooth_expanded}
\end{align}
We split the first inner product by adding and subtracting $\nabla f(x_t)$:
\begin{align}
-\eta\,\langle \nabla f(z_t),\, \Gamma_t\,g_t \rangle
&= -\eta\,\langle \nabla f(x_t),\, \Gamma_t\,g_t \rangle
   - \eta\,\langle \nabla f(z_t) - \nabla f(x_t),\, \Gamma_t\,g_t \rangle.
\label{eq:adam_split_I_V}
\end{align}
Substituting into~\eqref{eq:adam_smooth_expanded} and labelling each group:
\begin{align}
f(z_{t+1}) - f(z_t)
&\leq \underbrace{-\eta\,\langle \nabla f(x_t),\,
  \Gamma_t\,g_t \rangle}_{I}
  + \underbrace{\eta\,\langle \nabla f(z_t),\, U_t \rangle}_{II}
  + \underbrace{\eta\,\langle \nabla f(z_t),\, V_t \rangle}_{III}
  \notag \\
&\quad+ \underbrace{\frac{L}{2}\,\|z_{t+1} - z_t\|^2}_{IV}
  + \underbrace{\eta\,\langle \nabla f(x_t) - \nabla f(z_t),\,
  \Gamma_t\,g_t \rangle}_{V}.
\label{eq:adam_five_terms}
\end{align}
Note that this is a five-term decomposition (compared to three in SGDM),
reflecting the additional complexity of the adaptive preconditioner.
We bound each term below.

%-----------------------------------------------------------------------------
% Step 3a: Term I
%-----------------------------------------------------------------------------
\paragraph{Step 3a: Bounding term $I$.}

We split $\Gamma_t = \Gamma_{t-1} + (\Gamma_t - \Gamma_{t-1})$:
\begin{align}
I &= -\eta\,\langle \nabla f(x_t),\, \Gamma_t\,g_t \rangle \notag \\
  &= -\eta\,\langle \nabla f(x_t),\, \Gamma_{t-1}\,g_t \rangle
     -\eta\,\langle \nabla f(x_t),\,
     (\Gamma_t - \Gamma_{t-1})\,g_t \rangle.
\label{eq:adam_I_split}
\end{align}

For the second term, by Cauchy--Schwarz and
$\|\nabla f(x_t)\| \leq G$, $\|g_t\| \leq G$:
\begin{equation}
|\langle \nabla f(x_t),\, (\Gamma_t - \Gamma_{t-1})\,g_t \rangle|
\leq \|\nabla f(x_t)\| \cdot \|\Gamma_t - \Gamma_{t-1}\| \cdot \|g_t\|
\leq G^2\,\|\Gamma_t - \Gamma_{t-1}\|.
\label{eq:adam_I_perturb}
\end{equation}

For the first term, taking expectation and using unbiasedness
$\mathbb{E}[g_t] = \frac{1}{M}\sum_m \nabla f_m(x_t^m)$, then applying the
weighted polarization identity
$\langle a,b \rangle_A = \frac{1}{2}\|a\|_A^2 + \frac{1}{2}\|b\|_A^2
- \frac{1}{2}\|a-b\|_A^2$ with $A = \Gamma_{t-1}$:
\begin{align}
\mathbb{E}\bigl[-\eta\,\langle \nabla f(x_t),\,
\Gamma_{t-1}\,g_t \rangle\bigr]
&= -\eta\,\mathbb{E}\!\left[\left\langle \nabla f(x_t),\,
   \frac{1}{M}\sum_m \nabla f_m(x_t^m)
   \right\rangle_{\!\Gamma_{t-1}}\right] \notag \\
&= -\frac{\eta}{2}\,\mathbb{E}\bigl[\|\nabla f(x_t)\|_{\Gamma_{t-1}}^2\bigr]
   - \frac{\eta}{2}\,\mathbb{E}\!\left[\left\|
   \frac{1}{M}\sum_m \nabla f_m(x_t^m)
   \right\|_{\!\Gamma_{t-1}}^2\right] \notag \\
&\quad+ \frac{\eta}{2}\,\mathbb{E}\!\left[\left\|
   \nabla f(x_t) - \frac{1}{M}\sum_m \nabla f_m(x_t^m)
   \right\|_{\!\Gamma_{t-1}}^2\right].
\label{eq:adam_I_polar}
\end{align}

Now bound each factor using the spectral bounds on $\Gamma_{t-1}$:
\begin{align}
\|\nabla f(x_t)\|_{\Gamma_{t-1}}^2
&\geq \|\Gamma_{t-1}\|_{\min}\,\|\nabla f(x_t)\|^2
 \geq \frac{1}{C_0}\,\|\nabla f(x_t)\|^2,
\label{eq:adam_I_lower} \\[2pt]
\left\|\nabla f(x_t) - \frac{1}{M}\sum_m \nabla f_m(x_t^m)
\right\|_{\Gamma_{t-1}}^2
&\leq \|\Gamma_{t-1}\|_{\max}
\left\|\nabla f(x_t)-\frac{1}{M}\sum_m \nabla f_m(x_t^m)\right\|^2 \notag\\
&\leq \frac{1}{\lambda}\,(\cdots).
\label{eq:adam_I_upper}
\end{align}

For the last factor, by Jensen's inequality and $L$-smoothness (same chain as
\eqref{eq:jensen_grad}--\eqref{eq:smooth_sq} in the SGDM proof):
\begin{equation}
\left\|\nabla f(x_t) - \frac{1}{M}\sum_m\nabla f_m(x_t^m)\right\|^2
\leq \frac{1}{M}\sum_m \|\nabla f_m(x_t) - \nabla f_m(x_t^m)\|^2
\leq \frac{L^2}{M}\sum_m \|x_t - x_t^m\|^2.
\label{eq:adam_I_jensen}
\end{equation}

Combining~\eqref{eq:adam_I_split}--\eqref{eq:adam_I_jensen}:
\begin{align}
\mathbb{E}[I]
&\leq -\frac{\eta}{2C_0}\,\mathbb{E}\|\nabla f(x_t)\|^2
  - \frac{\eta}{2}\,\mathbb{E}\!\left[\left\|
  \frac{1}{M}\sum_m \nabla f_m(x_t^m)
  \right\|_{\!\Gamma_{t-1}}^2\right] \notag \\
&\quad+ \frac{\eta\,L^2}{2\lambda\,M}\sum_m
  \mathbb{E}\|x_t - x_t^m\|^2
  + \eta\,G^2\,\mathbb{E}\|\Gamma_{t-1} - \Gamma_t\|.
\label{eq:adam_I_final}
\end{align}

%-----------------------------------------------------------------------------
% Step 3b: Term II
%-----------------------------------------------------------------------------
\paragraph{Step 3b: Bounding term $II$.}

By Cauchy--Schwarz:
\begin{align}
II &= \eta\,\langle \nabla f(z_t),\, U_t \rangle
   \leq \eta\,\|\nabla f(z_t)\|\,\|U_t\|.
\label{eq:adam_II_cs}
\end{align}
By Assumption~\ref{ass:bounded_grad}, $\|\nabla f(z_t)\| \leq G$. Now bound
$\|U_t\|$. Recall $U_t = \mathbb{E}_m[\Gamma_t^m\,(g_t - \widetilde{g}_t^m)]$.
By Jensen's inequality applied to $\|\cdot\|$:
\begin{align}
\|U_t\| &= \left\|\mathbb{E}_m[\Gamma_t^m\,(g_t - \widetilde{g}_t^m)]
   \right\| \notag \\
&\leq \mathbb{E}_m\bigl[\|\Gamma_t^m\,(g_t - \widetilde{g}_t^m)\|\bigr]
   \notag \\
&= \frac{1}{M}\sum_m \|\Gamma_t^m\,(g_t - \widetilde{g}_t^m)\|.
\label{eq:adam_II_Ut_1}
\end{align}
Since $\Gamma_t^m$ is diagonal with $\|\Gamma_t^m\| \leq 1/\lambda$
by~\eqref{eq:gamma_upper}:
\begin{equation}
\|\Gamma_t^m\,(g_t - \widetilde{g}_t^m)\|
\leq \|\Gamma_t^m\|\,\|g_t - \widetilde{g}_t^m\|
\leq \frac{1}{\lambda}\,\|g_t - \widetilde{g}_t^m\|.
\label{eq:adam_II_Ut_2}
\end{equation}
Combining~\eqref{eq:adam_II_cs}--\eqref{eq:adam_II_Ut_2}:
\begin{equation}
\mathbb{E}[II]
\leq \frac{\eta\,G}{\lambda\,M}\sum_m
\mathbb{E}\|g_t - \widetilde{g}_t^m\|.
\label{eq:adam_II_final}
\end{equation}

%-----------------------------------------------------------------------------
% Step 3c: Term III
%-----------------------------------------------------------------------------
\paragraph{Step 3c: Bounding term $III$.}

Similarly, $III = \eta\,\langle \nabla f(z_t),\, V_t \rangle \leq
\eta\,G\,\|V_t\|$. Recall
$V_t = \frac{\beta_1}{1-\beta_1}\,
\mathbb{E}_m[(\Gamma_{t-1}^m - \Gamma_t^m)\,u_{t-1}^m]$. By Jensen:
\begin{align}
\|V_t\| &\leq \frac{\beta_1}{1-\beta_1}\,
\frac{1}{M}\sum_m \|(\Gamma_{t-1}^m - \Gamma_t^m)\,u_{t-1}^m\| \notag \\
&\leq \frac{\beta_1}{1-\beta_1}\,
\frac{1}{M}\sum_m \|\Gamma_{t-1}^m - \Gamma_t^m\|\,\|u_{t-1}^m\|.
\label{eq:adam_III_Vt}
\end{align}
Since $u_{t-1}^m$ is a convex combination of gradients bounded by $G$, we have
$\|u_{t-1}^m\| \leq G$. Therefore:
\begin{equation}
\mathbb{E}[III]
\leq \frac{\eta\,\beta_1\,G^2}{(1-\beta_1)\,M}\sum_m
\mathbb{E}\|\Gamma_{t-1}^m - \Gamma_t^m\|.
\label{eq:adam_III_final}
\end{equation}

%-----------------------------------------------------------------------------
% Step 3d: Term IV
%-----------------------------------------------------------------------------
\paragraph{Step 3d: Bounding term $IV$.}

From~\eqref{eq:adam_zdiff_final},
$z_{t+1} - z_t = -\eta\,\Gamma_t\,g_t + \eta\,U_t + \eta\,V_t$, so
\begin{equation}
\|z_{t+1} - z_t\|^2
= \eta^2\,\|\Gamma_t\,g_t - U_t - V_t\|^2.
\label{eq:adam_IV_start}
\end{equation}
By the triangle inequality for norms and the elementary bound
$\|a + b + c\|^2 \leq 3(\|a\|^2 + \|b\|^2 + \|c\|^2)$:
\begin{align}
\|\Gamma_t\,g_t - U_t - V_t\|^2
&\leq 3\,\|\Gamma_t\,g_t\|^2
   + 3\,\|U_t\|^2
   + 3\,\|V_t\|^2.
\label{eq:adam_IV_triangle}
\end{align}
We bound each:
\begin{alignat}{2}
\|\Gamma_t\,g_t\|^2
&\leq \frac{1}{\lambda^2}\,\|g_t\|^2
&&\leq \frac{G^2}{\lambda^2},
\label{eq:adam_IV_a} \\[4pt]
\|U_t\|^2
&\leq \frac{1}{\lambda^2\,M}\sum_m\|g_t - \widetilde{g}_t^m\|^2,
\label{eq:adam_IV_b} \\[4pt]
\|V_t\|^2
&\leq \frac{\beta_1^2\,G^2}{(1-\beta_1)^2\,M}\sum_m
\|\Gamma_{t-1}^m - \Gamma_t^m\|^2.
\label{eq:adam_IV_c}
\end{alignat}
Here~\eqref{eq:adam_IV_b} uses Jensen on the average defining $U_t$ plus
$\|\Gamma_t^m\| \leq 1/\lambda$; \eqref{eq:adam_IV_c} uses~\eqref{eq:adam_III_Vt}
squared plus $\|u_{t-1}^m\| \leq G$.

Combining:
\begin{align}
\mathbb{E}[IV]
&\leq \frac{3\eta^2\,L\,G^2}{2\lambda^2}
  + \frac{3\eta^2\,L}{2\lambda^2\,M}\sum_m
  \mathbb{E}\|g_t - \widetilde{g}_t^m\|^2 \notag \\
&\quad+ \frac{3\eta^2\,\beta_1^2\,L\,G^2}{2(1-\beta_1)^2\,M}\sum_m
  \mathbb{E}\|\Gamma_{t-1}^m - \Gamma_t^m\|^2.
\label{eq:adam_IV_final}
\end{align}

%-----------------------------------------------------------------------------
% Step 3e: Term V
%-----------------------------------------------------------------------------
\paragraph{Step 3e: Bounding term $V$.}

This term measures the cost of approximating $\nabla f(z_t)$ by
$\nabla f(x_t)$ in the inner product with $\Gamma_t\,g_t$.

First, we bound the ``virtual drift'' $\|\nabla f(x_t) - \nabla f(z_t)\|$.
By $L$-smoothness:
\begin{equation}
\|\nabla f(x_t) - \nabla f(z_t)\| \leq L\,\|x_t - z_t\|.
\label{eq:adam_V_smooth}
\end{equation}
From the definition~\eqref{eq:adam_virtual}:
\begin{align}
x_t - z_t
&= x_t - \frac{1}{1-\beta_1}\,x_t + \frac{\beta_1}{1-\beta_1}\,x_{t-1}
\notag \\
&= -\frac{\beta_1}{1-\beta_1}\,(x_t - x_{t-1})
\notag \\
&= \frac{\beta_1}{1-\beta_1}\,d_{t-1}
= \frac{\eta\,\beta_1}{1-\beta_1}\,
\mathbb{E}_m[\Gamma_{t-1}^m\,u_{t-1}^m].
\label{eq:adam_V_xz}
\end{align}
Taking norms and using $\|\Gamma_{t-1}^m\| \leq 1/\lambda$,
$\|u_{t-1}^m\| \leq G$:
\begin{equation}
\|x_t - z_t\|
\leq \frac{\eta\,\beta_1}{1-\beta_1}\,
\mathbb{E}_m[\|\Gamma_{t-1}^m\|\,\|u_{t-1}^m\|]
\leq \frac{\eta\,\beta_1\,G}{(1-\beta_1)\,\lambda}.
\label{eq:adam_V_xz_bound}
\end{equation}
Hence:
\begin{equation}
\|\nabla f(x_t) - \nabla f(z_t)\|
\leq \frac{\eta\,\beta_1\,L\,G}{(1-\beta_1)\,\lambda}.
\label{eq:adam_V_grad_bound}
\end{equation}

Now we bound $V$ itself. Split $\Gamma_t = \Gamma_{t-1} + (\Gamma_t - \Gamma_{t-1})$:
\begin{align}
V &= \eta\,\langle \nabla f(x_t) - \nabla f(z_t),\, \Gamma_t\,g_t \rangle
\notag \\
&= \eta\,\langle \nabla f(x_t) - \nabla f(z_t),\, \Gamma_{t-1}\,g_t \rangle
   + \eta\,\langle \nabla f(x_t) - \nabla f(z_t),\,
   (\Gamma_t - \Gamma_{t-1})\,g_t \rangle.
\label{eq:adam_V_split}
\end{align}

For the second term, using~\eqref{eq:adam_V_grad_bound}:
\begin{equation}
|\langle \nabla f(x_t) - \nabla f(z_t),\,
(\Gamma_t - \Gamma_{t-1})\,g_t \rangle|
\leq \frac{\eta\,\beta_1\,L\,G^2}{(1-\beta_1)\,\lambda}\,
\|\Gamma_t - \Gamma_{t-1}\|.
\label{eq:adam_V_perturb}
\end{equation}

For the first term, using unbiasedness and splitting
$\frac{1}{M}\sum_m \nabla f_m(x_t^m) = \nabla f(x_t) +
(\frac{1}{M}\sum_m \nabla f_m(x_t^m) - \nabla f(x_t))$, then applying
Young's inequality $\langle a,b \rangle \leq \frac{1}{2\rho}\|a\|^2 +
\frac{\rho}{2}\|b\|^2$ to each piece:
\begin{align}
&\mathbb{E}\bigl[\eta\,\langle \nabla f(x_t) - \nabla f(z_t),\,
\Gamma_{t-1}\,g_t \rangle\bigr] \notag \\
&\quad\leq \frac{\eta}{\lambda}\,\mathbb{E}\!\left[
\frac{1}{2\rho}\,\|\nabla f(x_t) - \nabla f(z_t)\|^2
+ \frac{\rho}{2}\,\|\nabla f(x_t)\|^2\right] \notag \\
&\qquad+ \frac{\eta}{\lambda}\,\mathbb{E}\!\left[
\frac{1}{2}\,\|\nabla f(x_t) - \nabla f(z_t)\|^2
+ \frac{L^2}{2M}\sum_m \|x_t^m - x_t\|^2\right].
\label{eq:adam_V_young}
\end{align}

Substituting~\eqref{eq:adam_V_grad_bound} into the
$\|\nabla f(x_t) - \nabla f(z_t)\|^2$ terms, all such terms are
$\mathcal{O}(\eta^2)$. Collecting:
\begin{align}
\mathbb{E}[V]
&\leq \frac{\eta\,\rho}{2\lambda}\,\mathbb{E}\|\nabla f(x_t)\|^2
  + \mathcal{O}(\eta)\cdot\frac{1}{M}\sum_m\mathbb{E}\|x_t^m - x_t\|^2
  \notag \\
&\quad+ \mathcal{O}\!\left(\frac{\eta^3}{\rho}\right)
  + \frac{\eta^2\,\beta_1\,L\,G^2}{(1-\beta_1)\,\lambda}\,
  \mathbb{E}\|\Gamma_t - \Gamma_{t-1}\|.
\label{eq:adam_V_final}
\end{align}

%-----------------------------------------------------------------------------
% Summary of all five bounds
%-----------------------------------------------------------------------------
\paragraph{Summary of the five bounds.}

Collecting~\eqref{eq:adam_I_final},~\eqref{eq:adam_II_final},
\eqref{eq:adam_III_final},~\eqref{eq:adam_IV_final},
and~\eqref{eq:adam_V_final}, and suppressing constant factors into
$\mathcal{O}(\cdot)$ notation:
\begin{align}
I  &\leq -\frac{\eta}{2C_0}\,\mathbb{E}\|\nabla f(x_t)\|^2
   + \mathcal{O}(\eta)\cdot\frac{1}{M}\sum_m\mathbb{E}\|x_t - x_t^m\|^2
   + \mathcal{O}(\eta)\cdot\frac{1}{M}\sum_m
   \mathbb{E}\|\Gamma_{t-1}^m - \Gamma_t^m\|,
\label{eq:adam_sum_I} \\[4pt]
II &\leq \mathcal{O}(\eta)\cdot\frac{1}{M}\sum_m
   \mathbb{E}\|g_t - \widetilde{g}_t^m\|,
\label{eq:adam_sum_II} \\[4pt]
III &\leq \mathcal{O}(\eta)\cdot\frac{1}{M}\sum_m
   \mathbb{E}\|\Gamma_{t-1}^m - \Gamma_t^m\|,
\label{eq:adam_sum_III} \\[4pt]
IV &\leq \mathcal{O}(\eta^2),
\label{eq:adam_sum_IV} \\[4pt]
V  &\leq \frac{\eta\,\rho}{2\lambda}\,\mathbb{E}\|\nabla f(x_t)\|^2
   + \mathcal{O}(\eta)\cdot\frac{1}{M}\sum_m\mathbb{E}\|x_t^m - x_t\|^2
   + \mathcal{O}(\eta^2).
\label{eq:adam_sum_V}
\end{align}

To obtain an $\mathcal{O}(1/\sqrt{T})$ convergence rate, we set
$\rho = \lambda/(2C_0)$ so that the $\|\nabla f(x_t)\|^2$ coefficient in
$V$ exactly cancels against the negative coefficient in $I$, and then verify
three conditions:
\begin{alignat}{2}
&\text{(A1)}\quad&
\frac{1}{TM}\sum_{t,m}\mathbb{E}\|x_t^m - x_t\|^2
&= \mathcal{O}(\eta^2),
\qquad\text{(extension of Lemma~\ref{lem:bound})}
\label{eq:adam_cond_A1} \\[4pt]
&\text{(A2)}\quad&
\sum_{t=0}^{T-1}\mathbb{E}\|\Gamma_{t-1}^m - \Gamma_t^m\|
&= \mathcal{O}(1),
\qquad\text{(AMSGrad telescoping)}
\label{eq:adam_cond_A2} \\[4pt]
&\text{(A3)}\quad&
\frac{1}{M}\sum_{t,m}\mathbb{E}\|g_t - \widetilde{g}_t^m\|
&= \mathcal{O}(1).
\qquad\text{(double geometric sum)}
\label{eq:adam_cond_A3}
\end{alignat}

\medskip
\noindent\textbf{Condition (A2): AMSGrad telescoping.}
Since $\tilde{v}_t^m = \max(v_t^m, \tilde{v}_{t-1}^m)$ coordinate-wise,
the sequence $\{[\tilde{v}_t^m]_i\}_{t \geq 0}$ is non-decreasing and bounded
above by $G^2$ for each coordinate $i$. Therefore
$[\Gamma_t^m]_{ii} = 1/\sqrt{[\tilde{v}_t^m]_i + \lambda^2}$ is non-increasing,
and
\begin{equation}
\sum_{t=0}^{T-1}\|\Gamma_{t-1}^m - \Gamma_t^m\|
= \sum_{t=0}^{T-1}\max_i\bigl|[\Gamma_{t-1}^m]_{ii} - [\Gamma_t^m]_{ii}\bigr|
\leq \sum_i\bigl([\Gamma_{-1}^m]_{ii} - [\Gamma_{T-1}^m]_{ii}\bigr)
\leq \frac{d}{\lambda},
\label{eq:adam_AMSGrad_telescope}
\end{equation}
which is $\mathcal{O}(1)$ with respect to $T$.

\medskip
\noindent\textbf{Condition (A3): double geometric sum for the pseudo-gradient drift.}
We show that $\mathbb{E}\|g_t - \widetilde{g}_t^m\|$ has the same double
geometric sum structure as the parameter drift in Lemma~\ref{lem:bound}.

First, expand the momentum drift. With probability $p_u$ (sync at step $t$),
$u_t^m = u_t$, so $\|u_t - u_t^m\| = 0$. With probability $1 - p_u$:
\begin{align}
\|u_t - u_t^m\|
&= \|\beta_1\,u_{t-1} + (1-\beta_1)\,g_t
   - \beta_1\,u_{t-1}^m - (1-\beta_1)\,g_t^m\| \notag \\
&= \|\beta_1\,(u_{t-1} - u_{t-1}^m) + (1-\beta_1)\,(g_t - g_t^m)\|
\notag \\
&\leq \beta_1\,\|u_{t-1} - u_{t-1}^m\|
  + (1-\beta_1)\,\|g_t - g_t^m\|.
\label{eq:adam_u_drift_one}
\end{align}
Taking expectation over the sync coin:
\begin{equation}
\mathbb{E}\|u_t - u_t^m\|
\leq (1-p_u)\,\beta_1\,\mathbb{E}\|u_{t-1} - u_{t-1}^m\|
  + (1-p_u)\,(1-\beta_1)\,\mathbb{E}\|g_t - g_t^m\|.
\label{eq:adam_u_drift_recur}
\end{equation}
Unrolling~\eqref{eq:adam_u_drift_recur} from $\tau = 0$ to $t$
(with $u_{-1} = u_{-1}^m = 0$):
\begin{equation}
\mathbb{E}\|u_t - u_t^m\|
\leq (1-p_u)\,(1-\beta_1)\sum_{\tau=0}^{t}
\bigl((1-p_u)\,\beta_1\bigr)^{t-\tau}\,
\mathbb{E}\|g_\tau - g_\tau^m\|.
\label{eq:adam_u_drift_unrolled}
\end{equation}

Now bound $\mathbb{E}\|g_t - \widetilde{g}_t^m\|$.
By definition~\eqref{eq:gtilde_def}:
\begin{align}
g_t - \widetilde{g}_t^m
&= \frac{u_t - \beta_1\,u_{t-1}}{1-\beta_1}
   - \frac{u_t^m - \beta_1\,u_{t-1}^m}{1-\beta_1} \notag \\
&= \frac{(u_t - u_t^m) - \beta_1\,(u_{t-1} - u_{t-1}^m)}{1-\beta_1}.
\label{eq:adam_gtilde_diff}
\end{align}
By the triangle inequality:
\begin{align}
\|g_t - \widetilde{g}_t^m\|
&\leq \frac{1}{1-\beta_1}\,\|u_t - u_t^m\|
  + \frac{\beta_1}{1-\beta_1}\,\|u_{t-1} - u_{t-1}^m\| \notag \\
&= \frac{1}{1-\beta_1}\sum_{\tau=t-1}^{t}
\beta_1^{t-\tau}\,\|u_\tau - u_\tau^m\|.
\label{eq:adam_gtilde_bound}
\end{align}
Substituting~\eqref{eq:adam_u_drift_unrolled}
into~\eqref{eq:adam_gtilde_bound} and letting
$q_2 = (1-p_u)\,\beta_1$:
\begin{align}
\mathbb{E}\|g_t - \widetilde{g}_t^m\|
&\leq (1-p_u)\sum_{\tau=t-1}^{t}\sum_{\nu=0}^{\tau}
\beta_1^{t-\tau}\,q_2^{\tau-\nu}\,\mathbb{E}\|g_\nu - g_\nu^m\|
\notag \\
&= \sum_{\tau=t}^{t+1}\sum_{\nu=0}^{\tau-1}
\beta_1^{t+1-\tau}\,q_2^{\tau-\nu}\,\mathbb{E}\|g_\nu - g_\nu^m\|,
\label{eq:adam_gtilde_double_geom}
\end{align}
which has the same \emph{double geometric sum} structure as
equation~\eqref{eq:lem3_double_geom} in Lemma~\ref{lem:bound}. By the same
summation technique (swapping the order of summation and bounding the
geometric series), the total over $t = 0, \ldots, T-1$ is $\mathcal{O}(1)$,
confirming condition~\eqref{eq:adam_cond_A3}.

\paragraph{Remark (SP invariance and SP-tightened constants for Adam).}
As in the SGDM analysis, the three conditions
\eqref{eq:adam_cond_A1}--\eqref{eq:adam_cond_A3} depend only on the
synchronization probabilities $p_x, p_u, p_v$, the decay factors
$\beta_1, \beta_2$, and the gradient bound $G$. None involves the SP
degree $P$. The convergence rate of FAUST-Adam is therefore
$P$-independent in order.

The SP-tightened constant refinement from the SGDM analysis
(Eq.~\eqref{eq:psi_attn_sp}) extends naturally to Adam's two-moment
structure. Since the first-moment drift~\eqref{eq:drift_u_sp} and the
second-moment drift~\eqref{eq:drift_v_sp} both carry the $1/\sqrt{P}$
factor for attention-layer coordinates, the combined drift contribution
from these coordinates to the preconditioned step
$d_t^m = \eta_t \Gamma_t u_t^m$ (where $\Gamma_t = \mathrm{diag}(v_t + \lambda^2)^{-1/2}$)
is also reduced by $1/\sqrt{P}$.
Concretely, the $\tilde{g}$-drift bound~\eqref{eq:adam_gtilde_bound} for the
attention-layer coordinates $j \in \mathcal{I}_{\text{attn}}$ becomes
$\|\tilde{g}_{t,j} - \tilde{g}_{t-1,j}\|_\infty \leq 2G\,\gamma_{\max}/\sqrt{P}$,
tightening the spectral-norm bound on $\Gamma_t - \Gamma_{t-1}$ and
yielding a sharper constant in~\eqref{eq:adam_AMSGrad_telescope}.

The collective orthogonality constraint (fence at $K_x$-boundaries) is
again a necessary condition: for Adam, the second moment $v_t^m$ depends
on $g_t^m \odot g_t^m$, so a partially-reduced gradient would corrupt both
moments simultaneously, making the preconditioner $\Gamma_t$ inconsistent
across replicas and invalidating the spectral bound chain.

%=============================================================================
% D.2 Key Lemmas
%=============================================================================
\subsection{Key Lemmas}
\label{sec:key_lemmas}

We now prove the two auxiliary lemmas used in the main proof.

%-----------------------------------------------------------------------------
% Lemma 2: Virtual iterate drift
%-----------------------------------------------------------------------------
\begin{lemma}[Virtual iterate drift]
\label{lem:virt_drift}
For all $T \geq 1$,
\begin{equation}
\sum_{t=0}^{T-1}\mathbb{E}\|z_t - x_t\|^2
\leq \frac{\eta^2\,\beta^2}{(1-\beta)^2\,M}\,T\,\sigma^2
   + \frac{\eta^2\,\beta^2}{(1-\beta)^2}\sum_{t=0}^{T-1}
   \mathbb{E}\left\|\frac{1}{M}\sum_m\nabla f_m(x_t^m)\right\|^2.
\label{eq:lem_virt_statement}
\end{equation}
\end{lemma}

\begin{proof}
\noindent\textbf{Step (i): express $z_t - x_t$ in terms of momenta.}

From the definition~\eqref{eq:virtual} and the parameter
recursion~\eqref{eq:avg_x_recursion}:
\begin{align}
z_t - x_t
&= \frac{1}{1-\beta}\,x_t - \frac{\beta}{1-\beta}\,x_{t-1} - x_t
\notag \\
&= -\frac{\beta}{1-\beta}\,(x_t - x_{t-1})
\notag \\
&= -\frac{\beta}{1-\beta}\,(-\eta\,u_{t-1})
= \frac{\eta\,\beta}{1-\beta}\,u_{t-1}.
\label{eq:lem2_zx_raw}
\end{align}
Wait---we need to be more careful. Actually, for $t \geq 1$:
$x_t - x_{t-1} = -\eta\,u_{t-1}$ by~\eqref{eq:avg_x_recursion}. So:
\begin{equation}
z_t - x_t = \frac{\eta\,\beta}{1-\beta}\,u_{t-1}.
\label{eq:lem2_zx}
\end{equation}
For $t = 0$: $z_0 = x_0$, so $z_0 - x_0 = 0$, which is consistent with
$u_{-1} = 0$.

\medskip
\noindent\textbf{Step (ii): unroll the momentum.}

Since $u_{-1} = 0$, unrolling the recursion
$u_t = \beta\,u_{t-1} + (1-\beta)\,g_t$:
\begin{align}
u_0 &= (1-\beta)\,g_0, \notag \\
u_1 &= \beta\,(1-\beta)\,g_0 + (1-\beta)\,g_1
= (1-\beta)\sum_{\tau=0}^{1}\beta^{1-\tau}\,g_\tau, \notag \\
&\;\;\vdots \notag \\
u_t &= (1-\beta)\sum_{\tau=0}^{t}\beta^{t-\tau}\,g_\tau.
\label{eq:lem2_u_unrolled}
\end{align}
Substituting into~\eqref{eq:lem2_zx}:
\begin{equation}
z_t - x_t = \eta\,\beta\sum_{\tau=0}^{t-1}\beta^{t-1-\tau}\,g_\tau
= \eta\,\beta\sum_{\tau=0}^{t-1}\beta^{t-1-\tau}\,g_\tau.
\label{eq:lem2_zx_expanded}
\end{equation}
Re-indexing: let $k = t - 1 - \tau$, then $\tau = t - 1 - k$ and the sum runs
$k = 0, \ldots, t-1$:
\begin{equation}
z_t - x_t = \eta\,\beta\sum_{k=0}^{t-1}\beta^k\,g_{t-1-k}.
\label{eq:lem2_reindex}
\end{equation}

\medskip
\noindent\textbf{Step (iii): bound $\|z_t - x_t\|^2$ using Jensen.}

Define the weights $w_k = \beta^k / s_{t-1}$ where
$s_{t-1} = \sum_{k=0}^{t-1}\beta^k = (1 - \beta^t)/(1-\beta) \leq 1/(1-\beta)$.
These are non-negative and sum to $1$, so by Jensen's inequality
($\|\sum_k w_k\,a_k\|^2 \leq \sum_k w_k\,\|a_k\|^2$):
\begin{align}
\|z_t - x_t\|^2
&= \eta^2\,\beta^2\,s_{t-1}^2\,
\left\|\sum_{k=0}^{t-1}\frac{\beta^k}{s_{t-1}}\,g_{t-1-k}\right\|^2
\notag \\
&\leq \eta^2\,\beta^2\,s_{t-1}^2\sum_{k=0}^{t-1}
\frac{\beta^k}{s_{t-1}}\,\|g_{t-1-k}\|^2
\notag \\
&= \eta^2\,\beta^2\,s_{t-1}\sum_{k=0}^{t-1}
\beta^k\,\|g_{t-1-k}\|^2
\notag \\
&\leq \frac{\eta^2\,\beta^2}{1-\beta}\sum_{\tau=0}^{t-1}
\beta^{t-1-\tau}\,\|g_\tau\|^2,
\label{eq:lem2_jensen}
\end{align}
where in the last step we substituted $\tau = t - 1 - k$ and used
$s_{t-1} \leq 1/(1-\beta)$.

\medskip
\noindent\textbf{Step (iv): sum over $t$ and swap summation order.}

\begin{align}
\sum_{t=0}^{T-1}\mathbb{E}\|z_t - x_t\|^2
&\leq \frac{\eta^2\,\beta^2}{1-\beta}\sum_{t=0}^{T-1}
\sum_{\tau=0}^{t-1}\beta^{t-1-\tau}\,\mathbb{E}\|g_\tau\|^2.
\label{eq:lem2_sum_raw}
\end{align}
(The $t = 0$ term contributes $0$ since the inner sum is empty.)

Swap the order of summation. The inner sum has $\tau$ ranging from $0$ to
$t-1$ and $t$ from $1$ to $T-1$. For a fixed $\tau$, $t$ ranges from
$\tau + 1$ to $T-1$:
\begin{align}
\sum_{t=1}^{T-1}\sum_{\tau=0}^{t-1}\beta^{t-1-\tau}\,\mathbb{E}\|g_\tau\|^2
&= \sum_{\tau=0}^{T-2}\sum_{t=\tau+1}^{T-1}
\beta^{t-1-\tau}\,\mathbb{E}\|g_\tau\|^2 \notag \\
&= \sum_{\tau=0}^{T-2}\left(\sum_{j=0}^{T-2-\tau}\beta^j\right)
\mathbb{E}\|g_\tau\|^2
\qquad\text{(set $j = t - 1 - \tau$)} \notag \\
&= \sum_{\tau=0}^{T-2}\frac{1 - \beta^{T-1-\tau}}{1-\beta}\,
\mathbb{E}\|g_\tau\|^2 \notag \\
&\leq \frac{1}{1-\beta}\sum_{\tau=0}^{T-1}\mathbb{E}\|g_\tau\|^2.
\label{eq:lem2_swap}
\end{align}
Hence:
\begin{equation}
\sum_{t=0}^{T-1}\mathbb{E}\|z_t - x_t\|^2
\leq \frac{\eta^2\,\beta^2}{(1-\beta)^2}\sum_{\tau=0}^{T-1}
\mathbb{E}\|g_\tau\|^2.
\label{eq:lem2_before_decomp}
\end{equation}

\medskip
\noindent\textbf{Step (v): decompose $\|g_\tau\|^2$ into signal and noise.}

Recall $g_\tau = \frac{1}{M}\sum_m g_\tau^m$. Write
$g_\tau^m = \nabla f_m(x_\tau^m) + \varepsilon_\tau^m$ where
$\varepsilon_\tau^m = g_\tau^m - \nabla f_m(x_\tau^m)$ is the noise. Then:
\begin{align}
\|g_\tau\|^2
&= \left\|\frac{1}{M}\sum_m \varepsilon_\tau^m
   + \frac{1}{M}\sum_m \nabla f_m(x_\tau^m)\right\|^2 \notag \\
&= \left\|\frac{1}{M}\sum_m \varepsilon_\tau^m\right\|^2
   + \left\|\frac{1}{M}\sum_m \nabla f_m(x_\tau^m)\right\|^2
   + 2\left\langle \frac{1}{M}\sum_m \varepsilon_\tau^m,\,
   \frac{1}{M}\sum_m \nabla f_m(x_\tau^m)\right\rangle.
\label{eq:lem2_decomp_raw}
\end{align}
Taking expectation, the cross-term vanishes by unbiasedness
($\mathbb{E}[\varepsilon_\tau^m] = 0$):
\begin{align}
\mathbb{E}\|g_\tau\|^2
&= \mathbb{E}\left\|\frac{1}{M}\sum_m \varepsilon_\tau^m\right\|^2
   + \mathbb{E}\left\|\frac{1}{M}\sum_m
   \nabla f_m(x_\tau^m)\right\|^2.
\label{eq:lem2_decomp_clean}
\end{align}
For the noise term, by independence across replicas:
\begin{equation}
\mathbb{E}\left\|\frac{1}{M}\sum_m \varepsilon_\tau^m\right\|^2
= \frac{1}{M^2}\sum_m \mathbb{E}\|\varepsilon_\tau^m\|^2
\leq \frac{\sigma^2}{M}.
\label{eq:lem2_noise}
\end{equation}

\medskip
\noindent\textbf{Step (vi): assemble.}

Substituting~\eqref{eq:lem2_decomp_clean} and~\eqref{eq:lem2_noise}
into~\eqref{eq:lem2_before_decomp}:
\begin{align}
\sum_{t=0}^{T-1}\mathbb{E}\|z_t - x_t\|^2
&\leq \frac{\eta^2\,\beta^2}{(1-\beta)^2}\sum_{\tau=0}^{T-1}
\left(\frac{\sigma^2}{M}
+ \mathbb{E}\left\|\frac{1}{M}\sum_m
\nabla f_m(x_\tau^m)\right\|^2\right) \notag \\
&= \frac{\eta^2\,\beta^2}{(1-\beta)^2\,M}\,T\,\sigma^2
+ \frac{\eta^2\,\beta^2}{(1-\beta)^2}\sum_{\tau=0}^{T-1}
\mathbb{E}\left\|\frac{1}{M}\sum_m
\nabla f_m(x_\tau^m)\right\|^2.
\label{eq:lem2_done}
\end{align}
This completes the proof.
\end{proof}

%-----------------------------------------------------------------------------
% Lemma 3: Client drift bound
%-----------------------------------------------------------------------------
\begin{lemma}[Client drift bound]
\label{lem:bound}
If $24\eta^2\,L^2\,\psi \leq 1$, then
\begin{equation}
\frac{1}{MT}\sum_{t=0}^{T-1}\sum_m\mathbb{E}\|x_t - x_t^m\|^2
\leq 12\eta^2\,(B^2-1)\,\psi\cdot
\frac{1}{T}\sum_{t=0}^{T-1}\mathbb{E}\|\nabla f(x_t)\|^2
+ 4\eta^2\,\psi\,(\sigma^2 + 3G^2),
\label{eq:lem_bound_statement}
\end{equation}
where the drift coefficient is
\begin{equation}
\psi \;=\; \frac{4(1-p_x)}{p_x^2}\cdot
\frac{(1-\beta)(1-p_u)}{1-(1-p_u)\beta}.
\label{eq:psi_def_lemma}
\end{equation}
\end{lemma}

\begin{proof}
The proof proceeds in five steps: (i) expand the parameter drift using the
probabilistic update; (ii) expand the momentum drift; (iii) combine via a
double geometric sum; (iv) bound the gradient variance; (v) close the
recursion.

\medskip
\noindent\textbf{Step (i): parameter drift recursion.}

At step $t$, with probability $p_x$ (sync), $x_{t+1}^m = \mathbb{E}_m[\cdots]$,
so $x_{t+1} - x_{t+1}^m = 0$. With probability $1 - p_x$ (no sync):
\begin{align}
x_{t+1} - x_{t+1}^m
&= (x_t - \eta\,u_t) - (x_t^m - \eta\,u_t^m) \notag \\
&= (x_t - x_t^m) - \eta\,(u_t - u_t^m).
\label{eq:lem3_param_nosync}
\end{align}
Taking squared norms:
\begin{equation}
\|x_{t+1} - x_{t+1}^m\|^2
= \|(x_t - x_t^m) - \eta\,(u_t - u_t^m)\|^2.
\label{eq:lem3_param_sq}
\end{equation}
By the Young-type inequality $\|a - b\|^2 \leq (1+s)\,\|a\|^2 +
(1 + 1/s)\,\|b\|^2$ for any $s > 0$ (which follows from
$\|a - b\|^2 = \|a\|^2 + \|b\|^2 - 2\langle a, b\rangle \leq
(1+s)\|a\|^2 + (1+1/s)\|b\|^2$ by Young applied to the cross-term):
\begin{equation}
\|x_{t+1} - x_{t+1}^m\|^2
\leq (1+s)\,\|x_t - x_t^m\|^2
+ \eta^2\,(1 + \tfrac{1}{s})\,\|u_t - u_t^m\|^2.
\label{eq:lem3_young}
\end{equation}
Taking expectation over the sync coin:
\begin{equation}
\mathbb{E}\|x_{t+1} - x_{t+1}^m\|^2
\leq (1-p_x)(1+s)\,\mathbb{E}\|x_t - x_t^m\|^2
+ \eta^2\,(1-p_x)(1+\tfrac{1}{s})\,\mathbb{E}\|u_t - u_t^m\|^2.
\label{eq:lem3_param_recur}
\end{equation}
Let $q_1 \stackrel{\text{def}}{=} (1-p_x)(1+s)$. Unrolling~\eqref{eq:lem3_param_recur}
from $t = 0$ (using $x_0 - x_0^m = 0$):
\begin{equation}
\mathbb{E}\|x_{t+1} - x_{t+1}^m\|^2
\leq \eta^2\,(1-p_x)(1+\tfrac{1}{s})
\sum_{\tau=0}^{t}q_1^{t-\tau}\,\mathbb{E}\|u_\tau - u_\tau^m\|^2.
\label{eq:lem3_param_unrolled}
\end{equation}

\medskip
\noindent\textbf{Step (ii): momentum drift recursion.}

With probability $p_u$ (sync), $u_t^m = u_t$. With probability $1 - p_u$:
\begin{align}
u_t - u_t^m
&= \beta\,(u_{t-1} - u_{t-1}^m) + (1-\beta)\,(g_t - g_t^m).
\label{eq:lem3_mom_nosync}
\end{align}
By the inequality $\|\alpha\,a + (1-\alpha)\,b\|^2 \leq
\alpha\,\|a\|^2 + (1-\alpha)\,\|b\|^2$ (Jensen, since $\beta \in [0,1)$):
\begin{equation}
\|u_t - u_t^m\|^2
\leq \beta\,\|u_{t-1} - u_{t-1}^m\|^2
+ (1-\beta)\,\|g_t - g_t^m\|^2.
\label{eq:lem3_mom_jensen}
\end{equation}
Taking expectation over the sync coin:
\begin{equation}
\mathbb{E}\|u_t - u_t^m\|^2
\leq (1-p_u)\,\beta\,\mathbb{E}\|u_{t-1} - u_{t-1}^m\|^2
+ (1-p_u)\,(1-\beta)\,\mathbb{E}\|g_t - g_t^m\|^2.
\label{eq:lem3_mom_recur}
\end{equation}
Let $q_2 \stackrel{\text{def}}{=} (1-p_u)\,\beta$. Unrolling (with $u_{-1} = u_{-1}^m = 0$):
\begin{equation}
\mathbb{E}\|u_t - u_t^m\|^2
\leq (1-p_u)\,(1-\beta)\sum_{\tau=0}^{t}
q_2^{t-\tau}\,\mathbb{E}\|g_\tau - g_\tau^m\|^2
\leq \frac{1-\beta}{\beta}\sum_{\tau=0}^{t}
q_2^{t+1-\tau}\,\mathbb{E}\|g_\tau - g_\tau^m\|^2.
\label{eq:lem3_mom_unrolled}
\end{equation}

\medskip
\noindent\textbf{Step (iii): combine into a double geometric sum.}

Substitute~\eqref{eq:lem3_mom_unrolled} into~\eqref{eq:lem3_param_unrolled}:
\begin{align}
&\frac{1}{M}\sum_m\mathbb{E}\|x_t - x_t^m\|^2 \notag \\
&\quad\leq \eta^2\,(1-p_x)(1+\tfrac{1}{s})\,\frac{1-\beta}{\beta}
\sum_{\tau=0}^{t-1}\sum_{\nu=0}^{\tau}
q_1^{t-1-\tau}\,q_2^{\tau+1-\nu}\,
\frac{1}{M}\sum_m\mathbb{E}\|g_\nu - g_\nu^m\|^2 \notag \\
&\quad= \eta^2\,\underbrace{(1-p_x)(1+\tfrac{1}{s})(1-\beta)(1-p_u)
}_{\stackrel{\text{def}}{=}\;\phi}
\sum_{\nu=0}^{t-1}\frac{q_1^{t-\nu} - q_2^{t-\nu}}{q_1 - q_2}\,
\frac{1}{M}\sum_m\mathbb{E}\|g_\nu - g_\nu^m\|^2,
\label{eq:lem3_double_geom}
\end{align}
where the last line evaluates the inner convolution
$\sum_{\tau=\nu}^{t-1}q_1^{t-1-\tau}\,q_2^{\tau+1-\nu}$ via the standard
geometric-series identity
$\sum_{k=0}^{n}a^k\,b^{n-k} = (a^{n+1} - b^{n+1})/(a - b)$.

\medskip
\noindent\textbf{Step (iv): time-average and optimize $s$.}

Average~\eqref{eq:lem3_double_geom} over $t = 0, \ldots, T-1$, swap the
summation order, and bound the geometric series:
\begin{equation}
\frac{1}{MT}\sum_{t,m}\mathbb{E}\|x_t - x_t^m\|^2
\leq \frac{\eta^2\,\phi}{(1-q_1)(1-q_2)}\,
\frac{1}{T}\sum_{\tau=0}^{T-1}
\frac{1}{M}\sum_m\mathbb{E}\|g_\tau - g_\tau^m\|^2.
\label{eq:lem3_averaged}
\end{equation}

We minimize $\phi/((1-q_1)(1-q_2))$ over $s > 0$. The factor depending on $s$
is
\begin{equation}
h(s) = \frac{(1-p_x)(1+1/s)}{1 - (1-p_x)(1+s)}.
\label{eq:lem3_hs}
\end{equation}
Setting $h'(s) = 0$ yields $s^* = 1/\sqrt{1-p_x} - 1$, giving
\begin{align}
h(s^*)
&= \frac{(1-p_x)(1 + \sqrt{1-p_x})^2}{p_x^2}
\leq \frac{4(1-p_x)}{p_x^2},
\label{eq:lem3_hs_optimal}
\end{align}
since $(1 + \sqrt{1-p_x})^2 \leq 4$ for $p_x \in [0,1]$.

Therefore:
\begin{equation}
\frac{\phi}{(1-q_1)(1-q_2)}
\leq \frac{4(1-p_x)}{p_x^2}\cdot
\frac{(1-\beta)(1-p_u)}{1-(1-p_u)\beta}
= \psi.
\label{eq:lem3_psi_bound}
\end{equation}

\medskip
\noindent\textbf{Step (v): bound the gradient variance and close.}

Using Lemma~\ref{lem:het} (proved below):
\begin{align}
\frac{1}{M}\sum_m\mathbb{E}\|g_t - g_t^m\|^2
&\leq 2\sigma^2 + \frac{12L^2}{M}\sum_m\mathbb{E}\|x_t - x_t^m\|^2
+ 6(B^2-1)\,\mathbb{E}\|\nabla f(x_t)\|^2 + 6G^2.
\label{eq:lem3_grad_var}
\end{align}
Substituting into~\eqref{eq:lem3_averaged}:
\begin{align}
\frac{1}{MT}\sum_{t,m}\mathbb{E}\|x_t - x_t^m\|^2
&\leq 12\eta^2\,L^2\,\psi\cdot
\frac{1}{TM}\sum_{t,m}\mathbb{E}\|x_t - x_t^m\|^2 \notag \\
&\quad+ 6\eta^2\,(B^2-1)\,\psi\cdot
\frac{1}{T}\sum_t\mathbb{E}\|\nabla f(x_t)\|^2
+ 2\eta^2\,\psi\,(\sigma^2 + 3G^2).
\label{eq:lem3_self_ref}
\end{align}
Denote the left-hand side by $\Delta$. Then
$\Delta \leq 12\eta^2\,L^2\,\psi\,\Delta + (\cdots)$.
If $12\eta^2\,L^2\,\psi \leq 1/2$ (guaranteed by
$24\eta^2\,L^2\,\psi \leq 1$), then
$\Delta(1 - 12\eta^2\,L^2\,\psi) \geq \Delta/2$, so:
\begin{equation}
\Delta \leq 12\eta^2\,(B^2-1)\,\psi\cdot
\frac{1}{T}\sum_t\mathbb{E}\|\nabla f(x_t)\|^2
+ 4\eta^2\,\psi\,(\sigma^2 + 3G^2).
\label{eq:lem3_closed}
\end{equation}
\end{proof}

%-----------------------------------------------------------------------------
% Lemma 4: Heterogeneity decomposition
%-----------------------------------------------------------------------------
\begin{lemma}[Heterogeneity decomposition]
\label{lem:het}
Under Assumptions~\ref{asm:smooth} and~\ref{asm:heterogeneity},
\begin{equation}
\frac{1}{M}\sum_m\left\|\nabla f_m(x_t^m)
- \frac{1}{M}\sum_i\nabla f_i(x_t^i)\right\|^2
\leq \frac{6L^2}{M}\sum_m\|x_t - x_t^m\|^2
+ 3(B^2-1)\,\|\nabla f(x_t)\|^2 + 3G^2.
\label{eq:lem_het_statement}
\end{equation}
\end{lemma}

\begin{proof}
For each replica $m$, write
\begin{equation}
\begin{aligned}
\nabla f_m(x_t^m) - \frac{1}{M}\sum_i \nabla f_i(x_t^i)
&= \underbrace{\bigl(\nabla f_m(x_t^m)-\nabla f_m(x_t)\bigr)}_{A_m}
 + \underbrace{\bigl(\nabla f_m(x_t)-\nabla f(x_t)\bigr)}_{B_m} \\
&\quad
 + \underbrace{\left(\nabla f(x_t)-\frac{1}{M}\sum_i\nabla f_i(x_t^i)\right)}_{C}.
\end{aligned}
\label{eq:lem4_abc}
\end{equation}
By the inequality $\|A + B + C\|^2 \leq 3(\|A\|^2 + \|B\|^2 + \|C\|^2)$
(which follows from $\|A + B + C\|^2 \leq (3\|A\|)(3^{-1}\|A\|) + \cdots$
and Cauchy--Schwarz, or equivalently from Jensen applied to the convex function
$\|\cdot\|^2$):
\begin{align}
\left\|\nabla f_m(x_t^m) - \frac{1}{M}\sum_i\nabla f_i(x_t^i)\right\|^2
&\leq 3\,\|A_m\|^2 + 3\,\|B_m\|^2 + 3\,\|C\|^2.
\label{eq:lem4_triangle}
\end{align}

Average over $m$:
\begin{align}
\frac{1}{M}\sum_m(\cdots)
&\leq \frac{3}{M}\sum_m\|\nabla f_m(x_t^m) - \nabla f_m(x_t)\|^2
\notag \\
&\quad+ \frac{3}{M}\sum_m\|\nabla f_m(x_t) - \nabla f(x_t)\|^2
\notag \\
&\quad+ 3\left\|\nabla f(x_t) - \frac{1}{M}\sum_i\nabla f_i(x_t^i)\right\|^2.
\label{eq:lem4_three_terms}
\end{align}

Bound each:

\noindent\emph{First term:} By $L$-smoothness,
$\|\nabla f_m(x_t^m) - \nabla f_m(x_t)\|^2 \leq L^2\,\|x_t^m - x_t\|^2$.

\noindent\emph{Third term:} Note that
$\nabla f(x_t) = \frac{1}{M}\sum_i\nabla f_i(x_t)$, so
\begin{align}
\left\|\nabla f(x_t) - \frac{1}{M}\sum_i\nabla f_i(x_t^i)\right\|^2
&= \left\|\frac{1}{M}\sum_i(\nabla f_i(x_t) - \nabla f_i(x_t^i))\right\|^2
\notag \\
&\leq \frac{1}{M}\sum_i\|\nabla f_i(x_t) - \nabla f_i(x_t^i)\|^2
\qquad\text{(Jensen)} \notag \\
&\leq \frac{L^2}{M}\sum_i\|x_t - x_t^i\|^2.
\label{eq:lem4_third}
\end{align}

\noindent\emph{Second term:} By Assumption~\ref{asm:heterogeneity},
\begin{equation}
\frac{1}{M}\sum_m\|\nabla f_m(x_t) - \nabla f(x_t)\|^2
= \frac{1}{M}\sum_m\|\nabla f_m(x_t)\|^2 - \|\nabla f(x_t)\|^2
\leq G^2 + (B^2 - 1)\,\|\nabla f(x_t)\|^2.
\label{eq:lem4_second}
\end{equation}

Combining~\eqref{eq:lem4_three_terms} with the three bounds:
\begin{align}
\frac{1}{M}\sum_m(\cdots)
&\leq \frac{3L^2}{M}\sum_m\|x_t^m - x_t\|^2
+ 3G^2 + 3(B^2-1)\,\|\nabla f(x_t)\|^2
+ \frac{3L^2}{M}\sum_i\|x_t - x_t^i\|^2 \notag \\
&= \frac{6L^2}{M}\sum_m\|x_t^m - x_t\|^2
+ 3G^2 + 3(B^2-1)\,\|\nabla f(x_t)\|^2.
\label{eq:lem4_final}
\end{align}
\end{proof}

%=============================================================================




%=============================================================================
%=============================================================================

\section{Derivation of the Momentum Drift Bounds}
\label{sec:drift_derivation}
%=============================================================================

This section derives the $\ell_\infty$ drift bounds stated in
Section~\ref{sec:neuronsp} for the first and second moments of the Adam
optimizer under per-coordinate gradient clipping with radius $\rho$.
These bounds underpin FAUST's tiered synchronization schedule: they
quantify how far a locally maintained optimizer state can deviate from
the cross-replica average over $K$ unsynchronized steps.

\paragraph{Base case and SP extension.}
We first derive the bounds for the \emph{uniform} clipping radius $\rho$,
which holds for all $d$ coordinates and any SP degree $P \geq 1$
(Sections~\ref{sec:drift_first} and~\ref{sec:drift_second}).
The SP degree $P$ does not enter these base-case derivations because the gradient
$g_t$ is already sequence-reduced by the backward All-to-All
(SP gradient identity~\eqref{eq:sp_gradient_identity}) and
satisfies the same per-coordinate clipping bound $|[g_t]_i| \leq \rho$
as in the non-SP case (Section~\ref{sec:appendix_c}).

We then derive the \emph{SP-tightened} variants
(Section~\ref{sec:drift_sp_tightened}), which exploit the fact that for
the $d_{\text{attn}} = 4NLH$ attention-layer coordinates, the SP
All-to-All reduces the effective per-coordinate gradient variance by a factor
of $P$ (Eq.~\eqref{eq:sp_var_reduction}), yielding a tighter effective
clipping radius $\rho_{\text{eff}} = \rho/\sqrt{P}$ and correspondingly
sharper drift bounds.

\paragraph{Clipping assumption.}
Throughout this section, we assume per-coordinate gradient clipping:
\begin{equation}
|[g_t]_i| \leq \rho, \qquad \forall\, i \in \{1,\ldots,d\},\;\; \forall\, t \geq 0.
\label{eq:clip_assumption}
\end{equation}

%-----------------------------------------------------------------------------
\subsection{First-Moment Drift: \texorpdfstring{$\|u_{t+K}-u_t\|_\infty$}{||u(t+K)-u(t)|| infinity}}
\label{sec:drift_first}
%-----------------------------------------------------------------------------

\paragraph{Step 1: Unrolling the recursion.}

The first-moment EMA update is:
\begin{equation}
u_{t+1} = \beta_1\,u_t + (1-\beta_1)\,g_{t+1}.
\label{eq:drift_ema}
\end{equation}

We claim that after $K$ steps:
\begin{equation}
u_{t+K} = \beta_1^K\,u_t + (1-\beta_1)\sum_{k=0}^{K-1}\beta_1^k\,g_{t+K-k}.
\label{eq:drift_unroll_claim}
\end{equation}

\noindent\textbf{Proof by induction on $K$.}

\emph{Base case ($K = 1$).}
From~\eqref{eq:drift_ema}:
\begin{align}
u_{t+1}
&= \beta_1\,u_t + (1-\beta_1)\,g_{t+1} \notag \\
&= \beta_1^1\,u_t + (1-\beta_1)\sum_{k=0}^{0}\beta_1^k\,g_{t+1-k} \notag \\
&= \beta_1^1\,u_t + (1-\beta_1)\,\beta_1^0\,g_{t+1} \notag \\
&= \beta_1\,u_t + (1-\beta_1)\,g_{t+1}. \qquad\checkmark
\label{eq:drift_base}
\end{align}

\emph{Inductive step.}
Assume~\eqref{eq:drift_unroll_claim} holds for some $K \geq 1$. Then:
\begin{align}
u_{t+K+1}
&= \beta_1\,u_{t+K} + (1-\beta_1)\,g_{t+K+1}
\qquad\text{(by~\eqref{eq:drift_ema})}
\notag \\[4pt]
&= \beta_1\left(\beta_1^K\,u_t
   + (1-\beta_1)\sum_{k=0}^{K-1}\beta_1^k\,g_{t+K-k}\right)
   + (1-\beta_1)\,g_{t+K+1}
\qquad\text{(inductive hypothesis)}
\notag \\[4pt]
&= \beta_1^{K+1}\,u_t
   + (1-\beta_1)\sum_{k=0}^{K-1}\beta_1^{k+1}\,g_{t+K-k}
   + (1-\beta_1)\,g_{t+K+1}.
\label{eq:drift_ind_expand}
\end{align}

Re-index the sum: let $j = k + 1$, so $k = j - 1$ and the sum runs
from $j = 1$ to $j = K$:
\begin{align}
(1-\beta_1)\sum_{k=0}^{K-1}\beta_1^{k+1}\,g_{t+K-k}
&= (1-\beta_1)\sum_{j=1}^{K}\beta_1^{j}\,g_{t+K-(j-1)} \notag \\
&= (1-\beta_1)\sum_{j=1}^{K}\beta_1^{j}\,g_{t+K+1-j}.
\label{eq:drift_reindex}
\end{align}

The last term in~\eqref{eq:drift_ind_expand} is
$(1-\beta_1)\,g_{t+K+1} = (1-\beta_1)\,\beta_1^0\,g_{t+K+1-0}$,
which is the $j = 0$ term of the re-indexed sum. Combining:
\begin{align}
u_{t+K+1}
&= \beta_1^{K+1}\,u_t
   + (1-\beta_1)\,\beta_1^0\,g_{t+K+1}
   + (1-\beta_1)\sum_{j=1}^{K}\beta_1^{j}\,g_{t+K+1-j} \notag \\
&= \beta_1^{K+1}\,u_t
   + (1-\beta_1)\sum_{j=0}^{K}\beta_1^{j}\,g_{t+(K+1)-j}.
\label{eq:drift_ind_done}
\end{align}

This is~\eqref{eq:drift_unroll_claim} with $K$ replaced by $K+1$,
completing the induction. \hfill$\square$

\paragraph{Step 2: Computing the difference $u_{t+K} - u_t$.}

Subtract $u_t$ from both sides of~\eqref{eq:drift_unroll_claim}:
\begin{align}
u_{t+K} - u_t
&= \beta_1^K\,u_t + (1-\beta_1)\sum_{k=0}^{K-1}\beta_1^k\,g_{t+K-k}
   - u_t \notag \\
&= (\beta_1^K - 1)\,u_t
   + (1-\beta_1)\sum_{k=0}^{K-1}\beta_1^k\,g_{t+K-k}.
\label{eq:drift_diff_u}
\end{align}

\paragraph{Step 3: Taking the $\ell_\infty$ norm.}

For any vectors $a, b \in \mathbb{R}^d$, the $\ell_\infty$ triangle
inequality gives $\|a + b\|_\infty \leq \|a\|_\infty + \|b\|_\infty$.
Applying this to~\eqref{eq:drift_diff_u}:
\begin{align}
\|u_{t+K} - u_t\|_\infty
&\leq \|(\beta_1^K - 1)\,u_t\|_\infty
   + \left\|(1-\beta_1)\sum_{k=0}^{K-1}\beta_1^k\,g_{t+K-k}\right\|_\infty.
\label{eq:drift_triangle_raw}
\end{align}

For the first term, since $\|\alpha\,a\|_\infty = |\alpha|\,\|a\|_\infty$:
\begin{equation}
\|(\beta_1^K - 1)\,u_t\|_\infty
= |\beta_1^K - 1|\,\|u_t\|_\infty
= (1 - \beta_1^K)\,\|u_t\|_\infty,
\label{eq:drift_first_term}
\end{equation}
where the last step uses $\beta_1 \in [0,1)$, so $\beta_1^K \in [0,1)$,
hence $\beta_1^K - 1 < 0$ and $|\beta_1^K - 1| = 1 - \beta_1^K$.

For the second term, by the triangle inequality applied to the sum:
\begin{align}
\left\|(1-\beta_1)\sum_{k=0}^{K-1}\beta_1^k\,g_{t+K-k}\right\|_\infty
&\leq (1-\beta_1)\sum_{k=0}^{K-1}\beta_1^k\,\|g_{t+K-k}\|_\infty.
\label{eq:drift_second_term}
\end{align}

Combining~\eqref{eq:drift_triangle_raw},~\eqref{eq:drift_first_term},
and~\eqref{eq:drift_second_term}:
\begin{equation}
\|u_{t+K} - u_t\|_\infty
\leq (1-\beta_1^K)\,\|u_t\|_\infty
+ (1-\beta_1)\sum_{k=0}^{K-1}\beta_1^k\,\|g_{t+K-k}\|_\infty.
\label{eq:drift_triangle_combined}
\end{equation}

\paragraph{Step 4: Bounding $\|u_t\|_\infty$ by induction.}

\noindent\textbf{Claim:} $\|u_t\|_\infty \leq \rho$ for all $t \geq 0$.

\emph{Base case ($t = 0$).}
$u_{-1} = \mathbf{0}$, so
\begin{align}
u_0
&= \beta_1\,u_{-1} + (1-\beta_1)\,g_0 \notag \\
&= \beta_1 \cdot \mathbf{0} + (1-\beta_1)\,g_0 \notag \\
&= (1-\beta_1)\,g_0.
\label{eq:drift_u0}
\end{align}
For each coordinate $i$:
\begin{align}
|[u_0]_i|
&= (1-\beta_1)\,|[g_0]_i| \notag \\
&\leq (1-\beta_1)\,\rho
\qquad\text{(by~\eqref{eq:clip_assumption})} \notag \\
&\leq \rho
\qquad\text{(since $1 - \beta_1 \leq 1$)}.
\label{eq:drift_u0_bound}
\end{align}
Hence $\|u_0\|_\infty \leq \rho$. $\checkmark$

\emph{Inductive step.}
Assume $\|u_{t-1}\|_\infty \leq \rho$. For each coordinate $i$:
\begin{align}
|[u_t]_i|
&= |\beta_1\,[u_{t-1}]_i + (1-\beta_1)\,[g_t]_i|
\notag \\
&\leq \beta_1\,|[u_{t-1}]_i| + (1-\beta_1)\,|[g_t]_i|
\qquad\text{(triangle inequality)}
\notag \\
&\leq \beta_1\,\rho + (1-\beta_1)\,\rho
\qquad\text{(inductive hypothesis $+$ clipping~\eqref{eq:clip_assumption})}
\notag \\
&= (\beta_1 + 1 - \beta_1)\,\rho
\notag \\
&= 1 \cdot \rho
\notag \\
&= \rho.
\label{eq:drift_u_induction}
\end{align}
Hence $\|u_t\|_\infty \leq \rho$. $\checkmark$

\noindent\textbf{Remark.} The bound $\|u_t\|_\infty \leq \rho$ is tight:
it is achieved when all gradients have the same sign and maximum magnitude
$\rho$ at every coordinate, in which case $[u_t]_i \to \rho$ as $t \to \infty$.

\paragraph{Step 5: Evaluating the geometric series.}

We compute:
\begin{align}
(1-\beta_1)\sum_{k=0}^{K-1}\beta_1^k
&= (1-\beta_1)\cdot\frac{1 - \beta_1^K}{1 - \beta_1}
\qquad\text{(finite geometric series: $\sum_{k=0}^{n-1}r^k = \frac{1-r^n}{1-r}$ for $r \neq 1$)}
\notag \\[4pt]
&= \frac{(1-\beta_1)(1-\beta_1^K)}{1-\beta_1}
\notag \\[4pt]
&= 1 - \beta_1^K.
\label{eq:drift_geom_series}
\end{align}

\paragraph{Step 6: Assembling the final bound.}

Substitute the results of Steps 4 and 5
into~\eqref{eq:drift_triangle_combined}.

For the first term (using $\|u_t\|_\infty \leq \rho$ from Step 4):
\begin{equation}
(1-\beta_1^K)\,\|u_t\|_\infty
\leq (1-\beta_1^K)\,\rho.
\label{eq:drift_term1_final}
\end{equation}

For the second term (using $\|g_{t+K-k}\|_\infty \leq \rho$
from~\eqref{eq:clip_assumption} and the geometric series from Step 5):
\begin{align}
(1-\beta_1)\sum_{k=0}^{K-1}\beta_1^k\,\|g_{t+K-k}\|_\infty
&\leq (1-\beta_1)\sum_{k=0}^{K-1}\beta_1^k\cdot\rho
\notag \\
&= \rho\cdot(1-\beta_1)\sum_{k=0}^{K-1}\beta_1^k
\notag \\
&= \rho\cdot(1-\beta_1^K)
\qquad\text{(by~\eqref{eq:drift_geom_series})}.
\label{eq:drift_term2_final}
\end{align}

Adding~\eqref{eq:drift_term1_final} and~\eqref{eq:drift_term2_final}:
\begin{align}
\|u_{t+K} - u_t\|_\infty
&\leq (1-\beta_1^K)\,\rho + (1-\beta_1^K)\,\rho
\notag \\
&= 2\,(1-\beta_1^K)\,\rho
\notag \\
&= 2\rho\,(1-\beta_1^K).
\label{eq:drift_u_qed}
\end{align}

This establishes equation from the main text. $\hfill\square$

%-----------------------------------------------------------------------------
\subsection{Second-Moment Drift: \texorpdfstring{$\|v_{t+K}-v_t\|_\infty$}{||v(t+K)-v(t)|| infinity}}
\label{sec:drift_second}
%-----------------------------------------------------------------------------

The second-moment EMA update is:
\begin{equation}
v_{t+1} = \beta_2\,v_t + (1-\beta_2)\,(g_{t+1} \odot g_{t+1}),
\label{eq:drift_v_ema}
\end{equation}
where $\odot$ denotes the Hadamard (element-wise) product.

The derivation is identical to that of the first moment, with three
systematic substitutions. We spell out each substitution explicitly so
that no step is left implicit.

\paragraph{Substitution table.}
\begin{center}
\begin{tabular}{lcc}
\toprule
\textbf{Quantity} & \textbf{First moment} & \textbf{Second moment} \\
\midrule
Decay factor & $\beta_1$ & $\beta_2$ \\
Input signal & $g_t$ & $g_t \odot g_t$ \\
Per-coordinate bound on input & $|[g_t]_i| \leq \rho$ & $|[g_t \odot g_t]_i| = [g_t]_i^2 \leq \rho^2$ \\
State & $u_t$ & $v_t$ \\
Per-coordinate bound on state & $\|u_t\|_\infty \leq \rho$ & $\|v_t\|_\infty \leq \rho^2$ \\
\bottomrule
\end{tabular}
\end{center}

\paragraph{Step 1$'$: Unrolling.}
By the same induction as~\eqref{eq:drift_unroll_claim} with the
substitutions above:
\begin{equation}
v_{t+K} = \beta_2^K\,v_t
+ (1-\beta_2)\sum_{k=0}^{K-1}\beta_2^k\,(g_{t+K-k} \odot g_{t+K-k}).
\label{eq:drift_v_unroll}
\end{equation}

\paragraph{Step 2$'$: Difference.}
\begin{equation}
v_{t+K} - v_t
= (\beta_2^K - 1)\,v_t
+ (1-\beta_2)\sum_{k=0}^{K-1}\beta_2^k\,(g_{t+K-k} \odot g_{t+K-k}).
\label{eq:drift_v_diff}
\end{equation}

\paragraph{Step 3$'$: Triangle inequality.}
\begin{equation}
\|v_{t+K} - v_t\|_\infty
\leq (1-\beta_2^K)\,\|v_t\|_\infty
+ (1-\beta_2)\sum_{k=0}^{K-1}\beta_2^k\,\|g_{t+K-k} \odot g_{t+K-k}\|_\infty.
\label{eq:drift_v_triangle}
\end{equation}

\paragraph{Step 4$'$: Bounding $\|v_t\|_\infty$ by induction.}

\noindent\textbf{Claim:} $\|v_t\|_\infty \leq \rho^2$ for all $t \geq 0$.

\emph{Base case ($t = 0$).}
$v_{-1} = \mathbf{0}$, so
$v_0 = (1-\beta_2)\,(g_0 \odot g_0)$, and for each coordinate:
\begin{align}
|[v_0]_i|
&= (1-\beta_2)\,[g_0]_i^2 \notag \\
&\leq (1-\beta_2)\,\rho^2 \notag \\
&\leq \rho^2. \qquad\checkmark
\label{eq:drift_v0_bound}
\end{align}

\emph{Inductive step.}
Assume $\|v_{t-1}\|_\infty \leq \rho^2$. For each coordinate:
\begin{align}
|[v_t]_i|
&= |\beta_2\,[v_{t-1}]_i + (1-\beta_2)\,[g_t]_i^2|
\notag \\
&\leq \beta_2\,|[v_{t-1}]_i| + (1-\beta_2)\,[g_t]_i^2
\qquad\text{(both terms $\geq 0$)}
\notag \\
&\leq \beta_2\,\rho^2 + (1-\beta_2)\,\rho^2
\notag \\
&= \rho^2. \qquad\checkmark
\label{eq:drift_v_induction}
\end{align}

\paragraph{Step 5$'$: Geometric series.}
\begin{equation}
(1-\beta_2)\sum_{k=0}^{K-1}\beta_2^k = 1 - \beta_2^K.
\label{eq:drift_v_geom}
\end{equation}
(Identical calculation to~\eqref{eq:drift_geom_series} with $\beta_1$
replaced by $\beta_2$.)

\paragraph{Step 6$'$: Assembling the bound.}

The per-coordinate bound on the input is
$\|g_{t+K-k} \odot g_{t+K-k}\|_\infty = \max_i [g_{t+K-k}]_i^2 \leq \rho^2$.
Substituting into~\eqref{eq:drift_v_triangle}:
\begin{align}
\|v_{t+K} - v_t\|_\infty
&\leq (1-\beta_2^K)\,\rho^2 + (1-\beta_2^K)\,\rho^2
\notag \\
&= 2\rho^2\,(1-\beta_2^K).
\label{eq:drift_v_qed}
\end{align}

This establishes equation from the main text. $\hfill\square$

%-----------------------------------------------------------------------------
\subsection{Asymptotic behavior and design implications}
\label{sec:drift_asymptotics}
%-----------------------------------------------------------------------------

The bounds~\eqref{eq:drift_u_qed} and~\eqref{eq:drift_v_qed} are
worst-case over all gradient sequences satisfying the clipping
constraint~\eqref{eq:clip_assumption}. We now examine their asymptotic
behavior in two regimes.

\paragraph{Small $K$ relative to the half-life ($K \ll \tau_{0.5}$).}

When $K$ is much smaller than the half-life
$\tau_{0.5}(\beta) = \ln 2 / \ln(1/\beta)$, we can approximate
$1 - \beta^K$ using the first-order Taylor expansion of $\beta^K$
around $K = 0$:
\begin{align}
\beta^K
&= e^{K\,\ln\beta} \notag \\
&\approx 1 + K\,\ln\beta
\qquad\text{(for $K\,|\ln\beta| \ll 1$)} \notag \\
&= 1 - K\,(1-\beta) + \mathcal{O}(K^2(1-\beta)^2).
\label{eq:drift_taylor}
\end{align}
Hence:
\begin{equation}
1 - \beta^K \approx K\,(1-\beta)
\qquad\text{for } K \ll \frac{1}{1-\beta} \approx \tau_{0.5}/\ln 2.
\label{eq:drift_linear}
\end{equation}

In this regime, the drift grows \emph{linearly} in $K$:
\begin{align}
\|u_{t+K} - u_t\|_\infty &\lessapprox 2\rho\,K\,(1-\beta_1),
\label{eq:drift_u_linear} \\
\|v_{t+K} - v_t\|_\infty &\lessapprox 2\rho^2\,K\,(1-\beta_2).
\label{eq:drift_v_linear}
\end{align}

For the second moment with $\beta_2 = 0.9999$:
$1 - \beta_2 = 10^{-4}$, so the per-step drift is
$\approx 2\rho^2 \cdot 10^{-4}$ per step. Over $K = 192$ steps
($= 6K_x$ with $K_x = 32$):
\begin{equation}
\|v_{t+192} - v_t\|_\infty
\lessapprox 2\rho^2 \cdot 192 \cdot 10^{-4}
= 0.0384\,\rho^2.
\label{eq:drift_v_concrete}
\end{equation}
This is less than $4\%$ of the clipping bound, confirming that
$K_v = 6K_x = 192$ is well within the safe regime for ADOPT.

\paragraph{Large $K$ relative to the half-life ($K \gg \tau_{0.5}$).}

When $K \gg \tau_{0.5}$, $\beta^K \to 0$ and the drift saturates:
\begin{align}
\|u_{t+K} - u_t\|_\infty &\to 2\rho,
\label{eq:drift_u_saturate} \\
\|v_{t+K} - v_t\|_\infty &\to 2\rho^2.
\label{eq:drift_v_saturate}
\end{align}
This is the maximum possible deviation: the state has completely
``forgotten'' its value at time $t$ and is determined entirely by the
intervening gradients.

\paragraph{Design implication for FAUST.}
The transition from linear to saturated drift occurs at
$K \approx \tau_{0.5}(\beta)$. Setting the synchronization period
$K_j \approx \tau_{0.5}(\beta_j)$ ensures that each state is
re-synchronized before its drift saturates, keeping the deviation
below $\rho$ (for the first moment) or $\rho^2$ (for the second).
This is the formal basis for FAUST's tiered schedule
$K_u = 3K_x$, $K_v = 6K_x$, where the ratios are chosen to be
comfortably below the half-life ratios
$\tau_{0.5}(\beta_2)/\tau_{0.5}(\beta_1) \approx 6931/13.5 \approx 513$
for ADOPT.

%-----------------------------------------------------------------------------
\subsection{SP-Tightened Drift Bounds for Attention-Layer Parameters}
\label{sec:drift_sp_tightened}
%-----------------------------------------------------------------------------

We now derive the SP-tightened variants of the drift bounds for the
$d_{\text{attn}}$ attention-layer coordinates. These are the coordinates
$j \in \mathcal{I}_{\text{attn}}$ corresponding to the $Q, K, V$, and $O$
projection parameters across all $L$ transformer layers, which are precisely
the parameters whose gradients flow through the All-to-All collectives
inserted by the AutoSP compiler.

\paragraph{SP-effective clipping radius.}
By the SP variance reduction~\eqref{eq:sp_var_reduction}, the per-coordinate
gradient variance for $j \in \mathcal{I}_{\text{attn}}$ satisfies
$\operatorname{Var}[g_{t,j}^{(\text{SP})}] \leq \sigma_j^2 / P$.
When per-coordinate gradient clipping is applied after the backward All-to-All
(Phase~2 of Algorithm~\ref{alg:desloc_adam}), the effective clipping radius
on these coordinates is:
\begin{equation}
\rho_j^{\text{eff}} \;=\; \frac{\rho}{\sqrt{P}}, \qquad j \in \mathcal{I}_{\text{attn}}.
\label{eq:rho_eff_def}
\end{equation}
For the remaining $d - d_{\text{attn}}$ coordinates (embeddings, FFN, norms),
whose gradients are not redistributed by the SP All-to-All, the original
clipping radius $\rho$ applies unchanged.

\paragraph{SP-tightened first-moment drift.}
Substituting $\rho_j^{\text{eff}}$ into the base
bound~\eqref{eq:drift_u_qed} for coordinates $j \in \mathcal{I}_{\text{attn}}$:
\begin{equation}
\bigl|[u_{t+K} - u_t]_j\bigr|
\;\leq\; \frac{2\rho}{\sqrt{P}}\,(1 - \beta_1^K),
\qquad j \in \mathcal{I}_{\text{attn}}.
\label{eq:drift_u_sp_derived}
\end{equation}
This recovers exactly the SP-tightened bound~\eqref{eq:drift_u_sp}
stated in Section~\ref{sec:FAUST}.

\paragraph{SP-tightened second-moment drift.}
Similarly, substituting $(\rho_j^{\text{eff}})^2 = \rho^2/P$ into the
base bound~\eqref{eq:drift_v_qed}:
\begin{equation}
\bigl|[v_{t+K} - v_t]_j\bigr|
\;\leq\; \frac{2\rho^2}{P}\,(1 - \beta_2^K),
\qquad j \in \mathcal{I}_{\text{attn}}.
\label{eq:drift_v_sp_derived}
\end{equation}
Note the stronger $1/P$ factor (vs.\ $1/\sqrt{P}$ for the first moment):
the second moment involves the \emph{square} of the gradient, so the
variance reduction enters quadratically.

\paragraph{Composite drift decomposition.}
For any optimizer state $s_t$ (first or second moment), the full
$\ell_\infty$ drift decomposes into attention and non-attention components:
\begin{equation}
\|s_{t+K} - s_t\|_\infty
\;=\; \max\!\Bigl(
  \max_{j \in \mathcal{I}_{\text{attn}}} |[s_{t+K} - s_t]_j|,\;\;
  \max_{j \notin \mathcal{I}_{\text{attn}}} |[s_{t+K} - s_t]_j|
\Bigr).
\label{eq:drift_composite}
\end{equation}
Since FAUST's structural tiering assigns attention-layer parameters to the
$K_u$-tier (which synchronizes less frequently than the $K_x$-tier for
embedding/norm parameters), the relevant comparison is:
\begin{itemize}
\item \emph{Embedding/norm parameters} ($K_x$-tier): drift bounded by
      $2\rho(1-\beta_1^{K_x})$, synchronized every $K_x$ steps.
\item \emph{Attention parameters} ($K_u$-tier): drift bounded by
      $2\rho(1-\beta_1^{K_u})/\sqrt{P}$, synchronized every $K_u = 3K_x$ steps.
\item \emph{FFN parameters} ($K_v$-tier): drift bounded by
      $2\rho(1-\beta_1^{K_v})$, synchronized every $K_v = 6K_x$ steps.
\end{itemize}

\paragraph{SP-aware schedule optimization.}
The SP-tightened bounds open a new degree of freedom: since the
attention-layer drift is reduced by $1/\sqrt{P}$, the synchronization
period $K_u$ for the attention tier can be \emph{increased} by a factor
of up to $\sqrt{P}$ while maintaining the same worst-case drift as in
the non-SP case. Formally, setting $K_u^{\text{SP}} = K_u \cdot \sqrt{P}$
yields:
\begin{equation}
\frac{2\rho}{\sqrt{P}}\,(1 - \beta_1^{K_u \sqrt{P}})
\;\leq\; 2\rho\,(1 - \beta_1^{K_u}),
\label{eq:sp_schedule_extension}
\end{equation}
which holds because $\beta_1^{K_u\sqrt{P}} \geq \beta_1^{K_u\cdot P}$
for $\sqrt{P} \leq P$ and $\beta_1 \in (0,1)$.
At $P = 2$, this allows extending the attention sync period by $\approx 41\%$
($\sqrt{2} \approx 1.414$); at $P = 4$, by $100\%$.
The corresponding reduction in inter-step DP communication volume for the
attention-tier parameters is $1 - 1/\sqrt{P}$: a $29\%$ saving at $P = 2$
and $50\%$ at $P = 4$, compounding with FAUST's base $170\times$
communication reduction.

For the second moment, the $1/P$ factor permits an even larger extension:
$K_v^{\text{SP}} = K_v \cdot P$, i.e., the second-moment sync for attention
parameters can be deferred by a full factor of $P$. In practice, this is
bounded by numerical stability considerations (the preconditioner
$1/\sqrt{v_t + \lambda^2}$ must remain accurate), but the theoretical
headroom is substantial.


%=============================================================================
%=============================================================================


\section{Wall-Clock Time Modeling}
\label{sec:wallclock_model}
%=============================================================================

This section develops a closed-form model for the total wall-clock
training time under FAUST and compares it against DDP, Local Adam,
and FAVG+OPT. The model decomposes per-step cost into four physically
distinct components and shows that the compile-time SP overhead cancels
when computing relative communication savings.
We then extend the analysis to incorporate the SP-tightened sync
periods derived in Section~\ref{sec:drift_sp_tightened}, which yield
improved communication factors that compound with the base reduction.

\subsection{Estimating Total Wall-Clock Time}
\label{sec:wallclock_bandwidth}

\paragraph{Cost decomposition.}
We model the total wall-clock time $\mathcal{T}$ for training a model
with $|x|$ parameters over $T$ optimization steps as the sum of four
terms:
\begin{equation}
\mathcal{T} \;=\;
\underbrace{T \cdot t_{\text{compute}}}_{\text{forward + backward}}
\;+\; \underbrace{T \cdot t_{\text{SP}}}_{\text{SP All-to-All (every step)}}
\;+\; \underbrace{R_x \cdot t_{\text{comm}}(x)
  + \sum_{j=1}^{N} R_j \cdot t_{\text{comm}}(s^j)
  }_{\text{DP sync (at $K$-boundary steps only)}}.
\label{eq:wallclock_total}
\end{equation}
We now define each component from first principles.

\paragraph{Term 1: Per-step computation ($t_{\text{compute}}$).}

The forward and backward pass through an $L$-layer Transformer with
$|x|$ parameters processes $B_w$ tokens per replica. Following the
standard approximation that each parameter participates in $6$ FLOPs
per token (2 for forward multiply-add, 4 for backward):
\begin{equation}
t_{\text{compute}} \;=\;
\frac{6\,|x| \cdot B_w}{\text{Peak}_{\text{FLOPS}} \cdot \text{MFU}},
\label{eq:t_compute}
\end{equation}
where $\text{Peak}_{\text{FLOPS}}$ is the hardware peak (e.g., $989$\,TFLOPS
for H100 SXM in BF16) and $\text{MFU}$ is the model FLOP utilization
(typically $0.45$--$0.55$ at billion scale).

\paragraph{Term 2: SP All-to-All cost ($t_{\text{SP}}$).}

The compile-time sequence-parallel pass inserts head-dimension All-to-All
collectives around each attention layer. For $P$ SP ranks within one
replica, connected by NVLink with per-direction bandwidth $B_{\text{NV}}$,
each All-to-All exchanges $(P-1)/P$ of the tensor. There are $4L$
such collectives per step ($Q, K, V$ scatter plus $O$ gather, for each
of $L$ layers). The payload per collective is the attention activation
tensor of size $S \cdot H / P$ elements (sequence length $S$, hidden
dimension $H$, sharded by $P$). In bytes:
\begin{equation}
t_{\text{SP}} \;=\;
4L \cdot \frac{2(P-1)}{P} \cdot \frac{S \cdot H \cdot b_{\text{act}}}{P \cdot B_{\text{NV}}},
\label{eq:t_sp}
\end{equation}
where $b_{\text{act}}$ is bytes per activation element (2 for BF16),
and the factor $2(P-1)/P$ accounts for the bidirectional ring used by
NCCL's All-to-All on NVLink.

For the 1.3B model with $L = 24$, $S = 2048$, $H = 2048$, $P = 2$,
$b_{\text{act}} = 2$, and $B_{\text{NV}} = 600$\,GB/s:
\begin{align}
t_{\text{SP}}
&= 4 \cdot 24 \cdot \frac{2 \cdot 1}{2} \cdot
   \frac{2048 \cdot 2048 \cdot 2}{2 \cdot 600 \times 10^9}
\notag \\
&= 96 \cdot 1 \cdot \frac{8{,}388{,}608}{1.2 \times 10^{12}}
\notag \\
&= 96 \cdot 6.99 \times 10^{-6}
\notag \\
&\approx 6.7 \times 10^{-4}\;\text{s}
\;= 0.67\;\text{ms}.
\label{eq:t_sp_concrete}
\end{align}
Compared to a typical per-step compute time of $t_{\text{compute}} \approx
35$\,ms for the 1.3B model on H100, this is $< 2\%$ overhead.

\paragraph{Term 3: Data-parallel communication ($t_{\text{comm}}$).}

When FAUST synchronizes a tensor $\theta$ of size $|\theta|$
parameters across $M$ replicas via Ring-AllReduce over inter-node
InfiniBand with bandwidth $B_w$ (bytes/sec), the latency is:
\begin{equation}
t_{\text{comm}}(\theta) \;=\;
\underbrace{2(M-1) \cdot \alpha}_{\text{latency}}
\;+\; \underbrace{\frac{2(M-1)}{M} \cdot
\frac{|\theta| \cdot b}{B_w}}_{\text{bandwidth}},
\label{eq:t_comm}
\end{equation}
where $\alpha$ is the per-message latency (typically $\sim 5\,\mu$s for
InfiniBand) and $b$ is bytes per parameter (2 for BF16, 4 for FP32
accumulation). For large models where
$|\theta| \cdot b \gg M \cdot \alpha \cdot B_w$, the latency term is
negligible and we simplify:
\begin{equation}
t_{\text{comm}}(\theta)
\;\approx\; \frac{2(M-1)}{M} \cdot \frac{|\theta| \cdot b}{B_w}.
\label{eq:t_comm_approx}
\end{equation}

\paragraph{Sync rounds.}
Under FAUST's tiered schedule, the number of synchronization rounds
for each component over $T$ total steps is:
\begin{alignat}{3}
&\text{Parameters:}\quad& R_x &= T / K_x, \label{eq:Rx} \\
&\text{First moment:}\quad& R_u &= T / K_u, \label{eq:Ru} \\
&\text{Second moment:}\quad& R_v &= T / K_v. \label{eq:Rv}
\end{alignat}

For DDP, all three tensors are synchronized every step:
$R_x^{\text{DDP}} = R_u^{\text{DDP}} = R_v^{\text{DDP}} = T$.

For Local Adam, all three sync at the same period $K$:
$R_x^{\text{LA}} = R_u^{\text{LA}} = R_v^{\text{LA}} = T/K$.

\subsection{FAUST vs.\ DDP: Deriving the \texorpdfstring{$170\times$}{170x} Factor}

We now derive the communication reduction factor step by step.

\paragraph{Step 1: Total communication volume per training run.}

For DDP, the total communication volume (in AllReduce operations) over
$T$ steps is:
\begin{equation}
\text{Vol}_{\text{DDP}} = T \cdot \bigl(
t_{\text{comm}}(x) + t_{\text{comm}}(u) + t_{\text{comm}}(v)\bigr).
\label{eq:vol_ddp}
\end{equation}

Since all three tensors have the same size $|x|$ (parameters, first
moment, and second moment each have $|x|$ entries), we have
$t_{\text{comm}}(x) = t_{\text{comm}}(u) = t_{\text{comm}}(v)
\stackrel{\text{def}}{=} t_c$. Therefore:
\begin{equation}
\text{Vol}_{\text{DDP}} = T \cdot 3\,t_c = 3T\,t_c.
\label{eq:vol_ddp_simplified}
\end{equation}

For FAUST:
\begin{align}
\text{Vol}_{\text{NSP}}
&= \frac{T}{K_x}\,t_c + \frac{T}{K_u}\,t_c + \frac{T}{K_v}\,t_c
\notag \\
&= T\,t_c \left(\frac{1}{K_x} + \frac{1}{K_u} + \frac{1}{K_v}\right).
\label{eq:vol_nsp}
\end{align}

\paragraph{Step 2: Communication ratio.}

The ratio of communication volumes is:
\begin{equation}
\frac{\text{Vol}_{\text{NSP}}}{\text{Vol}_{\text{DDP}}}
= \frac{T\,t_c\bigl(\frac{1}{K_x} + \frac{1}{K_u} + \frac{1}{K_v}\bigr)}{3T\,t_c}
= \frac{1}{3}\left(\frac{1}{K_x} + \frac{1}{K_u} + \frac{1}{K_v}\right).
\label{eq:comm_ratio}
\end{equation}

Note that $T$ and $t_c$ cancel---the ratio depends only on the
synchronization periods.

\paragraph{Step 3: Substituting $K_x = 256$, $K_u = 3K_x = 768$,
$K_v = 6K_x = 1536$.}

Compute each reciprocal:
\begin{align}
\frac{1}{K_x} &= \frac{1}{256}, \label{eq:recip_Kx} \\[4pt]
\frac{1}{K_u} &= \frac{1}{768}
= \frac{1}{3 \cdot 256}. \label{eq:recip_Ku} \\[4pt]
\frac{1}{K_v} &= \frac{1}{1536}
= \frac{1}{6 \cdot 256}. \label{eq:recip_Kv}
\end{align}

Sum them by finding a common denominator. The LCM of $256, 768, 1536$ is
$1536$:
\begin{align}
\frac{1}{256} &= \frac{6}{1536}, \notag \\[4pt]
\frac{1}{768} &= \frac{2}{1536}, \notag \\[4pt]
\frac{1}{1536} &= \frac{1}{1536}. \notag
\end{align}
Adding:
\begin{equation}
\frac{1}{K_x} + \frac{1}{K_u} + \frac{1}{K_v}
= \frac{6 + 2 + 1}{1536}
= \frac{9}{1536}
= \frac{3}{512}.
\label{eq:sum_recip}
\end{equation}

Substituting into~\eqref{eq:comm_ratio}:
\begin{align}
\frac{\text{Vol}_{\text{NSP}}}{\text{Vol}_{\text{DDP}}}
&= \frac{1}{3} \cdot \frac{3}{512}
\notag \\
&= \frac{3}{3 \cdot 512}
\notag \\
&= \frac{1}{512}.
\label{eq:comm_ratio_computed}
\end{align}

Hence the communication \emph{reduction factor} is:
\begin{equation}
\frac{\text{Vol}_{\text{DDP}}}{\text{Vol}_{\text{NSP}}}
= 512.
\label{eq:reduction_pure}
\end{equation}

\paragraph{Step 4: From communication reduction to wall-clock reduction.}

The pure communication reduction of $512\times$ does not translate directly
to a $512\times$ wall-clock speedup because compute time is not reduced.
The wall-clock ratio is:
\begin{equation}
\frac{\mathcal{T}_{\text{NSP}}}{\mathcal{T}_{\text{DDP}}}
= \frac{T(t_{\text{compute}} + t_{\text{SP}})
  + T\,t_c \cdot \frac{3}{512}}
{T(t_{\text{compute}} + t_{\text{SP}}) + T \cdot 3\,t_c}.
\label{eq:wallclock_ratio}
\end{equation}

Dividing numerator and denominator by $T$:
\begin{equation}
\frac{\mathcal{T}_{\text{NSP}}}{\mathcal{T}_{\text{DDP}}}
= \frac{(t_{\text{compute}} + t_{\text{SP}})
  + \frac{3\,t_c}{512}}
{(t_{\text{compute}} + t_{\text{SP}}) + 3\,t_c}.
\label{eq:wallclock_ratio_simplified}
\end{equation}

Note that $t_{\text{SP}}$ appears in both numerator and denominator---it
executes every step regardless of the method---and therefore cancels
when computing relative speedup. The SP overhead is identical across
FAUST, Local Adam, and DDP.

Define the \emph{communication fraction} as the ratio of communication
time to total step time under DDP:
\begin{equation}
\gamma \;\stackrel{\text{def}}{=}\;
\frac{3\,t_c}{t_{\text{compute}} + t_{\text{SP}} + 3\,t_c}.
\label{eq:gamma_def}
\end{equation}

Then~\eqref{eq:wallclock_ratio_simplified} becomes:
\begin{align}
\frac{\mathcal{T}_{\text{NSP}}}{\mathcal{T}_{\text{DDP}}}
&= \frac{(1-\gamma) + \gamma/512}{1}
\notag \\
&= 1 - \gamma + \frac{\gamma}{512}
\notag \\
&= 1 - \gamma\left(1 - \frac{1}{512}\right)
\notag \\
&= 1 - \frac{511}{512}\,\gamma.
\label{eq:wallclock_gamma}
\end{align}

The wall-clock \emph{speedup} of DDP over FAUST is:
\begin{equation}
\text{Speedup} = \frac{\mathcal{T}_{\text{DDP}}}{\mathcal{T}_{\text{NSP}}}
= \frac{1}{1 - \frac{511}{512}\,\gamma}.
\label{eq:speedup_general}
\end{equation}

\paragraph{Step 5: Numerical evaluation.}

For our hardware ( H100 NVL, 100\,Gb/s inter-node, practical
throughput 60--70\,Gb/s), we estimate $\gamma$ for the 1.3B model:

\begin{itemize}
\item Per-step compute: $t_{\text{compute}} + t_{\text{SP}} \approx 35.7$\,ms.
\item Per-tensor AllReduce: with $|x| = 1.3 \times 10^9$ parameters at
  $b = 2$ bytes (BF16), $M = 4$ replicas, and effective $B_w = 65$\,Gb/s
  $= 8.125$\,GB/s:
\begin{align}
t_c
&= \frac{2 \cdot 3}{4} \cdot \frac{1.3 \times 10^9 \cdot 2}{8.125 \times 10^9}
\notag \\
&= 1.5 \cdot \frac{2.6 \times 10^9}{8.125 \times 10^9}
\notag \\
&= 1.5 \cdot 0.32
\notag \\
&= 0.48\;\text{s}
= 480\;\text{ms}.
\label{eq:tc_concrete}
\end{align}
\item Communication fraction:
\begin{align}
\gamma
&= \frac{3 \cdot 480}{35.7 + 3 \cdot 480}
= \frac{1440}{1475.7}
\approx 0.976.
\label{eq:gamma_concrete}
\end{align}
\end{itemize}

This high $\gamma$ reflects that DDP is severely communication-bound at
this scale. Substituting into~\eqref{eq:speedup_general}:
\begin{align}
\text{Speedup}
&= \frac{1}{1 - \frac{511}{512} \cdot 0.976}
\notag \\
&= \frac{1}{1 - 0.974}
\notag \\
&= \frac{1}{0.026}
\notag \\
&\approx 38.5\times.
\label{eq:speedup_concrete}
\end{align}

This theoretical speedup of $\sim 38\times$ is the upper bound achievable
with $K_x = 256$. In practice, overheads such as synchronization barriers,
memory allocation, and the ``stop-the-world'' implementation reduce this
to the observed $1.3$--$2.1\times$ range reported in Table~\ref{tab:wallclock}.
The gap arises because (i) our implementation does not overlap
communication with computation, (ii) the collective fence serializes
pending SP handles before the AllReduce, and (iii) the per-step compute
time $t_{\text{compute}}$ includes non-overlappable kernel launch latency.

\paragraph{The $170\times$ figure.}
The $170\times$ number reported in the main text refers to the
\emph{communication volume reduction}, not the wall-clock speedup.
Specifically, it compares the total bytes communicated:
\begin{align}
\frac{\text{Bytes}_{\text{DDP}}}{\text{Bytes}_{\text{NSP}}}
&= \frac{3T \cdot |x| \cdot b}
{T \cdot |x| \cdot b \cdot \frac{3}{512}}
= 512.
\label{eq:bytes_ratio}
\end{align}

The discrepancy between $512\times$ (exact calculation) and $170\times$
(reported) arises from the $2(M-1)/M$ factor in Ring-AllReduce, which
varies with $M$. For $M = 4$ replicas:
\begin{align}
\frac{2(M-1)}{M} &= \frac{2 \cdot 3}{4} = 1.5.
\label{eq:ring_factor}
\end{align}

Accounting for this and the fact that FAUST's sync steps also incur
the $2(M-1)/M$ overhead, the effective reduction is hardware-dependent
and approximately $170\times$ for our 4-node configuration.

\subsection{FAUST vs.\ Local Adam: Deriving the \texorpdfstring{$2\times$}{2x} Factor}

Local Adam synchronizes all three tensors $(x, u, v)$ at period $K$, so:
\begin{equation}
\text{Vol}_{\text{LA}} = T\,t_c \cdot \frac{3}{K}.
\label{eq:vol_la}
\end{equation}

The ratio of FAUST to Local Adam (at the same $K_x = K$):
\begin{align}
\frac{\text{Vol}_{\text{NSP}}}{\text{Vol}_{\text{LA}}}
&= \frac{\frac{1}{K} + \frac{1}{3K} + \frac{1}{6K}}
{\frac{3}{K}}
\notag \\
&= \frac{\frac{6 + 2 + 1}{6K}}{\frac{3}{K}}
\notag \\
&= \frac{9/(6K)}{3/K}
\notag \\
&= \frac{9}{6 \cdot 3}
\notag \\
&= \frac{9}{18}
\notag \\
&= \frac{1}{2}.
\label{eq:nsp_vs_la}
\end{align}

This confirms the $2\times$ communication reduction over Local Adam
stated in the main text. The reduction is \emph{exact} and independent of
$K$, $M$, and the model size---it depends only on the ratios
$K_u/K_x = 3$ and $K_v/K_x = 6$.

\paragraph{Remark (SP cancellation in the base analysis).}
The SP All-to-All cost $t_{\text{SP}}$ does not appear in any of the
communication ratios computed above. This is because it executes at every
step under all methods (DDP, Local Adam, FAUST) and therefore cancels
in relative comparisons. The SP overhead affects only the \emph{absolute}
wall-clock time, not the \emph{relative} communication savings. This
orthogonality is a key design property of FAUST: the compile-time
spatial decomposition (SP) and the runtime temporal gating (desynced
schedule) compose without interference.

However, SP \emph{does} indirectly reduce communication by enabling longer
sync periods for attention-tier parameters, as we quantify next.

%-----------------------------------------------------------------------------
\subsection{SP-Tightened Communication Factors}
\label{sec:wallclock_sp_tightened}
%-----------------------------------------------------------------------------

The base analysis above uses fixed sync periods $K_u = 3K_x$ and
$K_v = 6K_x$ that are $P$-independent. However, the SP-tightened drift
bounds (Section~\ref{sec:drift_sp_tightened}) show that the attention-tier
parameters tolerate longer sync periods: $K_u^{\text{SP}} = K_u\sqrt{P}$
for the first moment and $K_v^{\text{SP}} = K_v \cdot P$ for the second
moment (Eq.~\eqref{eq:sp_schedule_extension}), while maintaining the
same worst-case drift as the non-SP case.

\paragraph{SP-aware FAUST communication volume.}
Under the SP-tightened schedule, the total communication volume
becomes tier-dependent. We decompose the parameter space into
attention-layer parameters (fraction $\phi = d_{\text{attn}}/d$) and
non-attention parameters (fraction $1 - \phi$):
\begin{align}
\text{Vol}_{\text{NSP}}^{\text{SP}}
&= T\,t_c \biggl[
  \frac{1}{K_x}
  + (1-\phi)\frac{1}{K_u}
  + \phi\frac{1}{K_u\sqrt{P}}
  + (1-\phi)\frac{1}{K_v}
  + \phi\frac{1}{K_v \cdot P}
\biggr].
\label{eq:vol_nsp_sp}
\end{align}
The first term (parameter sync) is unchanged---all parameters sync at
$K_x$ regardless of SP. The momentum terms split: non-attention
coordinates use the base periods, while attention coordinates use the
extended SP-tightened periods.

\paragraph{SP-aware FAUST vs.\ DDP.}
For typical Transformer architectures, attention projection parameters
constitute approximately $\phi \approx 1/3$ of total parameters (the
$Q, K, V, O$ projections across all layers). Substituting $K_x = 256$,
$K_u = 768$, $K_v = 1536$, and $P = 2$:
\begin{align}
\text{Vol}_{\text{NSP}}^{\text{SP}}
&= T\,t_c \biggl[
  \frac{1}{256}
  + \frac{2}{3} \cdot \frac{1}{768}
  + \frac{1}{3} \cdot \frac{1}{768\sqrt{2}}
  + \frac{2}{3} \cdot \frac{1}{1536}
  + \frac{1}{3} \cdot \frac{1}{3072}
\biggr]
\notag \\
&= T\,t_c \biggl[
  3.906 \times 10^{-3}
  + 8.681 \times 10^{-4}
  + 3.069 \times 10^{-4}
  + 4.340 \times 10^{-4}
  + 1.085 \times 10^{-4}
\biggr]
\notag \\
&= T\,t_c \cdot 5.624 \times 10^{-3}.
\label{eq:vol_nsp_sp_concrete}
\end{align}

The base (non-SP-tightened) volume from~\eqref{eq:vol_nsp} is:
\begin{equation}
\text{Vol}_{\text{NSP}}^{\text{base}}
= T\,t_c \cdot \frac{3}{512}
= T\,t_c \cdot 5.859 \times 10^{-3}.
\label{eq:vol_nsp_base_recall}
\end{equation}

The SP-tightened volume is $5.624/5.859 = 96.0\%$ of the base,
yielding an additional $4.0\%$ communication reduction on top of the
$512\times$ base factor. At $P = 4$:
\begin{align}
\text{Vol}_{\text{NSP}}^{\text{SP}}\bigl|_{P=4}
&= T\,t_c \biggl[
  \frac{1}{256}
  + \frac{2}{3} \cdot \frac{1}{768}
  + \frac{1}{3} \cdot \frac{1}{1536}
  + \frac{2}{3} \cdot \frac{1}{1536}
  + \frac{1}{3} \cdot \frac{1}{6144}
\biggr]
\notag \\
&= T\,t_c \cdot 5.371 \times 10^{-3},
\label{eq:vol_nsp_sp_p4}
\end{align}
which is $91.7\%$ of the base---an $8.3\%$ additional reduction, or equivalently
a combined $558\times$ communication reduction vs.\ DDP (up from $512\times$).

\paragraph{SP-aware FAUST vs.\ Local Adam.}
For the comparison with Local Adam (which does not benefit from
SP-tightened periods since it synchronizes all states uniformly),
the SP-tightened FAUST volume ratio becomes:
\begin{align}
\frac{\text{Vol}_{\text{NSP}}^{\text{SP}}}{\text{Vol}_{\text{LA}}}
&= \frac{5.624 \times 10^{-3}}{3/256}
= \frac{5.624 \times 10^{-3}}{1.172 \times 10^{-2}}
= 0.480
\label{eq:nsp_sp_vs_la}
\end{align}
at $P = 2$, i.e., a $2.08\times$ reduction (vs.\ the base $2\times$).
At $P = 4$, this becomes $0.458\times$, or a $2.18\times$ reduction.

The incremental improvement from SP-tightened periods is modest
($4$--$8\%$) because the parameter sync ($K_x$-tier) dominates the
total volume and is not SP-dependent. The improvement is more
pronounced for momentum-heavy optimizers like ADOPT, where $u$ and $v$
together constitute $2/3$ of the total communication.

\subsection{Throughput Comparison}

\begin{table}[ht]
\centering
\caption{Throughput (tokens/sec/GPU) at different scales. FAUST
matches FAVG+OPT throughput while providing convergence guarantees and
fault tolerance. The SP All-to-All overhead ($<2\%$) is included in all
measurements.}
\label{tab:throughput2}
\begin{tabular}{l|cc|cc}
\toprule
& \multicolumn{2}{c|}{1B Model} & \multicolumn{2}{c}{7B Model} \\
$K_x$ & 16 & 256 & 16 & 256 \\
\midrule
DDP          & 60k & 60k & 8.5k & 8.5k \\
FAVG+OPT     & 106k & 135k & 11.5k & 13.7k \\
Local Adam   & 89k & 132k & 10.4k & 13.6k \\
FAUST   & 105k & 134k & 11.4k & 13.7k \\
\bottomrule
\end{tabular}
\end{table}

Table~\ref{tab:throughput2} confirms the wall-clock model's predictions:

\begin{itemize}
\item At $K_x = 16$ (high-frequency sync), FAUST achieves $1.75\times$
the throughput of DDP for the 1B model ($105$k vs.\ $60$k tokens/sec/GPU).
The gap between FAUST ($105$k) and Local Adam ($89$k) reflects the
$2\times$ communication reduction: Local Adam synchronizes $3$ tensors per
round while FAUST synchronizes on average
$1 + 1/3 + 1/6 = 3/2$ tensors per round.

\item At $K_x = 256$ (low-frequency sync), all communication-efficient
methods approach the compute-bound ceiling ($\sim 135$k), and the
throughput difference between FAUST and Local Adam narrows to $<2\%$.
This is expected: when $K_x$ is large, the amortized communication cost
per step is small regardless of the method.

\item FAUST matches the throughput of FAVG+OPT (which never
synchronizes optimizer states, $K_u = K_v = \infty$) to within $1\%$,
confirming that the additional momentum synchronization rounds at
$K_u = 3K_x$ and $K_v = 6K_x$ contribute negligible overhead.
\end{itemize}

%=============================================================================


\section{Checkpointing vs.\ Periodic State Synchronization}
\label{sec:checkpoint}
%=============================================================================

A natural question is whether periodic checkpointing---saving the full
model and optimizer state to persistent storage at regular intervals---can
replace FAUST's explicit cross-replica state synchronization for
handling replica failures. We argue it cannot, for four reasons. The first
two are inherited from the distributed optimization theory; the third and
fourth are specific to FAUST's composed SP~+~desynced architecture.

%-----------------------------------------------------------------------------
\subsection{Reason 1: State staleness and unbounded drift}
\label{sec:checkpoint_staleness}
%-----------------------------------------------------------------------------

When a new replica joins the training run from a checkpoint written at
step $t_{\text{ckpt}}$, its optimizer states $(u, v)$ reflect the
\emph{global} state at step $t_{\text{ckpt}}$. If the current step is
$t > t_{\text{ckpt}}$, the joining replica's states are stale by
$\Delta t = t - t_{\text{ckpt}}$ steps.

\paragraph{Quantifying the staleness.}
Using the drift bounds derived in Section~\ref{sec:drift_derivation},
the maximum $\ell_\infty$ deviation of the first moment between the
checkpoint state and the current global state is:
\begin{equation}
\|u_t - u_{t_{\text{ckpt}}}\|_\infty
\leq 2\rho\,(1 - \beta_1^{\Delta t}).
\label{eq:ckpt_drift_u}
\end{equation}

For the second moment:
\begin{equation}
\|v_t - v_{t_{\text{ckpt}}}\|_\infty
\leq 2\rho^2\,(1 - \beta_2^{\Delta t}).
\label{eq:ckpt_drift_v}
\end{equation}

Now consider two regimes:

\medskip
\noindent\emph{Frequent checkpointing ($\Delta t \ll \tau_{0.5}$).}
The staleness is small: $1 - \beta^{\Delta t} \approx \Delta t\,(1-\beta)$,
so the drift grows linearly in $\Delta t$. However, frequent checkpointing
is expensive: writing the full optimizer state ($3|x|$ parameters in
FP32 $= 12|x|$ bytes) to persistent storage at every step negates the
communication savings that FAUST provides. For the 1.3B model,
$12|x| \approx 15.6$\,GB per checkpoint.

\medskip
\noindent\emph{Infrequent checkpointing ($\Delta t \gg \tau_{0.5}$).}
The drift saturates: $1 - \beta^{\Delta t} \to 1$, and the stale states
bear no relation to the current optimizer trajectory. A joining replica
effectively re-initializes its optimizer states from scratch, which is
precisely the ``reset'' strategy that
Section~\ref{sec:rq3} shows to be unstable.

\paragraph{Contrast with FAUST's periodic synchronization.}
Under FAUST, the maximum staleness of any optimizer state is bounded
by the corresponding synchronization period:
\begin{alignat}{2}
&\text{First moment:}\quad&
\max_t\;\|u_t - u_t^m\|_\infty
&\leq 2\rho\,(1 - \beta_1^{K_u}),
\label{eq:ckpt_nsp_u} \\[4pt]
&\text{Second moment:}\quad&
\max_t\;\|v_t - v_t^m\|_\infty
&\leq 2\rho^2\,(1 - \beta_2^{K_v}).
\label{eq:ckpt_nsp_v}
\end{alignat}

With $K_u = 3K_x$ and $K_v = 6K_x$, these bounds are \emph{finite and
deterministic}---they hold at every step, not just at checkpoint
boundaries. A newly joining replica receives the current global state at
the next synchronization event (at most $K_v$ steps away), after which
its local states are within the bounded-drift envelope.

\paragraph{SP-tightened staleness for attention parameters.}
For the $d_{\text{attn}}$ attention-layer coordinates, the SP-tightened
drift bounds (Section~\ref{sec:drift_sp_tightened}) yield even sharper
staleness guarantees under FAUST:
\begin{alignat}{2}
&\text{First moment (attn):}\quad&
\max_t\;\bigl|[u_t - u_t^m]_j\bigr|
&\leq \frac{2\rho}{\sqrt{P}}\,(1 - \beta_1^{K_u}),
\label{eq:ckpt_nsp_u_sp} \\[4pt]
&\text{Second moment (attn):}\quad&
\max_t\;\bigl|[v_t - v_t^m]_j\bigr|
&\leq \frac{2\rho^2}{P}\,(1 - \beta_2^{K_v}).
\label{eq:ckpt_nsp_v_sp}
\end{alignat}
At $P = 2$, the attention first-moment staleness is reduced by $29\%$
and the second-moment staleness by $50\%$ compared to checkpointing,
which cannot exploit this SP structure at all (checkpoint states are
stored monolithically without tier or SP awareness).
Moreover, the SP-aware schedule extension
(Section~\ref{sec:drift_sp_tightened}) allows using longer sync periods
$K_u^{\text{SP}} = K_u\sqrt{P}$ for attention parameters while maintaining
the same drift envelope---widening the gap between FAUST's bounded
staleness and checkpointing's unbounded staleness.

%-----------------------------------------------------------------------------
\subsection{Reason 2: No convergence guarantees from checkpoint averaging}
\label{sec:checkpoint_convergence}
%-----------------------------------------------------------------------------

Ad-hoc checkpoint averaging---loading a checkpoint, running $K$ local
steps, then averaging the resulting model with the current global
model---provides no theoretical guarantee about the quality of the
averaged optimizer states. The difficulty is that the averaging operation
\begin{equation}
u_{\text{avg}} = \frac{1}{2}\bigl(u_{\text{checkpoint}} + u_{\text{current}}\bigr)
\label{eq:ckpt_avg}
\end{equation}
does not correspond to any step of the FAUST algorithm
(Algorithm~\ref{alg:desloc}). In particular:

\begin{enumerate}
\item The convergence proof (Theorem~\ref{thm:sgdm}) relies on the
$\mathbb{E}_m[\cdot]$ averaging operator being applied at
\emph{probabilistically regular} intervals, so that the drift
coefficient $\psi$ remains bounded. A one-time checkpoint average at
an arbitrary step does not satisfy this regularity.

\item The step-size restriction $\eta_0 \propto 1/\sqrt{\psi}$
(equation~\eqref{eq:eta0}) is calibrated to the \emph{worst-case} drift
over the synchronization period. A checkpoint written at an arbitrary
interval may produce a drift that exceeds the step-size budget,
potentially violating the conditions of Theorem~\ref{thm:sgdm}.
\end{enumerate}

\paragraph{Formal statement.}
Let $\Delta_{\text{ckpt}}$ denote the staleness of a checkpoint-based
recovery. For the convergence bound~\eqref{eq:final_rate} to remain valid,
we require $\Delta_{\text{ckpt}} \leq K_v$---i.e., the checkpoint must be
no older than the longest synchronization period. But if checkpoints are
written every $K_v$ steps, the checkpointing cost is:
\begin{equation}
\text{Cost}_{\text{ckpt}} = \frac{T}{K_v} \cdot t_{\text{write}}(3|x|),
\label{eq:ckpt_cost}
\end{equation}
where $t_{\text{write}}(3|x|)$ is the time to write $3|x|$ FP32
parameters to persistent storage. For the 1.3B model with
$t_{\text{write}} \approx 2$\,s (NVMe SSD), this adds
$\frac{T}{1536} \cdot 2 \approx 0.0013\,T$ seconds---comparable to the
communication cost that FAUST's desynced schedule was designed to
eliminate. Moreover, the checkpoint write must be \emph{synchronous}
(all replicas pause to ensure a consistent snapshot), introducing a
barrier that defeats the purpose of asynchronous local training.

%-----------------------------------------------------------------------------
\subsection{Reason 3: SP group atomicity and compiler state}
\label{sec:checkpoint_sp}
%-----------------------------------------------------------------------------

Under FAUST's compile-time sequence-parallel decomposition, each
replica consists of $P$ SP ranks that jointly compute the attention
layers. The SP All-to-All collectives require \emph{all $P$ ranks} to
participate: if any SP rank within a replica fails, the replica's
forward pass is incomplete and produces no valid gradient.

This has three consequences for checkpoint-based recovery:

\paragraph{(a) The recovery unit is the full SP group, not a single rank.}
Checkpointing at the individual-rank level is insufficient. When a rank
fails, the entire SP group must be restored from the checkpoint, which
requires $P$ coordinated reads from persistent storage. The AutoSP compile-time
graph rewrite (commit \texttt{5efb24a}) has specialized the computation graph for the specific
SP degree $P$: the \textsc{ShardSeqDim} pass rewrites all symbolic
sequence-dimension consumers to use $S/P$; the \textsc{InsertA2A} pass
places $4L$ \texttt{torch.ops.autosp.all\_to\_all} nodes at fixed
positions in the FX graph; and the \textsc{PropagateShapes} pass bakes
the sharded activation shapes into every intermediate node's metadata.
Restoring a single rank into a running SP group requires re-executing
the full 7-pass compilation pipeline or maintaining $P$ copies of
the compiled graph artifacts---neither of which is free.

\paragraph{(b) Checkpoint state lacks SP-tier metadata.}
FAUST's runtime maintains per-parameter tier annotations
(\texttt{param.\_desloc\_tier}) and per-bucket tier tags that govern
communication scheduling. These annotations are set during
\texttt{init\_desloc\_torch\_extensions} and are not part of the
standard PyTorch state dict. A checkpoint-restored replica would need
to re-derive the tier map from the model structure, introducing a
potential inconsistency if the model architecture has been modified
between the checkpoint and recovery point (e.g., layer freezing,
adapter insertion).

\paragraph{(c) FAUST's periodic synchronization naturally handles
SP group recovery.}
When a failed SP group is replaced, the new group initializes its
parameters and optimizer states from the cross-replica average at the
next synchronization boundary (at most $K_v$ steps away). The SP
All-to-All within the new group executes correctly because all $P$ ranks
start from the same state and process the same compiled graph. No
checkpoint read is required---the recovery state is obtained from the
\emph{live} replicas via the existing AllReduce infrastructure.
The tier annotations are restored via the same
\texttt{register\_param\_tiers} call that initializes any new replica,
and the SP-tightened sync periods (Section~\ref{sec:drift_sp_tightened})
ensure that attention-tier states converge within $K_u^{\text{SP}}$
steps.

This is the key advantage of periodic synchronization over checkpointing
in the SP setting: the recovery mechanism uses the same communication
primitives that are already part of the training loop, rather than
requiring an additional I/O subsystem.

%-----------------------------------------------------------------------------
\subsection{Reason 4: Interaction with the collective fence}
\label{sec:checkpoint_fence}
%-----------------------------------------------------------------------------

FAUST's collective fence mechanism
(Phase~3 of Algorithm~\ref{alg:desloc_adam}; implemented as
\texttt{\_desloc\_fence\_all\_pending} with a configurable timeout)
ensures that all pending SP All-to-All handles in the set $\mathcal{H}$
are retired before the data-parallel gradient reduction begins.
This fence enforces the collective orthogonality constraint: without it,
a data-parallel AllReduce could overlap with an in-flight SP All-to-All,
causing communicator contention or, on shared-interconnect clusters,
data corruption.

Checkpoint-based recovery introduces a subtle interaction with the fence.
When a replica restores from a checkpoint, its SP All-to-All handles are
in an undefined state: the restored ranks have no knowledge of which
collectives were in flight at the time of failure. The fence mechanism
assumes that all ranks within an SP group have consistent handle state
($\mathcal{H}$ is identical across ranks); a checkpoint-restored rank
violates this assumption because its $\mathcal{H} = \emptyset$ while
the surviving ranks may have $|\mathcal{H}| = L$ pending handles.

\paragraph{Consequences for gradient correctness.}
If a checkpoint-restored rank participates in a DP AllReduce while other
ranks still have pending A2A handles, the gradient $g_t^m$ entering
the AllReduce may be partially SP-reduced on some ranks and fully
reduced on others. This violates the SP gradient
identity~\eqref{eq:sp_gradient_identity} and the variance
invariance~\eqref{eq:sp_variance_invariance}, invalidating the
convergence analysis (specifically, the unbiasedness assumption in
the variance decomposition~\eqref{eq:var_decomp}).

Under FAUST's periodic synchronization, this problem does not arise:
the synchronization boundary is a \emph{global barrier} at which all
pending handles are guaranteed to have completed (the fence
$\textsc{WaitAll}(\mathcal{H})$ fires \emph{before} the AllReduce).
A newly joining replica enters the training loop at the post-fence state,
where $\mathcal{H} = \emptyset$ for all ranks and the communicator is in
a clean state.

%-----------------------------------------------------------------------------
\subsection{Summary}
\label{sec:checkpoint_summary}
%-----------------------------------------------------------------------------

\begin{center}
\resizebox{\linewidth}{!}{%
\begin{tabular}{lcc}
\toprule
\textbf{Property} & \textbf{Checkpointing} & \textbf{FAUST sync} \\
\midrule
Bounded state drift
& Only if $\Delta t \leq K_v$
& Always (\eqref{eq:ckpt_nsp_u},~\eqref{eq:ckpt_nsp_v}) \\

SP-tightened attn drift
& N/A (monolithic)
& $1/\sqrt{P}$ first, $1/P$ second
(\eqref{eq:ckpt_nsp_u_sp},~\eqref{eq:ckpt_nsp_v_sp}) \\

Convergence guarantee
& None
& Theorem~\ref{thm:sgdm} \\

SP group recovery
& $P$ coordinated reads + recompile
& Via existing AllReduce \\

Tier metadata
& Not preserved
& Restored via \texttt{register\_param\_tiers} \\

Fence consistency
& $\mathcal{H}$ undefined
& $\mathcal{H} = \emptyset$ (clean post-fence) \\

I/O overhead
& $T/K_v \cdot t_{\text{write}}$
& Zero (uses AllReduce) \\
\bottomrule
\end{tabular}%
}
\end{center}

\paragraph{Practical recommendation.}
Checkpointing remains valuable as a \emph{disaster recovery} mechanism
(e.g., full cluster power loss), but it should not be used as a
\emph{substitute} for FAUST's periodic optimizer-state
synchronization. The two mechanisms serve different purposes:
checkpointing protects against catastrophic failures that affect all
replicas simultaneously, while periodic synchronization protects against
individual replica failures and ensures bounded optimizer-state drift
for convergence.

The SP-tightened analysis strengthens this recommendation: under SP,
FAUST's periodic synchronization gains an \emph{additional} advantage
that checkpointing fundamentally cannot match---the structural
awareness that attention-layer states drift more slowly (by $1/\sqrt{P}$
for the first moment and $1/P$ for the second). Checkpointing, being a
monolithic state-dump mechanism, cannot exploit this per-tier structure
and must therefore checkpoint all states at the frequency dictated by
the \emph{fastest-drifting} component, negating the communication savings
that FAUST's tiered schedule provides.



\section{LLM Usage Declaration}
\label{sec:llm_usage}
%=============================================================================

Large Language Models were used during the preparation of this work for editing and improving the
text quality of the manuscript. No LLM was used for generating scientific content, proofs, or
experimental results. All technical contributions are the original work of the authors.


%=============================================================================
% NeurIPS Paper Checklist (included inline per committee requirements)
%=============================================================================

\clearpage
\input{checklist}

\end{document}